\chapter{Rings and modules}
\label{chap:rings}


\section{Definition of ring}
\label{sec:rings:definition-of-ring}

In this chapter we will do for \emph{rings} and \emph{modules} what we
have done in \Cref{chap:groups} for groups: describe them in general terms,
with particular attention to distinguished subobjects and
quotients. More detailed information on these structures will be
deferred to later chapters: in \Cref{chap:irred} we will look more carefully
at several interesting classes of rings, and modules (over commutative
rings) will take center stage in our rapid overview of linear algebra
in \Cref{chap:linalg} and following. In this chapter I will also include a
brief jaunt into homological algebra, a topic that will entertain us
greatly in \Cref{chap:homol}.

\subsection{Definition}
\label{subsec:rings:definition}
Rings (and modules) are defined by `decorating' abelian groups with
additional data.  As motivation for the introduction of such
structures, note that all number-based examples of groups that we have
encountered, such as $\Z$ or $\R$, are endowed with an operation of
\emph{multiplication} as well as the `addition' making them into
(abelian) groups. The `ring axioms' will reflect closely the
properties and compatibilities of these two operations in such
examples.  These examples are, however, very special. A more
sophisticated motivation for the introduction of rings arises by
analyzing further the structure of homomorphisms of abelian
groups. Recall
(\S{}\ref{subsec:groups:homomorphisms-of-abelian-groups}) that if $G$,
$H$ are abelian groups, then $\Hom_{\categ{Ab}}(G, H)$ is also an
abelian group. In particular, if $G$ is an abelian group, then so is
the set of \emph{endomorphisms}
$\End_{\categ{Ab}}(G) = \Hom_{\categ{Ab}}(G, G)$. More is true:
morphisms from an object of a category to itself may be composed with
each other (by definition of category!). Thus, two operations coexist
in $\End_{\categ{Ab}}(G)$: addition (inherited from $G$, making
$\End_{\categ{Ab}}(G)$ an abelian group), and composition. These two
operations are compatible with each other in a sense captured by the
\emph{ring axioms}:
\begin{defn}{Ring}{defn:rings:ring}
  A \emph{ring} $(R, +, \cdot)$ is an abelian group $(R, +)$ endowed
  with a second binary operation $\cdot$, satisfying on its own the
  requirements of being associative and having a two-sided identity,
  i.e.,
  \begin{itemize}
  \item
    $(\forall r, s, t \in R) : (r \cdot s) \cdot t = r \cdot (s \cdot
    t)$,
  \item
    $(\exists 1_R \in R)(\forall r \in R) : r \cdot 1_R = r = 1_R
    \cdot r$
  \end{itemize}
  (which make $(R, \cdot)$ a \emph{monoid}), and further interacting
  with $+$ via the following distributive properties:
  \begin{itemize}
  \item
    $(\forall r, s, t \in R) : (r+s) \cdot t = r \cdot t + s \cdot t$
    and $t \cdot (r+s) = t \cdot r + t \cdot s$.
  \end{itemize}
\end{defn}
The notation $\cdot$ is often omitted in formulas, and we usually
refer to rings by the name of the underlying set.  \emph{Warning:}
What I am calling a `ring', others may call a `ring with identity' or
a `ring with 1': it is not uncommon to exclude the axiom of existence
of a `multiplicative identity' from the list of axioms defining a
ring. Rings without identity are sometimes called \emph{rngs}, but I
am not sure this should be encouraged.\footnote{The term `rng' was
  introduced with this meaning by Jacobson; but essentially at the
  same time Mac Lane introduced $\categ{Rng}$ as the name for the
  category of rings with identity. Hoping to steer clear of this clash
  of terminology I have opted to call this category `$\categ{Ring}$'.}
The reader should check conventions carefully when approaching the
literature.  Examples of structures without a multiplicative identity
abound: for example, the set $2\Z$ of even integers, with the usual
addition and multiplication, satisfies all the ring axioms given above
with the exception of the existence of $1$ (and is therefore a
rng). But in these notes all rings will have $1$.  Of course the
multiplicative identity is necessarily unique: the argument given for
\Cref{prop:groups:identity-unique} works verbatim.  The identity element of the
abelian group underlying a ring is denoted $0_R$ (or simply $0$, in
context) and is called the `additive' identity. This is a special
element with respect to multiplication:
\begin{lem}{In a ring $R$, $0 \cdot r = 0 = r \cdot 0$ for all
    $r \in R$.}{lem:rings:zero-mult}
\end{lem}
\begin{proof}
  Indeed, $0 = 0 + 0$; hence, applying distributivity,
  \[
    r \cdot 0 = r \cdot (0 + 0) = r \cdot 0 + r \cdot 0,
  \]
  from which $r \cdot 0 = 0$ by cancellation (in the group $(R,
  +)$). The equality $0 \cdot r = 0$ is proven similarly.
\end{proof}
It is equally easy to check that multiplication behaves as expected on
`subtraction'.  In fact, if $-1$ denotes the additive inverse of $1$,
then the additive inverse $-r$ of any $r \in R$ is the result of the
multiplication $(-1) \cdot r$: indeed, using distributivity,
\[
  r + (-1) \cdot r = 1 \cdot r + (-1) \cdot r = (1 - 1) \cdot r = 0
  \cdot r = 0
\]
(by \cref{lem:rings:zero-mult}) from which $(-1)r = -r$ follows by
(additive) cancellation.

\subsection{First examples and special classes of rings}
\label{subsec:rings:first-examples-and-special-classes-of-rings}

\begin{exam}{Zero-ring}{}
  We can define a ring structure on a trivial group $\{*\}$ by letting
  $* \cdot * = *$ (as well as $* + * = *$); this is often called the
  \emph{zero-ring}. Note that $0 = 1$ in this ring (cf.\
  \Cref{exer:rings:zero-ring}).
\end{exam}
\begin{exam}{Number rings}{}
  More interesting examples are the number-based groups such as $\Z$
  or $\R$, with the usual operations. These are very well known to our
  reader, who will realize immediately that they satisfy the
  requirements given in \cref{defn:rings:ring}; but they are very
  special. Why?
\end{exam}
To begin with, note that multiplication is commutative in these
examples; this is not among the requirements we have posed on rings in
the official definition given above.
\begin{defn}{Commutative ring}{defn:rings:commutative}
  A ring $R$ is \emph{commutative} if
  \begin{itemize}
  \item $(\forall r, s \in R) : r \cdot s = s \cdot r$.
  \end{itemize}
\end{defn}
\emph{Commutative rings} (with identity) form an extremely important
class of rings; \emph{commutative algebra} is the subfield of algebra
studying them. We will focus on commutative rings in later chapters;
in this chapter we will develop some of the basic theory for the more
general case of arbitrary rings (with $1$).
\begin{exam}{Matrix ring}{exam:rings:matrix-ring-noncommutative}
  An example of a noncommutative ring that is (likely) familiar to the
  readers is the ring of $2 \times 2$ matrices with, say, real
  entries: matrices can be added `componentwise', and they can be
  multiplied as recalled in \Cref{exam:groups:gln-matrices}; the two operations satisfy
  the requirements in \cref{defn:rings:ring}.

  Square matrices of any size, and with entries in any ring, form a
  ring (\Cref{exer:rings:matrix-ring}).
\end{exam}
\begin{exam}{$\Z/n\Z$}{}
  The reader is already familiar with a large class of (commutative)
  rings: the groups $\Z/n\Z$, endowed with the multiplication defined
  in \S{}\ref{subsec:groups:cyclic-groups-and-modular-arithmetic} (that is: $[a]_n \cdot [b]_n := [ab]_n$; this is
  well-defined (cf.\ \Cref{exer:groups:multiplication-well-defined})) satisfy the ring axioms listed
  above.
\end{exam}
The rings $\Z/n\Z$ prompt me to highlight an important point. Another
reason why rings such as $\Z$, $\Q$, $\R$, \dots\ are special is that
multiplicative cancellation by nonzero elements holds in these
rings. Of course additive cancellation is automatic, since rings are
in particular (abelian) groups; and multiplicative cancellation
clearly fails in general since one cannot `cancel $0$' (by
\cref{lem:rings:zero-mult}). But even the fact that
\[
  (\forall a \in R,\, a \neq 0) : a \cdot b = a \cdot c \implies b =
  c,
\]
which holds, for example, in $\Z$, does not follow from the ring
axioms.  Indeed, this cancellation property does not hold in all
rings; it may well fail in the rings $\Z/n\Z$. For example,
\[ [2]_6 \cdot [4]_6 = [8]_6 = [2]_6 = [2]_6 \cdot [1]_6
\]
even though $[4]_6 \neq [1]_6$.  The problem here is that in $\Z/6\Z$
there are elements $a \neq 0$ such that $a \cdot b = 0$ for some
$b \neq 0$ (take $a = [2]_6$, $b = [3]_6$).
\begin{defn}{Left-zero-divisor}{defn:rings:zero-divisor}
  An element $a$ in a ring $R$ is a \emph{left-zero-divisor} if there
  exist elements $b \neq 0$ in $R$ for which $ab = 0$.
\end{defn}
The reader will have no difficulty figuring out what a
\emph{right}-zero-divisor should be. The element $0$ is a zero-divisor
in all nonzero rings $R$; the zero ring is the only ring without
zero-divisors(!).
\begin{prop}{}{prop:rings:zero-divisor-injection}
  In a ring $R$, $a \in R$ is not a left- (resp., right-) zero-divisor
  if and only if left (resp., right) multiplication by $a$ is an
  injective function $R \to R$.
\end{prop}
In other words, $a$ is not a left- (resp., right-) zero-divisor if and
only if multiplicative left- (resp., right-) cancellation by the
element $a$ holds in $R$.
\begin{proof}
  Let's verify the `left' statement (the `right' statement is of
  course entirely analogous). Assume $a$ is not a left-zero-divisor
  and $ab = ac$ for $b, c \in R$.  Then, by distributivity,
  \[
    a(b - c) = ab - ac = 0,
  \]
  and this implies $b - c = 0$ since $a$ is not a left-zero-divisor;
  that is, $b = c$. This proves that left-multiplication is injective
  in this case.  Conversely, if $a$ is a left-zero-divisor, then
  $\exists b \neq 0$ such that $ab = 0 = a \cdot 0$; this shows that
  left-multiplication is not injective in this case, concluding the
  proof.
\end{proof}
Rings such as $\Z$, $\Q$, etc., are commutative rings \emph{without}
(nonzero) zero-divisors.  Such rings are very special, but very
important, and they deserve their own terminology:
\begin{defn}{Integral domain}{defn:rings:integral-domain}
  An \emph{integral domain} is a nonzero commutative ring $R$ (with
  $1$) such that
  \[
    (\forall a, b \in R) : ab = 0 \implies a = 0 \text{ or } b = 0.
  \]
\end{defn}
\Cref{chap:irred} will be entirely devoted to integral domains.  An element
which is not a zero-divisor is called a \emph{non-zero-divisor}. Thus,
integral domains are those nonzero commutative rings in which every
nonzero element is a non-zero-divisor. By
\Cref{prop:rings:zero-divisor-injection}, multiplicative cancellation
by nonzero elements holds in integral domains. The rings $\Z$, $\Q$,
$\R$, $\C$ are all integral domains.  As we have seen, some $\Z/n\Z$
are \emph{not} integral domains.  Here is one of those places where
the reader can do him/herself a great favor by pausing a moment and
figuring something out: answer the question, which $\Z/n\Z$ are
integral domains? This is entirely within reach, given what the reader
knows already. Don't read ahead before figuring this out---this
question will be answered within a few short paragraphs, spoiling all
the fun.  There are even subtler reasons why $\Z$ is a very special
ring: we will see in due time that it is a `UFD' (unique factorization
domain); in fact, it is a `PID' (principal ideal domain); in fact, it
is more special still!, as it is a `Euclidean domain'. All of this
will be discussed in \Cref{chap:irred}, particularly \S{}\ref{sec:irred:ufds-pids-euclidean-domains}.

However, $\Q$, $\R$, $\C$ are more special than all of that and then
some, since they are \emph{fields}.
\begin{defn}{Left-unit, right-unit, unit}{defn:rings:unit}
  An element $u$ of a ring $R$ is a \emph{left-unit} if
  $\exists v \in R$ such that $uv = 1$; it is a \emph{right-unit} if
  $\exists v \in R$ such that $vu = 1$. \emph{Units} are two-sided
  units.
\end{defn}
\begin{prop}{}{prop:rings:units}
  In a ring $R$:
  \begin{itemize}
  \item $u$ is a left- (resp., right-) unit if and only if left-
    (resp., right-) multiplication by $u$ is a surjective function
    $R \to R$;
  \item if $u$ is a left- (resp., right-) unit, then right- (resp.,
    left-) multiplication by $u$ is injective; that is, $u$ is not a
    right- (resp., left-) zero-divisor;
  \item the inverse of a two-sided unit is unique;
  \item two-sided units form a group under multiplication.
  \end{itemize}
\end{prop}
\begin{proof}
  These assertions are all straightforward. For example, denote by
  $\rho_u : R \to R$ right-multiplication by $u$, so that
  $\rho_u(r) = ru$. If $u$ is a right-unit, let $v \in R$ be such that
  $vu = 1$; then $\forall r \in R$
  \[
    \rho_u \circ \rho_v(r) = \rho_u(rv) = (rv)u = r(vu) = r \cdot 1_R
    = r.
  \]
  That is, $\rho_v$ is a right-inverse to $\rho_u$, and therefore
  $\rho_u$ is surjective (\cref{prop:inv-inj-sur}).  Conversely, if
  $\rho_u$ is surjective, then there exists a $v$ such that
  $1_R = \rho_u(v) = vu$, so that $u$ is a right-unit.

  This checks the first statement, for right-units.

  For the second statement, denote by $\lambda_u : R \to R$
  left-multiplication by $u$: $\lambda_u(r) = ur$. Assume $u$ is a
  right-unit, and let $v$ be such that $vu = 1_R$; then
  $\forall r \in R$
  \[
    \lambda_v \circ \lambda_u(r) = \lambda_v(ur) = v(ur) = (vu)r = 1_R
    r = r.
  \]
  That is, $\lambda_v$ is a left-inverse to $\lambda_u$, so
  $\lambda_u$ is injective (\cref{prop:inv-inj-sur} again).

  The rest of the proof is left to the reader (\Cref{exer:rings:prop112-proof}).
\end{proof}
Since the inverse of a two-sided unit $u$ is unique, we can give it a
name; of course we denote it by $u^{-1}$. The reader should keep in
mind that inverses of left- or right-units are \emph{not} unique in
general, so the `inverse notation' is not appropriate for them.
\begin{defn}{Division ring}{defn:rings:division-ring}
  A \emph{division ring} is a ring in which every nonzero element is a
  two-sided unit.
\end{defn}
We will mostly be concerned with the commutative case, which has its
own name:
\begin{defn}{Field}{defn:rings:field}
  A \emph{field} is a nonzero commutative ring $R$ (with $1$) in which
  every nonzero element is a unit.
\end{defn}
The whole of \Cref{chap:fields} will be devoted to studying fields.

By \Cref{prop:rings:units} (second part), every field is an integral
domain, but not conversely: indeed, $\Z$ is an integral domain, but it
is not a field. Remember:
\[
  \text{field} \Rightarrow \text{integral domain,}
\]
\[
  \text{integral domain} \not\Rightarrow \text{field.}
\]
There is a situation, however, in which the two notions coincide:
\begin{prop}{}{prop:rings:finite-field-domain}
  Assume $R$ is a finite commutative ring; then $R$ is an integral
  domain if and only if it is a field.
\end{prop}
\begin{proof}
  One implication holds for all rings, as pointed out above; thus we
  only have to verify that if $R$ is a finite integral domain, then it
  is a field. This amounts to verifying that if $a$ is a
  non-zero-divisor in a finite (commutative) ring $R$, then it is a
  unit in $R$.

  Now, if $a$ is a non-zero-divisor, then multiplication by $a$ in $R$
  is injective (\cref{prop:rings:zero-divisor-injection}); hence it is
  surjective, as the ring is finite, by the pigeon-hole principle;
  hence $a$ is a unit, by \cref{prop:rings:units}.
\end{proof}
\begin{rem}{Wedderburn's little theorem}{rem:rings:wedderburn}
  A little surprisingly, the hypothesis of commutativity in
  \Cref{prop:rings:finite-field-domain} is actually superfluous: a
  theorem known as \emph{Wedderburn's little theorem} shows that
  finite division rings are necessarily commutative. The reader will
  prove this fact in a distant future (\Cref{exer:fields:wedderburn-thm-witts-proof}).
\end{rem}
\begin{exam}{}{exam:rings:znz-units-group}
  The \emph{group of units} in the ring $\Z/n\Z$ is precisely the
  group $(\Z/n\Z)^*$ introduced in \S{}\ref{subsec:groups:cyclic-groups-and-modular-arithmetic}: indeed, a class $[m]_n$
  is a unit if and only if (right-) multiplication by $[m]_n$ is
  surjective (by \cref{prop:rings:units}), if and only if the map
  $a \mapsto a[m]_n$ is surjective, if and only if $[m]_n$ generates
  $\Z/n\Z$, if and only if $\gcd(m, n) = 1$
  (\cref{coro:groups:generator-coprime}), if and only if
  $[m]_n \in (\Z/n\Z)^*$.

  In particular, those $n$ for which all nonzero elements of $\Z/n\Z$
  are units (that is, for which $\Z/n\Z$ is a field) are precisely
  those $n \in \Z$ for which $\gcd(m, n) = 1$ for all $m$ that are not
  multiples of $n$; this is the case if and only if $n$ is
  prime. Putting this together with
  \Cref{prop:rings:finite-field-domain}, we get the pretty
  classification (for integers $p \neq 0$)
  \[
    \Z/p\Z \text{ integral domain} \iff \Z/p\Z \text{ field} \iff p
    \text{ prime,}
  \]
  which the reader is well advised to remember firmly.
\end{exam}
\begin{exam}{}{}
  The rings $\Z/p\Z$, with $p$ prime, are \emph{not} the only finite
  fields. In fact, for every prime $p$ and every integer $r > 0$ there
  is a (unique, in a suitable sense) multiplication on the product
  group
  \[
    \underbrace{\Z/p\Z \times \cdots \times \Z/p\Z}_{r \text{ times}}
  \]
  making it into a field. A discussion of these fields will have to
  wait until we have accumulated much more material (cf.\
  \S{}\ref{subsec:fields:finite-fields}), but the reader could already try to construct small
  examples `by hand' (cf.\ \Cref{exer:rings:four-element-field}).
\end{exam}

\subsection{Polynomial rings}
\label{subsec:rings:polynomial-rings}
We will study polynomial rings in some depth, especially over fields;
they are another class of examples that is to some extent already
familiar to our reader. I will capitalize on this familiarity and
avoid a truly formal (and truly tedious) definition.
\begin{defn}{Polynomial}{defn:rings:polynomial}
  Let $R$ be a ring. A \emph{polynomial} $f(x)$ in the
  \emph{indeterminate} $x$ and with \emph{coefficients} in $R$ is a
  finite linear combination of nonnegative `powers' of $x$ with
  coefficients in $R$:
  \[
    f(x) = \sum_{i \geq 0} a_i x^i = a_0 + a_1 x + a_2 x^2 + \cdots,
  \]
  where all $a_i$ are elements of $R$ (the \emph{coefficients}) and we
  require $a_i = 0$ for $i \gg 0$.  Two polynomials are taken to be
  equal if all the coefficients are equal:
  \[
    \sum_{i \geq 0} a_i x^i = \sum_{i \geq 0} b_i x^i \iff (\forall i
    \geq 0) : a_i = b_i.
  \]
\end{defn}
The set of polynomials in $x$ over $R$ is denoted by $R[x]$. Since all
but finitely many $a_i$ are assumed to be $0$, one usually employs the
notation
\[
  f(x) = a_0 + a_1 x + \cdots + a_n x^n
\]
for $\sum_{i \geq 0} a_i x^i$, if $a_i = 0$ for $i > n$.  At this
point the reader should just view all of this as a \emph{notation}: a
polynomial really stands for an element of an infinite direct sum of
the group $(R, +)$. The `polynomial' notation is more suggestive as it
hints at what \emph{operations} we are going to impose on $R[x]$: if
\[
  f(x) = \sum_{i \geq 0} a_i x^i \quad \text{and} \quad g(x) = \sum_{i
    \geq 0} b_i x^i,
\]
then we define
\[
  f(x) + g(x) := \sum_{i \geq 0} (a_i + b_i) x^i
\]
and
\[
  f(x) \cdot g(x) := \sum_{k \geq 0} \sum_{i+j=k} a_i b_j x^{i+j}.
\]
To clarify this latter definition, see how it works for small $k$:
$f(x) \cdot g(x)$ equals
\[
  a_0 b_0 + (a_0 b_1 + a_1 b_0)x + (a_0 b_2 + a_1 b_1 + a_2 b_0)x^2 +
  (a_0 b_3 + a_1 b_2 + a_2 b_1 + a_3 b_0)x^3 + \cdots,
\]
that is, business as usual.  It is essentially straightforward
(\Cref{exer:rings:rx-mult-assoc}) to check that $R[x]$, with these operations, is a
ring; the identity $1$ of $R$ is the identity of $R[x]$, when viewed
as a polynomial (that is, $1_{R[x]} = 1_R + 0x + 0x^2 + \cdots$).  The
\emph{degree} of a nonzero polynomial
$f(x) = \sum_{i \geq 0} a_i x^i$, denoted $\deg f(x)$, is the largest
integer $d$ for which $a_d \neq 0$. This notion is very useful, but
really behaves well (\Cref{exer:rings:degree-inequalities}) only if $R$ is an integral domain:
for example, note that over $R = \Z/6\Z$
\begin{align*}
  &\deg([1] + [2]x) = 1, \quad \deg([1] + [3]x) = 1, \quad \text{but} \\
  &\deg(([1] + [2]x) \cdot ([1] + [3]x)) = \deg([1] + [5]x) = 1 \neq 1 + 1.
\end{align*}
Polynomials of degree $0$ (together with $0$) are called
\emph{constants}; they form a `copy' of $R$ in $R[x]$, since the
operations $+$, $\cdot$ on constant polynomials are nothing but the
original operations in $R$, up to this identification. It is sometimes
convenient to assign to the polynomial $0$ the degree $-\infty$.
Polynomial rings in more indeterminates may be obtained by iterating
this construction:
\[
  R[x, y, z] := R[x][y][z];
\]
elements of this ring may be written as `ordinary' polynomials in
three indeterminates and are manipulated as usual. It can be checked
easily that the construction does not really depend on the order in
which the indeterminates are listed, in the sense that different
orderings lead to isomorphic rings (in the sense soon to be defined
officially). Different indeterminates commute with each other;
constructions analogous to polynomial rings, but with noncommuting
indeterminates, are also very important, but we will not develop them
in this book (we will glance at one such notion in \Cref{linrep:4.17}).
We will occasionally consider a polynomial ring in \emph{infinitely
  many} indeterminates: for example, we will denote by
$R[x_1, x_2, \ldots]$ the case of countably many indeterminates. Keep
in mind, however, that polynomials are \emph{finite} linear
combinations of \emph{finite} products of the indeterminates; in
particular, every given element of $R[x_1, x_2, \ldots]$ only involves
\emph{finitely many} indeterminates. An honest definition of this ring
involves \emph{direct limits}, which await us in \S{}\ref{subsec:linrep:limits-and-colimits}.  Rings
of \emph{power series} may be defined and are very useful; the ring of
series $\sum_{i=0}^{\infty} a_i x^i = a_0 + a_1 x + a_2 x^2 + \ldots$
in $x$ with coefficients in $R$, and evident operations, is denoted
$R[[x]]$. Regrettably, we will only occasionally encounter these rings
in this book.  The ring $R[x]$ is (clearly) commutative if $R$ is
commutative; it is an integral domain if $R$ is an integral domain
(\Cref{exer:rings:rx-integral-domain}); but it has no chances of being a field even if $R$ is
a field, since $x$ has no inverse in $R[x]$. The question of which
properties of $R$ are `inherited' by $R[x]$ is subtle and important,
and we will give it a great deal of attention in later sections.

\subsection{Monoid rings}
\label{subsec:rings:monoid-rings}
The polynomial ring is an instance of a rather general construction,
which is occasionally very useful. A \emph{semigroup} is a set endowed
with an associative operation; a \emph{monoid} is a semigroup with an
identity element. Thus a group is a monoid in which every element has
an inverse; positive integers with ordinary addition form a semigroup,
while the set $\N$ of \emph{natural numbers} (that is, nonnegative
integers\footnote{Some disagree, and insist that $\N$ should not
  include $0$.}) is a monoid under addition.  Given a monoid
$(M, \cdot)$ and a ring $R$, we can obtain a new ring $R[M]$ as
follows.  Elements of $R[M]$ are formal linear combinations
\[
  \sum_{m \in M} a_m \cdot m
\]
where the `coefficients' $a_m$ are elements of $R$ and $a_m \neq 0$
for at most finitely many summands (hence, as in \S{}\ref{subsec:rings:polynomial-rings}, as an
abelian group $R[M]$ is nothing but the direct sum $R^{\oplus
  M}$). Operations in $R[M]$ are defined by
\[
  \Bigl(\sum_{m \in M} a_m \cdot m\Bigr) + \Bigl(\sum_{m \in M} b_m
  \cdot m\Bigr) = \sum_{m \in M} (a_m + b_m) \cdot m,
\]
\[
  \Bigl(\sum_{m \in M} a_m \cdot m\Bigr) \cdot \Bigl(\sum_{m \in M}
  b_m \cdot m\Bigr) = \sum_{m \in M} \sum_{m_1 m_2 = m} (a_{m_1}
  b_{m_2}) \cdot m.
\]
The identity in $R[M]$ is $1_R \cdot 1_M$, viewed as a formal sum in
which all other summands have $0$ as coefficient.

The reader will hopefully see the similarity with the construction of
the polynomial ring $R[x]$ in \S{}\ref{subsec:rings:polynomial-rings}; in fact (\Cref{exer:rings:rx-monoid-ring}) the
polynomial ring $R[x]$ may be interpreted as $R[\N]$.

\emph{Group rings} are the result of this construction when $M$ is in
fact a group. The group ring $R[\Z]$ is a ring of `Laurent
polynomials' $R[x, x^{-1}]$, allowing for negative as well as positive
exponents.

\subsection*{Exercises}
\begin{exer}{\exerused}{exer:rings:zero-ring}
  Prove that if $0 = 1$ in a ring $R$, then $R$ is a
  zero-ring. [\S{}\ref{subsec:rings:first-examples-and-special-classes-of-rings}]
\end{exer}
\begin{exer}{$\lnot$}{exer:rings:power-set-ring}
  Let $S$ be a set, and define operations on the power set
  $\mathcal{P}(S)$ of $S$ by setting $\forall A, B \in \mathcal{P}(S)$
  \begin{align*}
    A + B &:= (A \cup B) \setminus (A \cap B), \\
    A \cdot B &= A \cap B
  \end{align*}
  (where the solid black contour indicates the set included in the
  operation).  Prove that $(\mathcal{P}(S), +, \cdot)$ is a
  commutative ring. [\ref{exer:rings:powerset-vs-z2s}, \ref{exer:rings:boolean-power-set}]
\end{exer}
\begin{exer}{$\lnot$}{exer:rings:function-ring}
  Let $R$ be a ring, and let $S$ be any set. Explain how to endow the
  set $R^S$ of set-functions $S \to R$ of two operations $+$, $\cdot$
  so as to make $R^S$ into a ring, such that $R^S$ is just a copy of
  $R$ if $S$ is a singleton. [\ref{exer:rings:powerset-vs-z2s}]
\end{exer}
\begin{exer}{\exerused}{exer:rings:matrix-ring}
  The set of $n \times n$ matrices with entries in a ring $R$ is
  denoted $\mathcal{M}_n(R)$. Prove that componentwise addition and
  matrix multiplication make $\mathcal{M}_n(R)$ into a ring, for any
  ring $R$. The notation $\mathfrak{gl}_n(R)$ is also commonly used,
  especially for $R = \R$ or $\C$ (although this indicates one is
  considering them as Lie algebras) in parallel with the analogous
  notation for the corresponding groups of units; cf.\
  \Cref{exer:groups:matrix-subgroups-inclusions}. In fact, the parallel continues with the definition
  of the following sets of matrices:
  \begin{itemize}
  \item
    $\mathfrak{sl}_n(\R) = \{M \in \mathfrak{gl}_n(\R) \mid
    \operatorname{tr}(M) = 0\}$;
  \item
    $\mathfrak{sl}_n(\C) = \{M \in \mathfrak{gl}_n(\C) \mid
    \operatorname{tr}(M) = 0\}$;
  \item
    $\mathfrak{so}_n(\R) = \{M \in \mathfrak{sl}_n(\R) \mid M + M^t =
    0\}$;
  \item
    $\mathfrak{su}(n) = \{M \in \mathfrak{sl}_n(\C) \mid M + M^\dagger
    = 0\}$.
  \end{itemize}
  Here $\operatorname{tr}(M)$ is the \emph{trace} of $M$, that is, the
  sum of its diagonal entries. The other notation matches the notation
  used in \Cref{exer:groups:matrix-subgroups-inclusions}. Can we make rings of these sets by endowing
  them with ordinary addition and multiplication of matrices? (These
  sets are all Lie algebras; cf.\ \Cref{exer:linalg:lie-bracket-axioms-and-category}.) [\S{}\ref{subsec:rings:first-examples-and-special-classes-of-rings}, \ref{exer:rings:quaternion-matrix-reps},
  \ref{exer:rings:endomorphism-algebra}, \ref{exer:linalg:lie-algebra-sets-are-r-vector-spaces}, \ref{exer:linalg:lie-bracket-axioms-and-category}]
\end{exer}
\begin{exer}{}{exer:rings:sum-zero-divisors}
  Let $R$ be a ring. If $a$, $b$ are zero-divisors in $R$, is $a+b$
  necessarily a zero-divisor?
\end{exer}
\begin{exer}{$\lnot$}{exer:rings:nilpotent}
  An element $a$ of a ring $R$ is \emph{nilpotent} if $a^n = 0$ for
  some $n$.
  \begin{itemize}
  \item Prove that if $a$ and $b$ are nilpotent in $R$ and $ab = ba$,
    then $a+b$ is also nilpotent.
  \item Is the hypothesis $ab = ba$ in the previous statement
    necessary for its conclusion to hold?
  \end{itemize} [\ref{exer:rings:nilpotents-form-ideal}]
\end{exer}
\begin{exer}{}{exer:rings:nilpotent-znz}
  Prove that $[m]$ is nilpotent in $\Z/n\Z$ if and only if $m$ is
  divisible by all prime factors of $n$.
\end{exer}
\begin{exer}{}{exer:rings:x-squared-one}
  Prove that $x = \pm 1$ are the only solutions to the equation
  $x^2 = 1$ in an integral domain. Find a ring in which the equation
  $x^2 = 1$ has more than $2$ solutions.
\end{exer}
\begin{exer}{\exerused}{exer:rings:prop112-proof}
  Prove \Cref{prop:rings:units}. [\S{}\ref{subsec:rings:first-examples-and-special-classes-of-rings}]
\end{exer}
\begin{exer}{}{exer:rings:right-unit-inverses}
  Let $R$ be a ring. Prove that if $a \in R$ is a right-unit and has
  two or more left-inverses, then $a$ is not a left-zero-divisor and
  is a right-zero-divisor.
\end{exer}
\begin{exer}{\exerused}{exer:rings:four-element-field}
  Construct a field with $4$ elements: as mentioned in the text, the
  underlying abelian group will have to be $\Z/2\Z \times \Z/2\Z$;
  $(0, 0)$ will be the zero element, and $(1, 1)$ will be the
  multiplicative identity. The question is what $(0, 1) \cdot (0, 1)$,
  $(0, 1) \cdot (1, 0)$, $(1, 0) \cdot (1, 0)$ must be, in order to
  get a field. [\S{}\ref{subsec:rings:first-examples-and-special-classes-of-rings}, \S{}\ref{subsec:irred:roots-and-reducibility}]
\end{exer}
\begin{exer}{\exerused}{exer:rings:quaternions}
  Just as complex numbers may be viewed as combinations $a + bi$,
  where $a, b \in \R$ and $i$ satisfies the relation $i^2 = -1$ (and
  commutes with $\R$), we may construct a ring\footnote{The letter
    $\field{H}$ is chosen in honor of William Rowan Hamilton.}
  $\field{H}$ by considering linear combinations $a + bi + cj + dk$
  where $a, b, c, d \in \R$ and $i, j, k$ commute with $\R$ and
  satisfy the following relations:
  \begin{align*}
    i^2 = j^2 = k^2 &= -1, \\
    ij = -ji &= k, \\
    jk = -kj &= i, \\
    ki = -ik &= j.
  \end{align*}
  Addition in $\field{H}$ is defined componentwise, while
  multiplication is defined by imposing distributivity and applying
  the relations. For example,
  \[
    (1 + i + j) \cdot (2 + k) = 1 \cdot 2 + i \cdot 2 + j \cdot 2 + 1
    \cdot k + i \cdot k + j \cdot k = 2 + 2i + 2j + k - j + i = 2 + 3i
    + j + k.
  \]
  \begin{enumerate}[label=(\roman*)]
  \item Verify that this prescription does indeed define a ring.
  \item Compute $(a + bi + cj + dk)(a - bi - cj - dk)$, where
    $a, b, c, d \in \R$.
  \item Prove that $\field{H}$ is a division ring.
  \end{enumerate}
  Elements of $\field{H}$ are called \emph{quaternions}. Note that
  $Q_8 := \{\pm 1, \pm i, \pm j, \pm k\}$ forms a subgroup of the
  group of units of $\field{H}$; it is a noncommutative group of
  order $8$, called the \emph{quaternionic group}.
  \begin{enumerate}[resume, label=(\roman*)]
  \item List all subgroups of $Q_8$, and prove that they are all
    normal.
  \item Prove that $Q_8$, $D_8$ are not isomorphic.
  \item Prove that $Q_8$ admits the presentation
    $\langle x, y \mid x^2 y^{-2}, y^4, xyx^{-1}y \rangle$.
  \end{enumerate} [\S{}\ref{subsec:groups:normal-subgroups}, \ref{exer:rings:quaternion-matrix-reps}, \ref{exer:groups2:class-formula-d8-q8-verify}, \ref{exer:groups2:q8-not-semidirect-prod}, \ref{exer:groups2:quaternions-as-semidirect-prod}, \ref{exer:irred:integral-quaternions-ring-basics}]
\end{exer}
\begin{exer}{\exerused}{exer:rings:rx-mult-assoc}
  Verify that the multiplication defined in $R[x]$ is
  associative. [\S{}\ref{subsec:rings:polynomial-rings}]
\end{exer}


\begin{exer}{\exerused}{exer:rings:degree-inequalities}
  Let $R$ be a ring, and let $f(x), g(x) \in R[x]$ be nonzero
  polynomials. Prove that
  \[
    \deg(f(x) + g(x)) \leq \max(\deg(f(x)), \deg(g(x))).
  \]
  Assuming that $R$ is an integral domain, prove that
  \[
    \deg(f(x) \cdot g(x)) = \deg(f(x)) + \deg(g(x)).
  \] [\S{}\ref{subsec:rings:polynomial-rings}]
\end{exer}

\begin{exer}{\exerused}{exer:rings:rx-integral-domain}
  Prove that $R[x]$ is an integral domain if and only if $R$ is an
  integral domain. [\S{}\ref{subsec:rings:polynomial-rings}]
\end{exer}

\begin{exer}{}{exer:rings:power-series-units}
  Let $R$ be a ring, and consider the ring of power series $R[[x]]$
  (cf.\ \S{}\ref{subsec:rings:polynomial-rings}).
  \begin{enumerate}[label=(\roman*)]
  \item Prove that a power series $a_0 + a_1 x + a_2 x^2 + \cdots$ is
    a unit in $R[[x]]$ if and only if $a_0$ is a unit in $R$. What is
    the inverse of $1 - x$ in $R[[x]]$?
  \item Prove that $R[[x]]$ is an integral domain if and only if $R$
    is.
  \end{enumerate}
\end{exer}

\begin{exer}{\exerused}{exer:rings:rx-monoid-ring}
  Explain in what sense $R[x]$ agrees with the monoid ring
  $R[\N]$. [\S{}\ref{subsec:rings:monoid-rings}]
\end{exer}

\section{The Category Ring}
\label{sec:rings:the-category-ring}
\subsection{Ring homomorphisms}
\label{subsec:rings:ring-homomorphisms}
\emph{Ring homomorphisms} are defined in the natural way: if $R$, $S$
are rings, a function $\varphi : R \to S$ is a \emph{ring
  homomorphism} if it preserves both operations and the identity
element. That is, $\varphi$ must be a homomorphism of the underlying
abelian groups,
\[
  (\forall a, b \in R) : \quad \varphi(a + b) = \varphi(a) +
  \varphi(b),
\]
it must preserve the operation of multiplication,
\[
  (\forall a, b \in R) : \quad \varphi(ab) = \varphi(a)\varphi(b),
\]
and finally
\[
  \varphi(1_R) = 1_S.
\]
It is evident that rings form a category, with ring homomorphisms as
morphisms. I will denote this category by $\categ{Ring}$.  The
zero-ring is clearly final in $\categ{Ring}$. However, note that it is
\emph{not} initial: because of the requirement that ring homomorphisms
send $1$ to $1$, the only rings to which zero-rings map
homomorphically are the zero-rings (\Cref{exer:rings:zero-ring-domain}).

The category $\categ{Ring}$ does have initial objects: the ring of
integers $\Z$ (with the usual operations $+$, $\cdot$) is initial in
$\categ{Ring}$. Indeed, for every ring $R$ we can define a group
homomorphism $\varphi : \Z \to R$ by
\[
  (\forall n \in \Z) : \quad \varphi(n) = n \cdot 1_R,
\]
that is, as the `exponential map' $\epsilon_{1_R}$ corresponding to
$1_R \in R$; cf.\ \S{}\ref{subsec:groups:examples}. But $\varphi$ is in fact a ring
homomorphism, since $\varphi(1) = 1_R$, and
\[
  \varphi(mn) = (mn)1_R = m(n1_R) \overset{!}{=} (m1_R) \cdot (n1_R) =
  \varphi(m) \cdot \varphi(n),
\]
where the equality $\overset{!}{=}$ holds by the distributivity
axiom. This ring homomorphism is unique, since it is determined by the
requirement that $\varphi(1) = 1_R$ and by the fact that $\varphi$
preserves addition. Thus for every ring $R$ there exists one and only
one ring homomorphism $\Z \to R$, showing that $\Z$ is initial in
$\categ{Ring}$.  This should already convince the reader that $\Z$ is
a very special ring. There is in fact nothing arbitrary about $\Z$:
knowledge of the group $\Z$ determines the ring structure of $\Z$, as
will be underscored below. This is one of the main reasons why I
included the `identity' axiom in the list in \Cref{defn:rings:ring}:
there are many nonisomorphic structures of `ring without identity' on
the group $(\Z, +)$ (cf.\ \Cref{exer:rings:z-mz-isomorphic-groups}), but only one structure of
ring with identity (\Cref{exer:rings:unique-ring-structure-on-z}).  Ring homomorphisms preserve
units: that is, if $u$ is a (left-, resp., right-) unit in $R$ and
$\varphi : R \to S$ is a ring homomorphism, then $\varphi(u)$ is a
(left-, resp., right-) unit. Indeed, if $v$ is a (right-, say) inverse
of $u$, then
\[
  \varphi(u)\varphi(v) = \varphi(uv) = \varphi(1_R) = 1_S,
\]
so that $\varphi(v)$ is a (right-) inverse of $\varphi(u)$.  On the
other hand, the image of a non-zero-divisor by a ring homomorphism may
well be a zero-divisor: for example, the canonical projection
$\pi : \Z \to \Z/6\Z$ is a ring homomorphism, and $2$ is a
non-zero-divisor in $\Z$, yet $\pi(2) = [2]_6$ is a zero-divisor.

\subsection{Universal property of polynomial rings}
\label{subsec:rings:universal-property-of-polynomial-rings}
Polynomial rings satisfy a universal property not unlike the one for
free groups explored in \S{}\ref{subsec:groups:universal-property}. The simplest (and possibly most
useful) case is for the polynomial rings $\Z[x_1, \ldots, x_n]$, and
with respect to \emph{commutative} rings; I will leave the reader the
pleasure of stating and proving fancier notions.  Let
$A = \{a_1, \ldots, a_n\}$ be a set of order $n$. Consider the
category $\mathscr{R}_A$ whose objects are pairs $(j, R)$, where $R$
is a commutative ring\footnote{For the reader interested in
  generalizations: only the requirement that $j(a_1), \ldots, j(a_n)$
  commute with one another is needed here.} and
\[
  j : A \to R
\]
is a set-function (cf.\ \S{}\ref{subsec:groups:universal-property}); morphisms
$(j_1, R_1) \to (j_2, R_2)$ are commutative diagrams
\[
  \begin{tikzcd}
    R_1 \arrow[rr, "\varphi"] & & R_2 \\
    & A \arrow[ul, "j_1"] \arrow[ur, "j_2"']
  \end{tikzcd}
\]
in which $\varphi$ is a ring homomorphism.

For example, $(i, \Z[x_1, \ldots, x_n])$ is an object of
$\mathscr{R}_A$, where $i : A \to \Z[x_1, \ldots, x_n]$ sends $a_k$ to
$x_k$.
\begin{prop}{}{prop:rings:poly-universal}
  $(i, \Z[x_1, \ldots, x_n])$ is initial in $\mathscr{R}_A$.
\end{prop}
\begin{proof}
  Let $(j, R)$ be an arbitrary object of $\mathscr{R}_A$; we have to
  show that there is a unique morphism
  $(i, \Z[x_1, \ldots, x_n]) \to (j, R)$, that is, there exists
  exactly one ring homomorphism $\varphi : \Z[x_1, \ldots, x_n] \to R$
  such that
  \[
    \begin{tikzcd}
      \Z[x_1, \cdots, x_n] \arrow[rr, "\varphi"] & & R \\
      & A \arrow[ul, "i"] \arrow[ur, "j"']
    \end{tikzcd}
  \]
  commutes.

  As usual in these verifications, the key point is that the
  requirements posed on $\varphi$ force its definition. The postulated
  commutativity of the diagram means that $\varphi(x_k) = j(a_k)$ for
  $k = 1, \ldots, n$. Then, since $\varphi$ must be a ring
  homomorphism, necessarily
  \begin{align*}
    \varphi\Bigl(\sum m_{i_1 \cdots i_n} x_1^{i_1} \cdots x_n^{i_n}\Bigr) &= \sum \varphi(m_{i_1 \cdots i_n}) \varphi(x_1)^{i_1} \cdots \varphi(x_n)^{i_n} \\
                                                                          &= \sum \iota(m_{i_1 \cdots i_n}) j(a_1)^{i_1} \cdots j(a_n)^{i_n},
  \end{align*}
  where $\iota : \Z \to R$ is the unique ring homomorphism (as $\Z$ is
  initial in $\categ{Ring}$).

  Thus, if $\varphi$ exists, then it is unique. On the other hand, the
  formula we just obtained clearly\footnote{Well, it is really clear
    that it preserves addition; multiplication requires a bit of work,
    which the reader would be well advised to perform. This is where
    the commutativity of $R$ is used.} preserves the operations and
  sends $1$ to $1$, so it does define a ring homomorphism, concluding
  the proof.
\end{proof}
\begin{exam}{}{exam:rings:poly-universal-n1}
  For $n = 1$, \Cref{prop:rings:poly-universal} says that if $s$ is
  any element of a ring $S$, then there is a unique ring homomorphism
  $\Z[x] \to S$ sending $x$ to $s$ and `extending' the unique ring
  homomorphism $\iota : \Z \to S$. In this case the commutativity of
  $S$ is immaterial (why?).
\end{exam}
\begin{exam}{}{exam:rings:poly-universal-general}
  More generally, let $\alpha : R \to S$ be a fixed ring homomorphism,
  and let $s \in S$ be an element commuting with $\alpha(r)$ for all
  $r \in R$. Then there is a unique ring homomorphism
  $\bar{\alpha} : R[x] \to S$ extending $\alpha$ and sending $x$ to
  $s$ (\Cref{exer:rings:poly-extension-property}).

  In particular, we get an `evaluation map' for polynomials over
  commutative rings as follows. Given a polynomial
  $f(x) = \sum_{i \geq 0} a_i x^i \in R[x]$, every $r \in R$
  determines an element
  \[
    f(r) = \sum_{i \geq 0} a_i r^i\,;
  \]
  this may be viewed as $\bar{\alpha}(f(x))$, where $\bar{\alpha}$ is
  obtained as above with $\alpha = \id_R : R \to R$ and
  $s = r$.

  Thus, every polynomial $f(x)$ determines a \emph{polynomial
    function} $f : R \to R$, defined by $r \mapsto f(r)$. It is a good
  idea to keep the two concepts of `polynomial' and `polynomial
  \emph{function}' well distinct (cf.\ \Cref{exer:rings:z2-poly-function}).
\end{exam}

\subsection{Monomorphisms and epimorphisms}
\label{subsec:rings:monomorphisms-and-epimorphisms}
The \emph{kernel} of a homomorphism $\varphi : R \to S$ of rings is
\[
  \ker \varphi := \{r \in R \mid \varphi(r) = 0\}\,;
\]
that is, it is nothing but the kernel of $\varphi$ when the latter is
viewed as a homomorphism of groups. As such, $\ker \varphi$ is a
subgroup of $(R, +)$; it satisfies an even stronger
requirement,\footnote{Note that $\ker \varphi$ cannot be a
  sub\emph{ring} in any reasonable sense, according to our definition
  of ring, since it does not contain a multiplicative identity in all
  but very pathological situations. It is a subring if the identity
  requirement is omitted.} which we will explore at great length in
\S{}\ref{sec:rings:ideals-and-quotient-rings}.

The reader may hope that an analog to \Cref{prop:groups:monomorphism-characterization} holds in
$\categ{Ring}$, and this is indeed the case:
\begin{prop}{}{prop:rings:mono-equiv}
  For a ring homomorphism $\varphi : R \to S$, the following are
  equivalent:
  \begin{enumerate}[label=(\alph*)]
  \item $\varphi$ is a monomorphism;
  \item $\ker \varphi = \{0\}$;
  \item $\varphi$ is injective (as a set-function).
  \end{enumerate}
\end{prop}
\begin{proof}
  We prove (a) $\implies$ (b), leaving the rest to the reader. Assume
  $\varphi : R \to S$ is a monomorphism and $r \in \ker
  \varphi$. Applying the extension property of \Cref{exam:rings:poly-universal-n1}, we obtain
  unique ring homomorphisms $\mathrm{ev}_r : \Z[x] \to R$ such that
  $\mathrm{ev}_r(x) = r$ and $\mathrm{ev}_0 : \Z[x] \to R$ such that
  $\mathrm{ev}_0(x) = 0$. Consider the parallel ring homomorphisms:
  \[
    \begin{tikzcd}
      \Z[x] \arrow[r, shift left, "\mathrm{ev}_r"] \arrow[r, shift
      right, "\mathrm{ev}_0"'] & R \arrow[r, "\varphi"] & S\,;
    \end{tikzcd}
  \]
  since $\varphi(r) = 0 = \varphi(0)$, the two compositions
  $\varphi \circ \mathrm{ev}_r$, $\varphi \circ \mathrm{ev}_0$ agree
  (because they agree on $\Z$ and they agree on $x$); hence
  $\mathrm{ev}_r = \mathrm{ev}_0$ since $\varphi$ is a monomorphism.
  Therefore
  \[
    r = \mathrm{ev}_r(x) = \mathrm{ev}_0(x) = 0.
  \]
  This proves $r \in \ker \varphi \implies r = 0$, that is, (b).
\end{proof}
By \Cref{prop:rings:mono-equiv}, if $S \to R$ is a monomorphism, then
$S$ may be identified with a subset of $R$; the following definition
formalizes this situation.
\begin{defn}{Subring}{defn:rings:subring}
  A \emph{subring} $S$ of a ring $R$ is a ring whose underlying set is
  a subset of $R$ and such that the inclusion function
  $S \hookrightarrow R$ is a ring homomorphism.
\end{defn}
Equivalently, a subring of $R$ is a subset $S \subseteq R$ which
contains $1_R$ and satisfies the ring axioms with respect to the
operations $+$, $\cdot$ induced from $R$. It is immediately checked
that $S \subseteq R$ is a subring if it is a subgroup of $(R, +)$, it
is closed with respect to $\cdot$, and it contains $1_R$. As a
nonexample, the zero-ring is \emph{not} a subring of a nonzero ring
$R$ (because it does not contain $1_R$).  \Cref{prop:rings:mono-equiv}
may induce the reader to believe that $\categ{Ring}$ and $\categ{Grp}$
are rather similar from this general categorical standpoint. But a new
phenomenon occurs concerning epimorphisms. A surjective map of rings
is certainly an epimorphism in $\categ{Ring}$, since it is already an
epimorphism in $\categ{Set}$---we have run into this argument earlier
(e.g., `(c) $\implies$ (a)' in \Cref{prop:groups:epimorphism-characterization}); but unlike as in
$\categ{Set}$, $\categ{Grp}$, and $\categ{Ab}$, epimorphisms need not
be surjective\footnote{Unfortunately, some references \emph{define}
  ring epimorphisms as `surjective ring homomorphisms'; this should be
  discouraged.} in $\categ{Ring}$. Indeed, consider the inclusion
homomorphism of rings
\[
  \iota : \Z \hookrightarrow \Q\,;
\]
$\iota$ is not surjective, and hence it is not an epimorphism in
$\categ{Set}$ or $\categ{Ab}$ (can the reader `describe'
$\coker\iota$? \Cref{exer:rings:z-in-q-cokernel}); but it \emph{is} an
epimorphism of rings.  Indeed, if $\alpha_1$, $\alpha_2$ are parallel
ring homomorphisms
\[
  \begin{tikzcd}
    \Z \arrow[r, hook, "\iota"] & \Q \arrow[r, shift left, "\alpha_1"]
    \arrow[r, shift right, "\alpha_2"'] & R
  \end{tikzcd}
\]
and $\alpha_1$, $\alpha_2$ agree\footnote{As they must, since $\Z$ is
  initial.} on $\Z$, say $\alpha = \alpha_1|_{\Z} = \alpha_2|_{\Z}$,
then they must agree on $\Q$: because for $p, q \in \Z$, $q \neq 0$,
\[
  \alpha_i\!\left(\frac{p}{q}\right) = \alpha_i(p)\alpha_i(q^{-1}) =
  \alpha(p)\alpha(q)^{-1} \qquad (i = 1, 2)
\]
is the same for both.\footnote{Note that $\alpha(q)$ must have a
  double-sided inverse in $R$ since $q$ has one in $\Q$, namely
  $1/q$. In particular, $\alpha_1(q^{-1}) = \alpha_2(q^{-1})$ must
  agree, since they must equal this unique inverse of $\alpha(q)$;
  cf.\ \cref{prop:rings:units}.} Thus, $\iota$ satisfies the
categorical requirement for epimorphisms (cf.\ \S{}\ref{subsec:prelims:functions:monomorphisms-and-epimorphisms}).
\emph{Warning:} Thus, in $\categ{Ring}$, a homomorphism may well be
both a monomorphism and an epimorphism without being an isomorphism!

\subsection{Products}
\label{subsec:rings:products}
\textit{Products} exist in $\categ{Ring}$: if $R_1$, $R_2$ are rings,
then $R_1 \times R_2$ may be defined by endowing the direct product of
groups $R_1 \times R_2$ (cf.\ \S{}\ref{subsec:groups:products-et-al}) with componentwise
multiplication. Thus, both operations on $R_1 \times R_2$ are defined
componentwise: $\forall (a_1, a_2), (b_1, b_2) \in R_1 \times R_2$,
\[
  (a_1, a_2) + (b_1, b_2) := (a_1 + b_1, a_2 + b_2),
\]
\[
  (a_1, a_2) \cdot (b_1, b_2) := (a_1 \cdot b_1, a_2 \cdot b_2).
\]
The identity in $R_1 \times R_2$ is $(1_{R_1}, 1_{R_2})$. The reader
will check (\Cref{exer:rings:componentwise-product}) that $R_1 \times R_2$ is indeed a
categorical product in $\categ{Ring}$.  The reader should keep in mind
that in general this is \emph{not} the only ring structure one can
define on the direct product of underlying groups. For example, as
mentioned in \S{}\ref{subsec:rings:first-examples-and-special-classes-of-rings}, one can define a field whose underlying group is
$\Z/p\Z \times \Z/p\Z$, while the product ring $\Z/p\Z \times \Z/p\Z$
is very far from being a field (why?).  The situation with
\emph{coproducts} is more unfortunate---dealing with this in any
generality (even for the case of \emph{commutative} rings) requires
\emph{tensor products}, and by the time we develop tensor products
(\S{}\ref{sec:linrep:tensor-products-and-the-tor-functors}), we will have almost forgotten this question. But the
universal property reviewed in \S{}\ref{subsec:rings:universal-property-of-polynomial-rings} suffices to deal with simple
examples, and the reader should work out one template case now, for
fun (\Cref{exer:rings:poly-two-var-coproduct}).  Incidentally---with the case of abelian groups
in mind (\S{}\ref{subsec:groups:abelian-groups}), the reader may be tempted to consider a `direct
sum of rings', agreeing with the product in the finite case and maybe
giving something interesting in general. This is less promising than
it looks: it does not satisfy the universal property of coproducts in
$\categ{Ring}$ (why?), and, further, the `infinite version' is not a
ring because it is missing the identity.

\subsection{\texorpdfstring{$\End_{\categ{Ab}}(G)$}{End\_Ab(G)}}
\label{subsec:rings:endab-g}
This is as good a place as any to expand a little on one of the main
examples of rings, and one that is not quite as special as the
examples reviewed in \S{}\ref{subsec:rings:first-examples-and-special-classes-of-rings}.  For every abelian group $G$, the group
$\End_{\categ{Ab}}(G) := \Hom_{\categ{Ab}}(G, G)$ of endomorphisms of
$G$ is a ring, under the operations of addition and composition (as
discussed in the paragraph preceding \Cref{defn:rings:ring}). Indeed,
associativity follows from the category axioms, and distributivity is
immediately checked.

It is instructive to examine carefully the endomorphism ring of the
abelian group $(\Z, +)$:
\begin{prop}{}{prop:rings:endab-Z}
  $\End_{\categ{Ab}}(\Z) \cong \Z$ as rings.
\end{prop}
\begin{proof}
  Consider the function
  \[
    \varphi : \End_{\categ{Ab}}(\Z) \to \Z
  \]
  defined by $\varphi(\alpha) = \alpha(1)$ for all group homomorphisms
  $\alpha : \Z \to \Z$. Then $\varphi$ is a group homomorphism: the
  addition in $\End_{\categ{Ab}}(\Z)$ is defined so that
  $\forall n \in \Z$
  \[
    (\alpha + \beta)(n) = \alpha(n) + \beta(n)
  \]
  (cf.\ \S{}\ref{subsec:groups:homomorphisms-of-abelian-groups}); in particular
  \[
    \varphi(\alpha + \beta) = (\alpha + \beta)(1) = \alpha(1) +
    \beta(1) = \varphi(\alpha) + \varphi(\beta).
  \]
  Further, $\varphi$ is a ring homomorphism. Indeed, for $\alpha$,
  $\beta$ in $\End_{\categ{Ab}}(\Z)$ denote $\alpha(1)$ by $a$; then
  \[
    \alpha(n) = n\alpha(1) = na = an
  \]
  for all $n \in \Z$; in particular,
  $\alpha(\beta(1)) = a\beta(1) = \alpha(1)\beta(1)$. Therefore,
  \[
    \varphi(\alpha \circ \beta) = (\alpha \circ \beta)(1) =
    \alpha(\beta(1)) = \alpha(1)\beta(1) =
    \varphi(\alpha)\varphi(\beta)
  \]
  as needed. Also,
  $\varphi(\id_{\Z}) = \id_{\Z}(1) = 1$.

  Finally, $\varphi$ has an inverse: for $a \in \Z$, let $\psi(a)$ be
  the homomorphism $\alpha : \Z \to \Z$ defined by
  \[
    (\forall n \in \Z) : \quad \alpha(n) = an;
  \]
  the reader will easily check that $\psi$ is a ring homomorphism and
  inverse to $\varphi$.

  Therefore $\varphi$ is a ring isomorphism, verifying the statement.
\end{proof}
\Cref{prop:rings:endab-Z} gives a sense in which the ring structure on
$\Z$ is truly `natural': this structure arises naturally from
categorical considerations, there is nothing `arbitrary' about it.
More generally, the rings $\End_{\categ{Ab}}(G)$ and other rings
arising as endomorphism rings of other structures (e.g., vector
spaces) are arguably the most important class of examples of
rings. The entire theory of modules is based on the observation that
$\End_{\categ{Ab}}(G)$ is a ring for every abelian group $G$.  Some of
the notions reviewed in \S{}\ref{subsec:rings:first-examples-and-special-classes-of-rings} are expressed very concretely in these
endomorphism rings; for example, the group of units in
$\End_{\categ{Ab}}(G)$ is nothing but
$\Aut_{\categ{Ab}}(G)$.  Every ring $R$ interacts with
the ring of endomorphisms of the underlying abelian group $(R,
+)$. For $r \in R$, define left- and right-multiplication by $r$ by
$\lambda_r$, $\mu_r$, respectively. That is, $\forall a \in R$,
\[
  \lambda_r(a) = ra, \qquad \mu_r(a) = ar.
\]
The following observation is nothing more than easy notation juggling,
but it is useful; it is a `ring analog' of Cayley's theorem
(\Cref{thm:groups:cayleys-theorem}):
\begin{prop}{}{prop:rings:Cayley-ring}
  Let $R$ be a ring. Then the function $r \mapsto \lambda_r$ is an
  injective ring homomorphism
  \[
    \lambda : R \to \End_{\categ{Ab}}(R).
  \]
\end{prop}
\begin{proof}
  For any $r \in R$ and for all $a, b \in R$, distributivity gives
  \[
    \lambda_r(a + b) = r(a + b) = ra + rb = \lambda_r(a) +
    \lambda_r(b):
  \]
  this shows that $\lambda_r$ is indeed an endomorphism of the group
  $(R, +)$, that is, $\lambda_r \in \End_{\categ{Ab}}(R)$.

  The function $\lambda : R \to \End_{\categ{Ab}}(R)$ defined by the
  assignment $r \mapsto \lambda_r$ is clearly injective, since if
  $r \neq s$, then $\lambda_r(1) = r \neq s = \lambda_s(1)$, so that
  $\lambda_r \neq \lambda_s$.

  We have to verify that $\lambda$ is in fact a homomorphism of
  rings. Recall that the addition in $\End_{\categ{Ab}}(R)$ is
  inherited from $R$ (cf.\ \S{}\ref{subsec:groups:homomorphisms-of-abelian-groups}): for all $r, s \in R$,
  $\lambda_r + \lambda_s$ is defined by
  \[
    (\forall a \in R) : (\lambda_r + \lambda_s)(a) = \lambda_r(a) +
    \lambda_s(a).
  \]
  Thus, the fact that $\lambda$ preserves addition is an immediate
  consequence of distributivity: for all $r, s, a \in R$,
  \[
    \lambda_{r+s}(a) = (r + s)a = ra + sa = \lambda_r(a) +
    \lambda_s(a).
  \]
  As for multiplication, it is associativity's turn to do its job:
  \[
    \lambda_{rs}(a) = (rs)a = r(sa) = r\lambda_s(a) =
    \lambda_r(\lambda_s(a)) = (\lambda_r \circ \lambda_s)(a).
  \]
  Of course $\lambda_1$ is the identity, completing the verification.
\end{proof}
The function $\mu : R \to \End_{\categ{Ab}}(R)$ defined by
$r \mapsto \mu_r$ is `almost' a ring homomorphism: the reader will
check (\Cref{exer:rings:right-mult-cayley}) that
\[
  \mu_{r+s} = \mu_r + \mu_s, \qquad \mu_{rs} = \mu_s \circ \mu_r,
\]
and $\mu_1 = \id_R$. That is, $\mu$ would in fact be a
ring homomorphism if we `reversed' multiplication\footnote{This issue
  is entirely analogous to the business of right-actions vs.\
  left-actions; cf.\ \S{}\ref{subsec:groups:transitive-actions-and-the-category-g-set}, especially \Cref{exer:groups:opposite-group}.} in
$R$. Of course this issue disappears if $R$ happens to be commutative
(since then $\lambda = \mu$).

\subsection*{Exercises}

\begin{exer}{\exerused}{exer:rings:zero-ring-domain}
  Prove that if there is a homomorphism from a zero-ring to a ring
  $R$, then $R$ is a zero-ring [\S{}\ref{subsec:rings:ring-homomorphisms}]
\end{exer}
\begin{exer}{}{exer:rings:additive-multiplicative-preserving}
  Let $R$ and $S$ be rings, and let $\varphi : R \to S$ be a function
  preserving both operations $+$, $\cdot$.
  \begin{itemize}
  \item Prove that if $\varphi$ is \emph{surjective}, then necessarily
    $\varphi(1_R) = 1_S$.
  \item Prove that if $\varphi \neq 0$ and $S$ is an integral domain,
    then $\varphi(1_R) = 1_S$.
  \end{itemize}
  (Therefore, in both cases $\varphi$ is in fact a ring homomorphism).
\end{exer}
\begin{exer}{}{exer:rings:powerset-vs-z2s}
  Let $S$ be a set, and consider the power set ring $\mathscr{P}(S)$
  (\Cref{exer:rings:power-set-ring}) and the ring $(\Z/2\Z)^S$ you constructed in
  \Cref{exer:rings:function-ring}. Prove that these two rings are isomorphic. (Cf.\
  \Cref{exer:prelim:power-set-power-two}.)
\end{exer}
\begin{exer}{}{exer:rings:quaternion-matrix-reps}
  Define functions $\field{H} \to \mathfrak{gl}_4(\R)$ and
  $\field{H} \to \mathfrak{gl}_2(\C)$ (cf.\ Exercises~1.4 and~1.12)
  by
  \[
    a + bi + cj + dk \mapsto \begin{pmatrix} a & b & c & d \\ -b & a &
      -d & c \\ -c & d & a & -b \\ -d & -c & b & a \end{pmatrix},
  \]
  \[
    a + bi + cj + dk \mapsto \begin{pmatrix} a+bi & c+di \\ -c+di &
      a-bi \end{pmatrix}
  \]
  for all $a, b, c, d \in \R$. Prove that both functions are injective
  ring homomorphisms. Thus, quaternions may be viewed as real or
  complex matrices.
\end{exer}
\begin{exer}{\lnot}{exer:rings:quaternion-norm}
  The \emph{norm} of a quaternion $w = a + bi + cj + dk$, with
  $a, b, c, d \in \R$, is the real number
  $N(w) = a^2 + b^2 + c^2 + d^2$.

  Prove that the function from the multiplicative group $\field{H}
^*$
  of nonzero quaternions to the multiplicative group $\R^+$ of
  positive real numbers, defined by assigning to each nonzero
  quaternion its norm, is a homomorphism. Prove that the kernel of
  this homomorphism is isomorphic to $\mathrm{SU}(2)$ (cf.\
  \Cref{exer:groups:su2-sphere}). [\ref{exer:rings:z-sqrt-d-subring-of-c}, \ref{exer:groups2:quaternions-as-semidirect-prod}, \ref{exer:irred:integral-quaternions-ring-basics}]
\end{exer}
\begin{exer}{\exerused}{exer:rings:poly-extension-property}
  Verify the `extension property' of polynomial rings, stated in
  \Cref{exam:rings:poly-universal-general}. [\S{}\ref{subsec:rings:universal-property-of-polynomial-rings}]
\end{exer}
\begin{exer}{\exerused}{exer:rings:z2-poly-function}
  Let $R = \Z/2\Z$, and let $f(x) = x^2 - x$; note $f(x) \neq 0$. What
  is the polynomial function $R \to R$ determined by $f(x)$? [\S{}\ref{subsec:rings:universal-property-of-polynomial-rings},
  \S{}\ref{subsec:irred:the-field-of-fractions-of-an-integral-domain}, \S{}\ref{subsec:irred:roots-and-reducibility}]
\end{exer}
\begin{exer}{}{exer:rings:subring-of-field-domain}
  Prove that every subring of a field is an integral domain.
\end{exer}
\begin{exer}{\lnot}{exer:rings:center-is-subring}
  The \emph{center} of a ring $R$ consists of the elements $a$ such
  that $ar = ra$ for all $r \in R$. Prove that the center is a subring
  of $R$.

  Prove that the center of a division ring is a field. [\ref{exer:rings:division-ring-p2-commutative}, \ref{exer:groups2:division-ring-ord-64-comm},
  \ref{exer:fields:wedderburn-thm-witts-proof}, \ref{exer:fields:wedderburn-thm-birkhoff-vandiver-proof}]
\end{exer}
\begin{exer}{\lnot}{exer:rings:centralizer-is-subring}
  The \emph{centralizer} of an element $a$ of a ring $R$ consists of
  the elements $r \in R$ such that $ar = ra$. Prove that the
  centralizer of $a$ is a subring of $R$, for every $a \in R$.

  Prove that the center of $R$ is the intersection of all its
  centralizers.

  Prove that every centralizer in a division ring is a division
  ring. [\ref{exer:rings:division-ring-p2-commutative}, \ref{exer:groups2:division-ring-ord-64-comm}, \ref{exer:fields:wedderburn-thm-birkhoff-vandiver-proof}]
\end{exer}
\begin{exer}{\lnot}{exer:rings:division-ring-p2-commutative}
  Let $R$ be a division ring consisting of $p^2$ elements, where $p$
  is a prime. Prove that $R$ is commutative, as follows:
  \begin{itemize}
  \item If $R$ is not commutative, then its center $C$ (\Cref{exer:rings:center-is-subring})
    is a proper subring of $R$. Prove that $C$ would then consist of
    $p$ elements.
  \item Let $r \in R$, $r \notin C$. Prove that the centralizer of $r$
    (\Cref{exer:rings:centralizer-is-subring}) contains both $r$ and $C$.
  \item Deduce that the centralizer of $r$ is the whole of $R$.
  \item Derive a contradiction, and conclude that $R$ had to be
    commutative (hence, a field).
  \end{itemize}
  This is a particular case of Wedderburn's theorem: every finite
  division ring is a field. [\ref{exer:groups2:division-ring-ord-64-comm}, \ref{exer:fields:wedderburn-thm-birkhoff-vandiver-proof}]
\end{exer}
\begin{exer}{\exerused}{exer:rings:z-in-q-cokernel}
  Consider the inclusion map $\iota : \Z \hookrightarrow \Q$. Describe
  the cokernel of $\iota$ in $\categ{Ab}$ and its cokernel in
  $\categ{Ring}$ (as defined by the appropriate universal property in
  the style of the one given in \S{}\ref{subsec:groups:epimorphisms-and-cokernels}). [\S{}\ref{subsec:rings:monomorphisms-and-epimorphisms}, \S{}\ref{sec:rings:modules-over-a-ring}]
\end{exer}
\begin{exer}{\exerused}{exer:rings:componentwise-product}
  Verify that the `componentwise' product $R_1 \times R_2$ of two
  rings satisfies the universal property for products in a category,
  given in \S{}\ref{subsec:prelims:products}. [\S{}\ref{subsec:rings:products}]
\end{exer}
\begin{exer}{\exerused}{exer:rings:poly-two-var-coproduct}
  Verify that $\Z[x_1, x_2]$ (along with the evident morphisms)
  satisfies the universal property for the coproduct of two copies of
  $\Z[x]$ in the category of \emph{commutative} rings. Explain why it
  does not satisfy it in $\categ{Ring}$. [\S{}\ref{subsec:rings:products}]
\end{exer}
\begin{exer}{\exerused}{exer:rings:z-mz-isomorphic-groups}
  For $m > 1$, the abelian groups $(\Z, +)$ and $(m\Z, +)$ are
  manifestly isomorphic: the function $\varphi : \Z \to m\Z$,
  $n \mapsto mn$ is a group isomorphism. Use this isomorphism to
  transfer the structure of `ring without identity' $(m\Z, +, \cdot)$
  back onto $\Z$: give an explicit formula for the `multiplication'
  $\bullet$ this defines on $\Z$ (that is, such that
  $\varphi(a \bullet b) = \varphi(a) \cdot \varphi(b)$). Explain why
  structures induced by different positive integers $m$ are
  nonisomorphic as `rings without 1'.

  (This shows that there are many different ways to give a structure
  of ring without identity to the group $(\Z, +)$. Compare this
  observation with \Cref{exer:rings:unique-ring-structure-on-z}.) [\S{}\ref{subsec:rings:ring-homomorphisms}]
\end{exer}
\begin{exer}{\exerused}{exer:rings:unique-ring-structure-on-z}
  Prove that there is (up to isomorphism) only one structure of ring
  \emph{with identity} on the abelian group $(\Z, +)$. (Hint: Let $R$
  be a ring whose underlying group is $\Z$. By
  \cref{prop:rings:Cayley-ring}, there is an injective ring
  homomorphism $\lambda : R \to \End_{\categ{Ab}}(R)$, and the latter
  is isomorphic to $\Z$ by \cref{prop:rings:endab-Z}. Prove that
  $\lambda$ is surjective.) [\S{}\ref{subsec:rings:ring-homomorphisms}, \ref{exer:rings:z-mz-isomorphic-groups}]
\end{exer}
\begin{exer}{\lnot}{exer:rings:endab-r}
  Let $R$ be a ring, and let $E = \End_{\categ{Ab}}(R)$ be the ring of
  endomorphisms of the underlying abelian group $(R, +)$. Prove that
  the center of $E$ is isomorphic to a subring of the center of
  $R$. (Prove that if $\alpha \in E$ commutes with all
  right-multiplications by elements of $R$, then $\alpha$ is
  left-multiplication by an element of $R$; then use
  \cref{prop:rings:Cayley-ring}.)
\end{exer}
\begin{exer}{\exerused}{exer:rings:right-mult-cayley}
  Verify the statements made about right-multiplication $\mu$,
  following \cref{prop:rings:Cayley-ring}. [\S{}\ref{subsec:rings:endab-g}]
\end{exer}
\begin{exer}{}{exer:rings:endab-znz}
  Prove that for $n \in \Z$ a positive integer,
  $\End_{\categ{Ab}}(\Z/n\Z)$ is isomorphic to $\Z/n\Z$ as a ring.
\end{exer}

\section{Ideals and quotient rings}
\label{sec:rings:ideals-and-quotient-rings}


\subsection{Ideals}
\label{subsec:rings:ideals}
In both $\categ{Set}$ and $\categ{Grp}$ we have been able to
`classify' surjective morphisms: in both cases, a surjective morphism
is, up to natural identifications, a quotient by a (suitable)
equivalence relation. In $\categ{Grp}$, we have seen that such
equivalence relations arise in fact from certain substructures: normal
subgroups.

The situation in $\categ{Ring}$ is analogous. We will establish a
canonical decomposition for rings, modeled after Theorems~I.2.7
and~II.8.1; the corresponding version of the `first isomorphism
theorem' (\Cref{coro:groups:first-isomorphism-theorem}) will identify every surjective ring
homomorphism with a quotient by a suitable substructure. The role of
normal subgroups will be played by \emph{ideals}.
\begin{defn}{}{defn:rings:ideal}
  Let $R$ be a ring. A subgroup $I$ of $(R, +)$ is a \emph{left-ideal}
  of $R$ if $rI \subseteq I$ for all $r \in R$; that is,
  \[
    (\forall r \in R)(\forall a \in I) : \quad ra \in I;
  \]
  it is a \emph{right-ideal} if $Ir \subseteq I$ for all $r \in R$;
  that is,
  \[
    (\forall r \in R)(\forall a \in I) : \quad ar \in I.
  \]
  A \emph{two-sided ideal} is a subgroup $I$ which is both a left- and
  a right-ideal.
\end{defn}
Of course in a commutative ring there is no distinction between left-
and right-ideals. Even in the general setting, we will almost
exclusively be concerned with two-sided ideals; thus I will omit
qualifiers, and an \emph{ideal} of a ring will implicitly be
two-sided.
\begin{rem}{}{}
  As seen in \S{}\ref{subsec:rings:monomorphisms-and-epimorphisms}, a subring of $R$ is a subset $S \subseteq R$
  which contains $1_R$ and satisfies the ring axioms with respect to
  the operations $+$, $\cdot$ induced from $R$. Ideals are close to
  being subrings: they are subgroups, and they are closed with respect
  to multiplication. But the only ideal of a ring $R$ containing $1_R$
  is $R$ itself: this is an immediate consequence of the `absorption
  properties' stated in \cref{defn:rings:ideal}. Thus ideals are in
  general not subrings; they are `rngs'.
\end{rem}
Ideals are considerably more important than subrings in the
development of the theory of rings.\footnote{Arguably, the reason is
  that ideals are precisely the submodules of a ring $R$.} Of course
the image of a ring homomorphism is necessarily a subring of the
target; but a lesson learned from \S{}\ref{subsec:groups:canonical-decomposition} and following is that
kernels really capture the structure of a homomorphism, and kernels
are ideals:
\begin{exam}{}{exam:rings:kernel-is-ideal}
  Let $\varphi : R \to S$ be any ring homomorphism. Then
  $\ker \varphi$ is an ideal of $R$.

  Indeed, we know already that $\ker \varphi$ is a subgroup; we have
  to verify the absorption properties. These are an immediate
  consequence of \cref{lem:rings:zero-mult}: for all $r \in R$, all
  $a \in \ker \varphi$, we have
  \[
    \varphi(ra) = \varphi(r)\varphi(a) = \varphi(r) \cdot 0 = 0,
    \qquad \varphi(ar) = \varphi(a)\varphi(r) = 0 \cdot \varphi(r) =
    0.
  \]
  More generally, it is easy to verify that the inverse image of an
  ideal is an ideal (\Cref{exer:rings:preimage-ideal}) and $\{0_S\}$ is clearly an ideal
  of $S$.

  Similarly to the situation with normal subgroups in the context of
  groups, we will soon see that `kernels of ring homomorphisms' and
  `ideals' are in fact equivalent concepts.
\end{exam}

\subsection{Quotients}
\label{subsec:rings:quotients}
Let $I$ be a subgroup of the abelian group $(R, +)$ of a ring
$R$. Subgroups of abelian groups are automatically normal, so we have
a quotient \emph{group} $R/I$, whose elements are the cosets of $I$:
\[
  r + I
\]
(written, of course, in additive notation). Further, we have a
surjective \emph{group} homomorphism
\[
  \pi : R \to \frac{R}{I}, \quad r \mapsto r + I.
\]

As we have explored in great detail for groups, this construction
satisfies a suitable universal property with respect to group
homomorphisms (\Cref{thm:groups:quotient-universal-property}). \emph{Of course} we are now going to
ask under what circumstances this construction can be performed in
$\categ{Ring}$, satisfying the analogous universal property with
respect to ring homomorphisms.  That is, what should we ask of $I$, in
order to have a ring structure on $R/I$, so that $\pi$ becomes a ring
homomorphism? Go figure this out for yourself, before reading ahead!
As so often is the case, the requirement is precisely what tells us
the answer. Since $\pi$ will have to be a ring homomorphism, the
multiplication in $R/I$ \emph{must} be as follows: for all $(a + I)$,
$(b + I)$ in $R/I$,
\[
  (a + I) \cdot (b + I) = \pi(a) \cdot \pi(b) \stackrel{?}{=} \pi(ab)
  = ab + I,
\]
where the equality $\stackrel{?}{=}$ is forced if $\pi$ is to be a
ring homomorphism.  This says that there is only one sensible ring
structure on $R/I$, given by
\[
  (\forall a, b \in R) : \quad (a + I)(b + I) := ab + I.
\]
The reader should realize right away that \emph{if} this operation is
well-defined, then it does make $R/I$ into a ring: associativity will
be inherited from the associativity in $R$, and the identity will
simply be the coset $1 + I$ of $1$.

So all is well, if the proposed operation is well-defined. \emph{But}
is it going to be well-defined?
\begin{exam}{}{}
  It need not be, if $I$ is an arbitrary subgroup of $R$. For example,
  take $\Z$ as a subgroup of $\Q$; then
  \[
    0 + \Z = 1 + \Z
  \]
  $(= \Z)$ as elements of the group $\Q/\Z$, and $\tfrac{1}{2} + \Z$
  is another coset, yet
  \[
    0 \cdot \tfrac{1}{2} + \Z = \Z \neq \tfrac{1}{2} + \Z = 1 \cdot
    \tfrac{1}{2} + \Z.
  \]
\end{exam}
Assume then that the operation \emph{is} well-defined, so that $R/I$
is a ring and $\pi : R \to R/I$ is a ring homomorphism. What does this
say about $I$? Answer: $I$ is the kernel of $R \to R/I$, so
necessarily $I$ must be an ideal, as seen in \Cref{exam:rings:kernel-is-ideal}!  Conversely,
let us assume $I$ is an ideal of $R$, and verify that the proposed
prescription for the operation in $R/I$ is well-defined. For this,
suppose
\[
  a' + I = a'' + I \quad \text{and} \quad b' + I = b'' + I;
\]
recall that this means that $a'' - a' \in I$, $b'' - b' \in I$; then
\[
  a''b'' - a'b' = a''b'' - a''b' + a''b' - a'b' = a''(b'' - b') + (a''
  - a')b' \in I,
\]
using both the left-absorption and right-absorption properties of
\cref{defn:rings:ideal}. This says precisely that
$a'b' + I = a''b'' + I$, proving that the operation is well-defined.
Summarizing, we have verified that $R/I$ is a ring, in such a way that
the canonical projection $\pi : R \to R/I$ is a \emph{ring}
homomorphism, if and only if $I$ is an \emph{ideal} of $R$.
\begin{defn}{}{defn:rings:quotient-ring}
  This ring $R/I$ is called the \emph{quotient ring} of $R$
  \emph{modulo} $I$.
\end{defn}
\begin{exam}{}{}
  We know that all subgroups of $(\Z, +)$ are of the form $n\Z$ for a
  nonnegative integer $n$ (\Cref{prop:groups:subgroups-of-z}). It is immediately
  verified that all subgroups of $\Z$ are in fact \emph{ideals} of the
  ring $(\Z, +, \cdot)$. The quotients $\Z/n\Z$ are of course nothing
  but the rings so-denoted in \S{}\ref{subsec:rings:first-examples-and-special-classes-of-rings} (and earlier).
\end{exam}
The fact that $\Z$ is initial in $\categ{Ring}$ now prompts a natural
definition. For a ring $R$, let $f : \Z \to R$ be the unique ring
homomorphism, defined by $a \mapsto a \cdot 1_R$. Then $\ker f = n\Z$
for a well-defined nonneg­ative integer $n$ determined by $R$.

\begin{defn}{}{defn:rings:characteristic}
  The \emph{characteristic} of $R$ is this nonneg­ative integer $n$.
\end{defn}

Thus, the characteristic of $R$ is $n > 0$ if the order of $1_R$
\emph{as an element of} $(R, +)$ is a positive integer $n$, while the
characteristic is $0$ if the order of $1_R$ is $\infty$.  I have now
fulfilled my promise to identify the notions of ideal and kernel of
ring homomorphisms: every kernel is an ideal (cf.\ \Cref{exam:rings:kernel-is-ideal}); and
on the other hand every ideal $I$ is the kernel of the ring
homomorphisms $R \to R/I$. We have the slogan:
\[
  \text{kernel} \iff \text{ideal}
\]
in the context of ring theory.  The key universal property holds in
this context as it does for groups; cf.\ \Cref{thm:groups:quotient-universal-property}. I reproduce
the statement here for reference, but the reader should realize that
very little needs to be proven at this point: the needed (group)
homomorphism exists and is unique by \Cref{thm:groups:quotient-universal-property}, and verifying it
is a \emph{ring} homomorphism is immediate.
\begin{thm}{}{thm:rings:quotient-universal}
  Let $I$ be a two-sided ideal of a ring $R$. Then for every ring
  homomorphism $\varphi : R \to S$ such that
  $I \subseteq \ker \varphi$ there exists a unique ring homomorphism
  $\tilde{\varphi} : R/I \to S$ so that the diagram
  \[
    \begin{tikzcd}
      R \arrow[r, "\varphi"] \arrow[d, "\pi"'] & S \\
      R/I \arrow[ur, "\exists!\tilde{\varphi}"']
    \end{tikzcd}
  \]
  commutes.
\end{thm}

As a reminder to the lazy reader, $\tilde{\varphi}$ is defined by
\[
  \tilde{\varphi}(r + I) := \varphi(r);
\]
(part of) the content of the theorem is that this function is
well-defined (if $I \subseteq \ker \varphi$), and it is a ring
homomorphism.

\subsection{Canonical decomposition and consequences}
\label{subsec:rings:canonical-decomposition-and-consequences}
As the reader should now expect, \cref{thm:rings:quotient-universal}
is the key element in a standard decomposition of every ring
homomorphism. This is entirely analogous to the decomposition of
set-functions studied in \S{}\ref{subsec:prelims:canonical-decomposition} and the decomposition of group
homomorphisms obtained in \S{}\ref{subsec:groups:canonical-decomposition}. Here is the statement:
\begin{thm}{}{thm:rings:canonical-decomp}
  Every ring homomorphism $\varphi : R \to S$ may be decomposed as
  follows:
  \[
    \begin{tikzcd}
      R \arrow[r, two heads] \arrow[rrr, bend left, "\varphi"] &
      R/\ker\varphi \arrow[r, "\tilde{\varphi}"', "\sim"] &
      \image\varphi \arrow[r, hook] & S
    \end{tikzcd}
  \]
  where the isomorphism $\tilde{\varphi}$ in the middle is the
  homomorphism induced by $\varphi$ (as in
  \cref{thm:rings:quotient-universal}).
\end{thm}
The reader will realize that this statement requires no proof at this
point: the decomposition holds at the level of groups (by
\Cref{thm:groups:canonical-decomposition}) and the maps are all ring homomorphisms as observed
earlier in this section.

The `first isomorphism theorem' for rings is an immediate corollary:
\begin{coro}{}{coro:rings:first-iso}
  Suppose $\varphi : R \to S$ is a surjective ring homomorphism. Then
  $S \cong R/\ker\varphi$.
\end{coro}
As in the group case, the reader should develop the healthy
instinctive reaction of viewing every surjective homomorphism of rings
as a quotient (by an ideal), up to a natural identification.

Equally instinctive should be the realization that ideals of a
quotient $R/I$ are in one-to-one correspondence with ideals of $R$
containing $I$. Once more, not a whole lot is new here: we already
know (cf.\ \S{}\ref{subsec:groups:subgroups-of-quotients}) that the function
\[
  u : \{\text{subgroups } J \text{ of } R \text{ containing } I\} \to
  \{\text{subgroups of } R/I\}
\]
defined by $u(J) = J/I$ is a bijection preserving inclusions; it takes
a moment to check that $J/I$ is an ideal of $R/I$ if and only if $J$
is an ideal of $R$. The main content of this observation is best
packaged in the corresponding `third isomorphism theorem':
\begin{prop}{}{prop:rings:third-iso}
  Let $I$ be an ideal of a ring $R$, and let $J$ be an ideal of $R$
  containing $I$. Then $J/I$ is an ideal of $R/I$, and
  \[
    \frac{R/I}{J/I} \cong \frac{R}{J}.
  \]
\end{prop}
\begin{proof}
  Since $I \subseteq J = \ker(R \to R/J)$, we have an induced ring
  homomorphism $\varphi : R/I \to R/J$ by
  \cref{thm:rings:quotient-universal}. Explicitly,
  $\varphi(r + I) = r + J$; $\varphi$ is manifestly surjective. Since
  \[
    \ker \varphi = \{r + I \mid \varphi(r + I) = J\} = \{r + I \mid r
    + J = J\} = \{r + I \mid r \in J\} = J/I,
  \]
  we see that $J/I$ is an ideal (since it is a kernel), and the stated
  isomorphism follows from \cref{coro:rings:first-iso}.
\end{proof}
What about the `second' isomorphism theorem? This would be a relation
between the ideals
\[
  \frac{I + J}{I}, \qquad \frac{J}{I \cap J}
\]
of the rings $R/I$, $R/(I \cap J)$, respectively, assuming $I$ and $J$
are ideals of $R$ (and where it is immediately checked that $I + J$
and $I \cap J$ are indeed ideals\footnote{The notation $I + J$ should
  not surprise the reader: $I$, $J$ are subgroups of the abelian group
  $R$, so we know what it means to `add' them; cf.\ \S{}\ref{subsec:groups:normal-subgroups}.}). The
reader may want to go back to \S{}\ref{subsec:groups:hk-h-vs-k-h-cap-k} for the version of this story
for groups. My feeling is that $\categ{Ring}$ is not the best place to
play this game, since $(I+J)/I$, $J/(I \cap J)$ are not even rings
according to my conventions. \emph{Modules} will be a more natural
context for this result.

In any case, the reader will benefit the most from exploring the
matter on his/her own; cf.\ \Cref{exer:rings:third-iso-ideals}.

\subsection*{Exercises}

\begin{exer}{}{exer:rings:image-is-subring}
  Prove that the image of a ring homomorphism $\varphi : R \to S$ is a
  subring of $S$. What can you say about $\varphi$ if its image is an
  ideal of $S$? What can you say about $\varphi$ if its kernel is a
  subring of $R$?
\end{exer}
\begin{exer}{\exerused}{exer:rings:preimage-ideal}
  Let $\varphi : R \to S$ be a ring homomorphism, and let $J$ be an
  ideal of $S$. Prove that $I = \varphi^{-1}(J)$ is an ideal of
  $R$. [\S{}\ref{subsec:rings:ideals}]
\end{exer}
\begin{exer}{\lnot}{exer:rings:image-ideal-not-necessarily}
  Let $\varphi : R \to S$ be a ring homomorphism, and let $J$ be an
  ideal of $R$.
  \begin{itemize}
  \item Show that $\varphi(J)$ need not be an ideal of $S$.
  \item Assume that $\varphi$ is surjective; then prove that
    $\varphi(J)$ is an ideal of $S$.
  \item Assume that $\varphi$ is surjective, and let
    $I = \ker\varphi$; thus we may identify $S$ with $R/I$. Let
    $\overline{J} = \varphi(J)$, an ideal of $R/I$ by the previous
    point. Prove that
    $\frac{R/I}{\overline{J}} \cong \frac{R}{I + J}$.
  \end{itemize}
  (Of course this is just a rehash of \cref{prop:rings:third-iso}.)
  [\ref{exer:rings:principal-ideal-with-a}]
\end{exer}
\begin{exer}{}{exer:rings:all-subgroups-ideals-znz}
  Let $R$ be a ring such that every subgroup of $(R, +)$ is in fact an
  ideal of $R$. Prove that $R \cong \Z/n\Z$, where $n$ is the
  characteristic of $R$.
\end{exer}
\begin{exer}{\lnot}{exer:rings:matrix-ideal-entries}
  Let $J$ be a two-sided ideal of the ring $\mathcal{M}_n(R)$ of
  $n \times n$ matrices over a ring $R$. Prove that a matrix
  $A \in \mathcal{M}_n(R)$ belongs to $J$ if and only if the matrices
  obtained by placing any entry of $A$ in any position, and $0$
  elsewhere, belong to $J$. (Hint: Carefully contemplate the operation
  \[
    \begin{pmatrix} 0 & 0 & 0 \\ 0 & 0 & 0 \\ 1 & 0 & 0 \end{pmatrix}
    \begin{pmatrix} a & b & c \\ d & e & f \\ g & h & i \end{pmatrix}
    \begin{pmatrix} 0 & 0 & 0 \\ 1 & 0 & 0 \\ 0 & 0 & 0 \end{pmatrix}
    = \begin{pmatrix} 0 & 0 & 0 \\ 0 & 0 & 0 \\ b & 0 &
      0 \end{pmatrix}.)
  \] [\ref{exer:rings:matrix-ideal-entries-converse}]
\end{exer}
\begin{exer}{\lnot}{exer:rings:matrix-ideal-entries-converse}
  Let $J$ be a two-sided ideal of the ring $\mathcal{M}_n(R)$ of
  $n \times n$ matrices over a ring $R$, and let $I \subseteq R$ be
  the set of $(1,1)$ entries of matrices in $J$. Prove that $I$ is a
  two-sided ideal of $R$ and $J$ consists precisely of those matrices
  whose entries all belong to $I$. (Hint: \Cref{exer:rings:matrix-ideal-entries}.) [\ref{exer:rings:division-ring-two-sided-counterexample}]
\end{exer}
\begin{exer}{}{exer:rings:ra-ar-ideals}
  Let $R$ be a ring, and let $a \in R$. Prove that $Ra$ is a
  left-ideal of $R$ and $aR$ is a right-ideal of $R$. Prove that $a$
  is a left-, resp.\ right-, unit if and only if $R = aR$, resp.\
  $R = Ra$.
\end{exer}
\begin{exer}{\exerused}{exer:rings:division-ring-ideals}
  Prove that a ring $R$ is a division ring if and only if its only
  left-ideals and right-ideals are $\{0\}$ and $R$.

  In particular, a commutative ring $R$ is a field if and only if the
  only ideals of $R$ are $\{0\}$ and $R$. [\ref{exer:rings:division-ring-two-sided-counterexample}, \S{}\ref{subsec:rings:prime-and-maximal-ideals}]
\end{exer}
\begin{exer}{\lnot}{exer:rings:division-ring-two-sided-counterexample}
  Counterpoint to \Cref{exer:rings:division-ring-ideals}: It is not true that a ring $R$ is a
  division ring if and only if its only two-sided ideals are $\{0\}$
  and $R$. A nonzero ring with this property is said to be
  \emph{simple}; by \Cref{exer:rings:division-ring-ideals}, fields are the only simple
  commutative rings.

  Prove that $\mathcal{M}_n(\R)$ is simple. (Use \Cref{exer:rings:matrix-ideal-entries-converse}.) [\ref{exer:rings:maximal-iff-simple-quotient}]
\end{exer}
\begin{exer}{\exerused}{exer:rings:field-hom-injective}
  Let $\varphi : k \to R$ be a ring homomorphism, where $k$ is a field
  and $R$ is a nonzero ring. Prove that $\varphi$ is
  injective. [\S{}\ref{subsec:irred:the-field-of-fractions-of-an-integral-domain}, \S{}\ref{subsec:irred:adding-roots-algebraically-closed-fields}]
\end{exer}
\begin{exer}{}{exer:rings:c-subring-no-hom-to-r}
  Let $R$ be a ring containing $\C$ as a subring. Prove that there are
  no ring homomorphisms $R \to \R$.
\end{exer}
\begin{exer}{\exerused}{exer:rings:nilpotents-form-ideal}
  Let $R$ be a commutative ring. Prove that the set of nilpotent
  elements of $R$ is an ideal of $R$. (Cf.\ \Cref{exer:rings:nilpotent}. This ideal
  is called the \emph{nilradical} of $R$.)

  Find a noncommutative ring in which the set of nilpotent elements is
  not an ideal. [\ref{exer:rings:nilradical-quotient-reduced}, \ref{exer:rings:nilradical-in-every-prime}, \ref{exer:irred:nilradical-subset-intersection-of-primes}, \S{}\ref{subsec:fields:a-little-affine-algebraic-geometry}]
\end{exer}
\begin{exer}{\lnot}{exer:rings:nilradical-quotient-reduced}
  Let $R$ be a commutative ring, and let $N$ be its nilradical (cf.\
  \Cref{exer:rings:nilpotents-form-ideal}). Prove that $R/N$ contains no nonzero nilpotent
  elements. (Such a ring is said to be \emph{reduced}.) [\ref{exer:rings:ij-reduced-implies-i-cap-j}, \ref{exer:fields:radical-ideal-nilradical-correspondence}]
\end{exer}
\begin{exer}{\lnot}{exer:rings:domain-characteristic-0-or-prime}
  Prove that the characteristic of an integral domain is either $0$ or
  a prime integer. Do you know any ring of characteristic $1$?
  [\ref{exer:irred:prime-subfield-char-0-or-p}]
\end{exer}
\begin{exer}{\lnot}{exer:rings:boolean-power-set}
  A ring $R$ is \emph{Boolean}\footnote{After George Boole.} if
  $a^2 = a$ for all $a \in R$. Prove that $\mathcal{P}(S)$ is Boolean,
  for every set $S$ (cf.\ \Cref{exer:rings:power-set-ring}). Prove that every nonzero
  Boolean ring is commutative and has characteristic $2$. Prove that
  if an integral domain $R$ is Boolean, then $R \cong \Z/2\Z$. [\ref{exer:rings:krull-dimension-zero},
  \ref{exer:irred:finite-boolean-ring-is-product-of-z2}]
\end{exer}
\begin{exer}{\lnot}{exer:rings:subsets-of-t-form-ideal}
  Let $S$ be a set and $T \subseteq S$ a subset. Prove that the
  subsets of $S$ contained in $T$ form an ideal of the power set ring
  $\mathcal{P}(S)$. Prove that if $S$ is finite, then every ideal of
  $\mathcal{P}(S)$ is of this form. For $S$ infinite, find an ideal of
  $\mathcal{P}(S)$ that is not of this form. [\ref{exer:irred:power-set-ring-noeth-iff}]
\end{exer}
\begin{exer}{\exerused}{exer:rings:third-iso-ideals}
  Let $I$, $J$ be ideals of a ring $R$. State and prove a precise
  result relating the ideals $(I + J)/I$ of $R/I$ and $J/(I \cap J)$
  of $R/(I \cap J)$. [\S{}\ref{subsec:rings:canonical-decomposition-and-consequences}]
\end{exer}

\section{Ideals and quotients: Remarks and examples. Prime and maximal
  ideals}
\label{sec:rings:ideals-and-quotients-remarks-and-examples-prime-and}

\subsection{Basic operations}
\label{subsec:rings:basic-operations}
In the commutative case, these two subsets coincide and are denoted
$(a)$. This is the \emph{principal ideal generated by $a$}. For
example, the zero-ideal $\{0\} = (0)$ and the whole ring $R = (1)$ are
both principal ideals.

In general (\Cref{exer:rings:family-of-ideals-sum}), if $\{I_\alpha\}_{\alpha \in A}$ is a
family of ideals of a ring $R$, then the sum $\sum_\alpha I_\alpha$ is
an ideal of $R$. If $a_\alpha$ is any collection of elements of a
commutative ring $R$, then
\[
  (a_\alpha)_{\alpha \in A} := \sum_{\alpha \in A} (a_\alpha)
\]
is the ideal \emph{generated by the elements $a_\alpha$}. In
particular,
\[
  (a_1, \ldots, a_n) = (a_1) + \cdots + (a_n)
\]
is the smallest ideal of $R$ containing $a_1, \ldots, a_n$; this ideal
consists of the elements of $R$ that may be written as
$r_1 a_1 + \cdots + r_n a_n$ for $r_1, \ldots, r_n \in R$. An ideal
$I$ of a commutative ring $R$ is \emph{finitely generated} if
$I = (a_1, \ldots, a_n)$ for some $a_1, \ldots, a_n \in R$.

\begin{exam}{}{exam:rings:calculus-of-ideals}
  It is a good idea to get used to a bit of `calculus' of ideals and
  quotients in terms of generators; judicious use of the isomorphism
  theorems yields convenient statements. For example, let $R$ be a
  commutative ring, and let $a, b \in R$; denote by $\bar{b}$ the
  class of $b$ in $R/(a)$. Then
  \[
    (R/(a))/(\bar{b}) \cong R/(a, b).
  \]
  Indeed, this is a particular case of \cref{prop:rings:third-iso},
  since
  \[
    (\bar{b}) = \frac{(a,b)}{(a)}
  \]
  as ideals of $R/(a)$.
\end{exam}

Note that principal ideals are (very special) finitely generated
ideals. These notions are so important that we give special names to
rings in which they are satisfied by every ideal.

\begin{defn}{Noetherian ring}{defn:rings:noetherian}
  A commutative ring $R$ is \emph{Noetherian} if every ideal of $R$ is
  finitely generated.
\end{defn}

\begin{defn}{PID}{defn:rings:PID}
  An integral domain $R$ is a \emph{PID} (`Principal Ideal Domain') if
  every ideal of $R$ is principal.
\end{defn}

Thus, PIDs are (very special) Noetherian rings. In due time we will
deal at length with these classes of rings (cf.\ \Cref{chap:irred});
Noetherian rings are very important in number theory and algebraic
geometry.

The reader is already familiar with an important PID:

\begin{prop}{}{prop:rings:Z-PID}
  $\Z$ is a PID.
\end{prop}
\begin{proof}
  Let $I \subseteq \Z$ be an ideal. Since $I$ is a subgroup, $I = n\Z$
  for some $n \in \Z$, by \Cref{prop:groups:subgroups-of-z}. Since $n\Z = (n)$, this
  shows that $I$ is principal.
\end{proof}

The fact that $\Z$ is a PID captures precisely `why' greatest common
divisors behave as they do in $\Z$: if $m$, $n$ are integers, then the
ideal $(m, n)$ must be principal, and hence
\[
  (m, n) = (d)
\]
for some (positive) integer $d$. This integer is manifestly the gcd of
$m$ and $n$: since $m \in (d)$ and $n \in (d)$, then $d \mid m$ and
$d \mid n$, etc.

If $k$ is a field, the ring of polynomials $k[x]$ is also a PID;
proving this is easy, using the `division with remainder' that we will
run into very soon (\S{}\ref{subsec:rings:quotients-of-polynomial-rings}); the reader should work this out on
his/her own (\Cref{exer:rings:field-poly-ring-pid}) now. This fact will be absorbed in the
general theory when we review the general notion of `Euclidean
domain', in \S{}\ref{subsec:irred:euclidean-domain-implies-pid}.

By contrast, the ring $\Z[x]$ is not a PID: indeed, the reader should
be able to verify that the ideal $(2, x)$ cannot be generated by a
single element. As we will see in due time, greatest common divisors
make good sense in a ring such as $\Z[x]$, but the matter is a little
more delicate, since this ring is not a PID\footnote{It is, however, a
  `UFD', that is, a `unique factorization domain'. This suffices for a
  good notion of gcd; cf.\ \S{}\ref{subsec:irred:irreducible-factors-and-greatest-common-divisor}.}.

There are several more basic operations involving ideals; for now, the
following two will suffice.
\begin{itemize}
\item Again assume that $\{I_\alpha\}_{\alpha \in A}$ is a collection
  of ideals of a ring $R$. Then the intersection
  $\bigcap_{\alpha \in A} I_\alpha$ is (clearly) an ideal of $R$; it
  is the largest ideal contained in all of the ideals $I_\alpha$.
\item If $I$, $J$ are ideals of $R$, then $IJ$ denotes the ideal
  \emph{generated} by all products $ij$ with $i \in I$, $j \in
  J$. More generally, if $I_1, \ldots, I_n$ are ideals in $R$, then
  the `product' $I_1 \cdots I_n$ denotes the ideal generated by all
  products $i_1 \cdots i_n$ with $i_k \in I_k$.
\end{itemize}
The reader should note the clash of notation: in the context of groups
(especially \S{}\ref{subsec:groups:hk-h-vs-k-h-cap-k}) $IJ$ would mean \emph{something else}. Watch
out!

It is clear that $IJ \subseteq I \cap J$: every element $ij$ with
$i \in I$ and $j \in J$ is in $I$ (because $I$ is a right-ideal) and
in $J$ (because $J$ is a left-ideal); therefore $I \cap J$ contains
all products $ij$, and hence it must contain the ideal $IJ$ they
generate. Sometime the product agrees with the intersection:
\[
  (4) \cap (3) = (12) = (4) \cdot (3) \quad \text{in } \Z;
\]
and sometime it does not:
\[
  (4) \cap (6) = (12) \neq (24) = (4) \cdot (6).
\]
The matter of whether $IJ = I \cap J$ is often subtle; a prototype
situation in which this equality holds is given in \Cref{exer:rings:coprime-ideals-product-intersection}.

\subsection{Quotients of polynomial rings}
\label{subsec:rings:quotients-of-polynomial-rings}

$\Z/n\Z$ is our familiar ring of congruence classes modulo
$n$. Quotients of polynomial rings by principal ideals are a good
source of `concrete', but maybe less familiar, examples.

Let $R$ be a (nonzero) ring, and let
\[
  f(x) = x^d + a_{d-1}x^{d-1} + \cdots + a_1 x + a_0 \in R[x]
\]
be a polynomial; for convenience, I am assuming that $f(x)$ is
\emph{monic}, that is, its leading coefficient (the coefficient of the
highest power of $x$ appearing in $f(x)$) is $1$. In terms of ideals,
this is not a serious requirement if the coefficient ring $R$ is a
field (\Cref{exer:rings:field-poly-ideal-monic-generator}), but it may be substantial otherwise: for
example, $(2x) \subseteq \Z[x]$ cannot be generated by a monic
polynomial. Also note that a monic polynomial is necessarily a
non-zero-divisor (\Cref{exer:rings:monic-poly-not-zero-divisor}) and that if $f(x)$ is monic, then
$\deg(f(x)q(x)) = \deg f(x) + \deg q(x)$ for all polynomials $q(x)$.

It is convenient to assume that $f(x)$ is monic because we can then
\emph{divide} by $f(x)$, with remainder. That is, if $g(x) \in R[x]$
is another polynomial, then there exist polynomials
$q(x), r(x) \in R[x]$ such that
\[
  g(x) = f(x)q(x) + r(x)
\]
and\footnote{Note: With the convention that the degree of the
  polynomial $0$ is $-\infty$, the condition $\deg r(x) < \deg f(x)$
  is satisfied by $r(x) = 0$.} $\deg r(x) < \deg f(x)$. This is simply
the process of `long division' of polynomials, which is surely
familiar to the reader, and can be performed over any ring when
dividing by monic polynomials\footnote{The key point is that if
  $n > d$, then for all $a \in R$ we have $ax^n = ax^{n-d}f(x) + h(x)$
  for some polynomial $h(x)$ of degree $< n$. Arguing inductively,
  this shows that we may perform division by $f(x)$ with remainder for
  all `monomials' $ax^n$, and hence (by linearity) for all polynomials
  $g(x) \in R[x]$.}.

The situation appears then to be similar to the situation in $\Z$,
where we also have division with remainder. Quotients and remainders
are uniquely\footnote{This assertion has to be taken with a grain of
  salt in the noncommutative case, as different quotients and
  remainders may arise if we divide `on the left' rather than `on the
  right'.} determined by $g(x)$ and $f(x)$:

\begin{lem}{}{lem:rings:poly-division-unique}
  Let $f(x)$ be a monic polynomial, and assume
  \[
    f(x)q_1(x) + r_1(x) = f(x)q_2(x) + r_2(x)
  \]
  with both $r_1(x)$ and $r_2(x)$ polynomials of degree $< \deg
  f(x)$. Then $q_1(x) = q_2(x)$ and $r_1(x) = r_2(x)$.
\end{lem}
\begin{proof}
  Indeed, we have
  \[
    f(x)(q_1(x) - q_2(x)) = r_2(x) - r_1(x);
  \]
  if $r_2(x) \neq r_1(x)$, then $r_2(x) - r_1(x)$ has degree
  $< \deg f(x)$, while $f(x)(q_1(x) - q_2(x))$ has degree
  $\geq \deg f(x)$, giving a contradiction. Therefore
  $r_1(x) = r_2(x)$, and $q_1(x) = q_2(x)$ follows right away since
  monic polynomials are non-zero-divisors.
\end{proof}

The preceding considerations may be summarized in a rather efficient
way in the language of ideals and cosets. We will now restrict
ourselves to the \emph{commutative} case, mostly for notational
convenience, but also because this will guarantee that ideals are
two-sided ideals, so that quotients are defined as rings (cf.\
\S{}\ref{subsec:rings:quotients}).

Assume then that $R$ is a commutative ring. What we have shown is
that, if $f(x)$ is monic, then for every $g(x) \in R[x]$ there exists
a unique polynomial $r(x)$ of degree $< \deg f(x)$ and such that
\[
  g(x) + (f(x)) = r(x) + (f(x))
\]
as cosets of the principal ideal $(f(x))$ in $R[x]$.

Refining this observation leads to a useful group-theoretic
statement. Note that polynomials of degree $< d$ may be seen as
elements of a direct sum
\[
  R^{\oplus d} = \underbrace{R \oplus \cdots \oplus R}_{d \text{
      times}};
\]
indeed, the function $\psi : R^{\oplus d} \to R[x]$ defined by
\[
  \psi((r_0, r_1, \ldots, r_{d-1})) = r_0 + r_1 x + \cdots + r_{d-1}
  x^{d-1}
\]
is clearly an injective homomorphism of abelian groups, hence an
isomorphism onto its image, and this consists precisely of the
polynomials of degree $< d$. I will glibly identify $R^{\oplus d}$
with this set of polynomials for the purpose of the discussion that
follows.

The next result may be seen as a way to concoct many different and
interesting ring structures on the direct sum $R^{\oplus d}$:

\begin{prop}{}{prop:rings:poly-group-iso}
  Let $R$ be a commutative ring, and let $f(x) \in R[x]$ be a monic
  polynomial of degree $d$. Then the function
  $\varphi : R[x] \to R^{\oplus d}$ defined by sending $g(x) \in R[x]$
  to the remainder of the division of $g(x)$ by $f(x)$ induces an
  isomorphism of abelian groups
  \[
    \frac{R[x]}{(f(x))} \cong R^{\oplus d}.
  \]
\end{prop}
\begin{proof}
  The given function $\varphi$ is well-defined by
  \cref{lem:rings:poly-division-unique}, and it is surjective since it
  has a right inverse (that is, the function
  $\psi : R^{\oplus d} \to R[x]$ defined above).

  I claim that $\varphi$ is a homomorphism of abelian groups. Indeed,
  if
  \[
    g_1(x) = f(x)q_1(x) + r_1(x) \quad \text{and} \quad g_2(x) =
    f(x)q_2(x) + r_2(x)
  \]
  with $\deg r_1(x) < d$, $\deg r_2(x) < d$, then
  \[
    g_1(x) + g_2(x) = f(x)(q_1(x) + q_2(x)) + (r_1(x) + r_2(x))
  \]
  and $\deg(r_1(x) + r_2(x)) < d$: this implies (again by
  \cref{lem:rings:poly-division-unique})
  \[
    \varphi(g_1(x) + g_2(x)) = r_1(x) + r_2(x) = \varphi(g_1(x)) +
    \varphi(g_2(x)).
  \]
  By the first isomorphism theorem for abelian groups, then, $\varphi$
  induces an isomorphism
  \[
    \frac{R[x]}{\ker\varphi} \cong R^{\oplus d}.
  \]
  On the other hand, $\varphi(g(x)) = 0$ if and only if
  $g(x) = f(x)q(x)$ for some $q(x) \in R[x]$, that is, if and only if
  $g(x)$ is in the principal ideal generated by $f(x)$. This shows
  $\ker\varphi = (f(x))$, concluding the proof.
\end{proof}

\begin{exam}{}{exam:rings:remainder-degree-1}
  Assume $f(x)$ is monic of degree $1$: $f(x) = x - a$ for some
  $a \in R$. Then the remainder of $g(x)$ after division by $f(x)$ is
  simply the `evaluation' $g(a)$ (cf.\ \Cref{exam:rings:poly-universal-general}): indeed,
  \[
    g(x) = (x - a)q(x) + r
  \]
  for some $r \in R$ (the remainder must have degree $< 1$; hence it
  is a constant); evaluating at $a$ gives
  \[
    g(a) = (a - a)q(a) + r = 0 \cdot q(a) + r = r
  \]
  as claimed. In particular, $g(a) = 0$ if and only if
  $g(x) \in (x - a)$.

  The content of \cref{prop:rings:poly-group-iso} in this case is that
  the evaluation map
  \[
    R[x] \to R, \quad g(x) \mapsto g(a)
  \]
  induces an isomorphism
  \[
    \frac{R[x]}{(x - a)} \cong R
  \]
  of abelian groups; the reader will verify (either by hand or by
  invoking \cref{coro:rings:first-iso}) that this is in fact an
  isomorphism of rings.
\end{exam}

\begin{exam}{}{exam:rings:higher-degree-quotient}
  It is fun to analyze higher-degree examples. For every monic
  $f(x) \in R[x]$ of degree $d$, \cref{prop:rings:poly-group-iso}
  gives a potentially \emph{different} ring (that is, $R[x]/(f(x))$)
  isomorphic to $R^{\oplus d}$ as a group; one can then use this
  isomorphism to define a new ring structure onto the group
  $R^{\oplus d}$.

  For $d = 1$ all these structures are isomorphic (as seen in the
  previous example); but interesting structures already arise in
  degree $2$. For a concrete example, apply this procedure with
  $f(x) = x^2 + 1$: \cref{prop:rings:poly-group-iso} gives an
  isomorphism of groups
  \[
    R \oplus R \cong \frac{R[x]}{(x^2 + 1)};
  \]
  what multiplication does this isomorphism induce on $R \oplus R$?
  Take two elements $(a_0, a_1)$, $(b_0, b_1)$ of $R \oplus R$. With
  the notation used in \cref{prop:rings:poly-group-iso}, we have
  \[
    (a_0, a_1) = \varphi(a_0 + a_1 x), \qquad (b_0, b_1) = \varphi(b_0
    + b_1 x).
  \]
  Now a bit of high-school algebra gives
  \begin{align*}
    (a_0 + a_1 x)(b_0 + b_1 x) &= a_0 b_0 + (a_0 b_1 + a_1 b_0)x + a_1 b_1 x^2 \\
                               &= (x^2 + 1)a_1 b_1 + ((a_0 b_0 - a_1 b_1) + (a_0 b_1 + a_1 b_0)x)
  \end{align*}
  which shows
  \[
    \varphi((a_0 + a_1 x)(b_0 + b_1 x)) = (a_0 b_0 - a_1 b_1,\; a_0
    b_1 + a_1 b_0).
  \]
  \emph{Therefore}, the multiplication induced on $R \oplus R$ by this
  procedure is defined by
  \[
    (a_0, a_1) \cdot (b_0, b_1) = (a_0 b_0 - a_1 b_1,\; a_0 b_1 + a_1
    b_0).
  \]
  This recipe may seem somewhat arbitrary, but note that upon taking
  $R = \R$, the ring of real numbers, and identifying pairs
  $(x, y) \in \R \oplus \R$ with complex numbers $x + iy$, the
  multiplication we obtained on $\R \oplus \R$ matches precisely the
  ordinary multiplication in $\C$. Therefore,
  \[
    \frac{\R[x]}{(x^2 + 1)} \cong \C
  \]
  as rings. In other words, this procedure constructs the ring $\C$
  `from scratch', starting from $\R[x]$.
\end{exam}

The point is that the polynomial equation $x^2 + 1 = 0$ has no
solutions in $\R$; the quotient $\R[x]/(x^2 + 1)$ produces a ring
containing a copy of $\R$ and in which the polynomial \emph{does} have
roots (that is, $\pm$ the class of $x$ in the quotient). The fact that
this turns out to be isomorphic to $\C$ may not be too surprising,
considering that $\C$ is precisely a ring containing a copy of $\R$
and in which $x^2 + 1$ does have roots (that is, $\pm i$).

Such constructions are the ``algebraist's way'' to solve equations. We
will come back to them in \S{}\ref{subsec:irred:adding-roots-algebraically-closed-fields}, and in a sense the whole of
\Cref{chap:fields} will be devoted to this topic.

\subsection{Prime and maximal ideals}
\label{subsec:rings:prime-and-maximal-ideals}
Other `qualities' of ideals are best expressed in terms of
quotients. I am still assuming that our rings are commutative---both
due to unforgivable laziness and because that is the only context in
which we will use these notions.

\begin{defn}{Prime and maximal ideals}{defn:rings:prime-maximal-ideal}
  Let $I \neq (1)$ be an ideal of a commutative ring $R$.
  \begin{itemize}
  \item $I$ is a \emph{prime ideal} if $R/I$ is an integral domain.
  \item $I$ is a \emph{maximal ideal} if $R/I$ is a field.
  \end{itemize}
\end{defn}

\begin{exam}{}{}
  For all $a \in R$, the ideal $(x - a)$ is prime in $R[x]$ if and
  only if $R$ is an integral domain; it is maximal if and only if $R$
  is a field. Indeed, $R[x]/(x - a) \cong R$, as we have seen in
  \Cref{exam:rings:remainder-degree-1}.

  The ideal $(2, x)$ is maximal in $\Z[x]$, since
  \[
    \frac{\Z[x]}{(2,x)} \stackrel{!}{\cong} \frac{\Z[x]/(x)}{(2)}
    \cong \frac{\Z}{(2)} = \Z/2\Z
  \]
  is a field (for the isomorphism $\stackrel{!}{\cong}$, cf.\
  \Cref{exam:rings:calculus-of-ideals}).
\end{exam}
Of course these notions may be translated into terms not involving
quotients at all, and it is largely a matter of \ae sthetic preference
whether prime and maximal ideals should be defined as in
\cref{defn:rings:prime-maximal-ideal} or by the following equivalent
conditions:
\begin{prop}{}{prop:rings:prime-maximal-equiv}
  Let $I \neq (1)$ be an ideal of a commutative ring $R$. Then
  \begin{itemize}
  \item $I$ is prime if and only if for all $a, b \in R$:
    $ab \in I \implies (a \in I$ or $b \in I)$;
  \item $I$ is maximal if and only if for all ideals $J$ of $R$:
    $I \subseteq J \implies (I = J$ or $J = R)$.
  \end{itemize}
\end{prop}
\begin{proof}
  The ring $R/I$ is an integral domain if and only if
  $\forall \bar{a}, \bar{b} \in R/I$:
  $\bar{a} \cdot \bar{b} = 0 \implies (\bar{a} = 0$ or $\bar{b} =
  0)$. This condition translates immediately to the given condition in
  $R$, with $\bar{a} = a + I$, $\bar{b} = b + I$, since the $0$ in
  $R/I$ is $I$.

  As for maximality, the given condition follows from the
  correspondence between ideals of $R/I$ and ideals of $R$ containing
  $I$ (\S{}\ref{subsec:rings:canonical-decomposition-and-consequences}) and the observation that a commutative ring is a field
  if and only if its only ideals are $(0)$ and $(1)$ (\Cref{exer:rings:division-ring-ideals}).
\end{proof}
From the formulation in terms of quotients, it is completely clear
that maximal $\implies$ prime; indeed, fields are integral
domains. This fact is of course easy to check in terms of the other
description, but the argument is a little more cumbersome
(\Cref{exer:rings:maximal-implies-prime-by-hand}). Prime ideals are not necessarily maximal, but note
the following:
\begin{prop}{}{prop:rings:finite-quotient-prime-max}
  Let $I$ be an ideal of a commutative ring $R$. If $R/I$ is finite,
  then $I$ is prime if and only if it is maximal.
\end{prop}
\begin{proof}
  This follows immediately from \cref{prop:rings:finite-field-domain}.
\end{proof}
For example, let $(n)$ be an ideal of $\Z$, with $n > 0$; then $(n)$
prime $\iff$ $(n)$ maximal $\iff$ $n$ is prime as an integer. Indeed,
for nonzero $n$ the ring $\Z/n\Z$ is finite, so
\cref{prop:rings:finite-quotient-prime-max} applies; cf.\
\Cref{exam:rings:znz-units-group}.  In general, the set of prime ideals of a commutative
ring $R$ is called\footnote{Believe it or not, the term is borrowed
  from functional analysis.} the \emph{spectrum} of $R$, denoted
$\operatorname{Spec} R$. We can `draw' $\operatorname{Spec} \Z$ as
follows:
% TODO: Spec Z diagram
Actually, attributing the fact that nonzero prime ideals of $\Z$ are
maximal to the finiteness of the quotients (as I have just done) is
slightly misleading; a better `explanation' is that $\Z$ is a PID, and
this phenomenon is common to all PIDs:
\begin{prop}{}{prop:rings:PID-prime-maximal}
  Let $R$ be a PID, and let $I$ be a nonzero ideal in $R$. Then $I$ is
  prime if and only if it is maximal.
\end{prop}
\begin{proof}
  Maximal ideals are prime in every ring, so we only need to verify
  that nonzero prime ideals are maximal in a PID; we will use the
  characterization of prime and maximal ideals obtained in
  \cref{prop:rings:prime-maximal-equiv}. Let $I = (a)$ be a prime
  ideal in $R$, with $a \neq 0$, and assume $I \subseteq J$ for an
  ideal of $R$. As $R$ is a PID, $J = (b)$ for some $b \in R$. Since
  $I = (a) \subseteq (b) = J$, we have that $a = bc$ for some
  $c \in R$. But then $b \in (a)$ or $c \in (a)$, since $I = (a)$ is
  prime.  If $b \in (a)$, then $(b) \subseteq (a)$; and $I = J$
  follows. If $c \in (a)$, then $c = da$ for some $d \in R$. But then
  $a = bc = bda$, from which $bd = 1$ since cancellation by the
  nonzero $a$ holds in $R$ (since $R$ is an integral domain). This
  implies that $b$ is a unit, and hence $J = (b) = R$.  That is, we
  have shown that if $I \subseteq J$, then either $I = J$ or $J = R$:
  thus $I$ is maximal, by \cref{prop:rings:prime-maximal-equiv}.
\end{proof}
\begin{exam}{}{exam:rings:speck-prime-ideals-maximal}
  Let $k$ be a field. Then nonzero prime ideals in $k[x]$ are maximal,
  since $k[x]$ is a PID (as the reader has hopefully checked by now;
  cf.\ \Cref{exer:rings:field-poly-ring-pid}). Therefore, a picture of
  $\operatorname{Spec} k[x]$ would look pretty much like the picture
  of $\operatorname{Spec} \Z$ shown above: with the maximal ideals at
  one level, all containing the prime (and nonmaximal) ideal $(0)$.
  This picture is particularly appealing for fields such as $k = \C$,
  which are \emph{algebraically closed}, that is, in which for every
  nonconstant $f(x) \in k[x]$ there exists $r \in k$ such that
  $f(r) = 0$. It would take us too far afield to discuss this notion
  at any length now; but the reader should be aware that $\C$ is
  algebraically closed. We will come back to this\footnote{A
    particularly pleasant proof may be given using elementary complex
    analysis, as a consequence of Liouville's theorem, or the maximum
    modulus principle; cf.\ \S{}\ref{subsec:irred:irreducibility-in-c-x-r-x-q-x}.}. Assuming this fact, it is
  easy to verify (\Cref{exer:rings:maximal-ideal-linear-poly}) that the maximal ideals in $\C[x]$
  are all and only the ideals $(x - z)$ where $z$ ranges over all
  complex numbers. That is, stretching our imagination a little, we
  could come up with the following picture for
  $\operatorname{Spec} \C[x]$:
  % TODO: Spec C[x] diagram
  There is a `complex line'\footnote{I know: it looks like a
    plane. But it is a line, as a complex entity.} worth of maximal
  ideals: for each $z \in \C$ we have the maximal $(x - z)$; the prime
  $(0)$ is contained in all the maximal ideals; and there are no other
  prime ideals.  The picture for $\operatorname{Spec} \C[x]$ may serve
  as justification for the fact that, in algebraic geometry, $\C[x]$
  is the ring corresponding to the `affine line' $\C$; it is the ring
  of an algebraic \emph{curve}. It turns out that the fact that there
  is exactly `one level' of maximal ideals over $(0)$ in $\C[x]$
  reflects precisely the fact that the corresponding geometric object
  has dimension $1$.

  In general, the \emph{(Krull) dimension} of a commutative ring $R$
  is the length of the longest chain of prime ideals in $R$. Thus,
  \cref{prop:rings:PID-prime-maximal} tells us that PIDs other than
  fields, such as $\Z$, have `dimension $1$'. In the lingo of
  algebraic geometry, they all correspond to curves.
\end{exam}

\begin{exam}{}{exam:rings:higher-dimension-rings}
  For examples of rings of higher dimension, consider
  $k[x_1, \ldots, x_n]$, where $k$ is a field. Note that there are
  chains of prime ideals of length $n$
  \[
    (0) \subsetneq (x_1) \subsetneq (x_1, x_2) \subsetneq \cdots
    \subsetneq (x_1, \ldots, x_n)
  \]
  in this ring. (Why are these ideals prime? Cf.\ \Cref{exer:rings:xk-prime-in-polyring}.) This
  says that $k[x_1, \ldots, x_n]$ has dimension $\geq n$. One can show
  that $n$ is in fact the longest length of a chain of prime ideals in
  $k[x_1, \ldots, x_n]$; that is, the Krull dimension of
  $k[x_1, \ldots, x_n]$ is precisely $n$. In algebraic geometry, the
  ring $\C[x_1, \ldots, x_n]$ corresponds to the $n$-dimensional
  complex space $\C^n$.

  Dealing precisely with these notions is not so easy, however. Even
  the seemingly simple statement that the maximal ideals of
  $\C[x_1, \ldots, x_n]$ are all and only the ideals
  $(x_1 - z_1, \ldots, x_n - z_n)$, for $(z_1, \ldots, z_n) \in \C^n$,
  requires a rather deep result, known as Hilbert's Nullstellensatz.
\end{exam}

We will come back to all of this and get a very small taste of
algebraic geometry in \S{}\ref{sec:fields:algebraic-closure-nullstellensatz-and-a-little-algebraic}, after we develop (much) more
machinery.

\subsection*{Exercises}
\begin{exer}{\exerused}{exer:rings:family-of-ideals-sum}
  Let $R$ be a ring, and let $\{I_\alpha\}_{\alpha \in A}$ be a family
  of ideals of $R$. We let
  \[
    \sum_{\alpha \in A} I_\alpha := \left\{ \sum_{\alpha \in A}
      r_\alpha \;\middle|\; r_\alpha \in I_\alpha \text{ and }
      r_\alpha = 0 \text{ for all but finitely many } \alpha \right\}.
  \]
  Prove that $\sum_\alpha I_\alpha$ is an ideal of $R$ and that it is
  the smallest ideal containing all of the ideals
  $I_\alpha$. [\S{}\ref{subsec:rings:basic-operations}]
\end{exer}
\begin{exer}{\exerused}{exer:rings:noetherian-homomorphic-image}
  Prove that the homomorphic image of a Noetherian ring is
  Noetherian. That is, prove that if $\varphi : R \to S$ is a
  surjective ring homomorphism and $R$ is Noetherian, then $S$ is
  Noetherian. [\S{}\ref{subsec:rings:submodule-generated-by-a-subset-noetherian-modules}]
\end{exer}
\begin{exer}{}{exer:rings:2x-not-principal}
  Prove that the ideal $(2, x)$ of $\Z[x]$ is not principal.
\end{exer}
\begin{exer}{\exerused}{exer:rings:field-poly-ring-pid}
  Prove that if $k$ is a field, then $k[x]$ is a PID. (Hint: Let
  $I \subseteq k[x]$ be any ideal. If $I = (0)$, then $I$ is
  principal. If $I \neq (0)$, let $f(x)$ be a monic polynomial in $I$
  of minimal degree. Use division with remainder to construct a proof
  that $I = (f(x))$, arguing as in the proof of \Cref{prop:groups:subgroups-of-z}.)
  [\S{}\ref{subsec:rings:basic-operations}, \S{}\ref{subsec:rings:prime-and-maximal-ideals}, \S{}\ref{subsec:irred:euclidean-domain-implies-pid}, \S{}\ref{subsec:irred:primitivity-and-content-gauss-s-lemma}, \S{}\ref{subsec:linalg:k-t-modules-and-the-rational-canonical-form}, \S{}\ref{subsec:fields:simple-extensions}]
\end{exer}
\begin{exer}{\exerused}{exer:rings:coprime-ideals-product-intersection}
  Let $I$, $J$ be ideals in a commutative ring $R$, such that
  $I + J = (1)$. Prove that $IJ = I \cap J$. [\S{}\ref{subsec:rings:basic-operations}, \S{}\ref{subsec:irred:chinese-remainder-theorem}]
\end{exer}
\begin{exer}{}{exer:rings:ij-reduced-implies-i-cap-j}
  Let $I$, $J$ be ideals in a commutative ring $R$. Assume that
  $R/(IJ)$ is reduced (that is, it has no nonzero nilpotent elements;
  cf.\ \Cref{exer:rings:nilradical-quotient-reduced}). Prove that $IJ = I \cap J$.
\end{exer}
\begin{exer}{\exerused}{exer:rings:field-poly-ideal-monic-generator}
  Let $R = k$ be a field. Prove that every nonzero (principal) ideal
  in $k[x]$ is generated by a unique monic polynomial. [\S{}\ref{subsec:rings:quotients-of-polynomial-rings},
  \S{}\ref{subsec:linalg:k-t-modules-and-the-rational-canonical-form}]
\end{exer}
\begin{exer}{\exerused}{exer:rings:monic-poly-not-zero-divisor}
  Let $R$ be a ring and $f(x) \in R[x]$ a monic polynomial. Prove that
  $f(x)$ is not a (left- or right-) zero-divisor. [\S{}\ref{subsec:rings:quotients-of-polynomial-rings}, \ref{exer:rings:zero-divisor-poly-generalization}]
\end{exer}
\begin{exer}{}{exer:rings:zero-divisor-poly-generalization}
  Generalize the result of \Cref{exer:rings:monic-poly-not-zero-divisor}, as follows. Let $R$ be a
  commutative ring, and let $f(x)$ be a zero-divisor in $R[x]$. Prove
  that $\exists\, b \in R$, $b \neq 0$, such that $f(x)b = 0$. (Hint:
  Let $f(x) = a_d x^d + \cdots + a_0$, and let
  $g(x) = b_e x^e + \cdots + b_0$ be a nonzero polynomial of minimal
  degree $e$ such that $f(x)g(x) = 0$. Deduce that $a_d g(x) = 0$, and
  then prove $a_{d-i} g(x) = 0$ for all $i$. What does this say about
  $b_e$?)
\end{exer}
\begin{exer}{\lnot}{exer:rings:z-sqrt-d-subring-of-c}
  Let $d$ be an integer that is not the square of an integer, and
  consider the subset of $\C$ defined by\footnote{Of course there are
    two `square roots of $d$'; but the definition of $\Q(\sqrt{d})$
    does not depend on which one is used.}
  \[
    \Q(\sqrt{d}) := \{a + b\sqrt{d} \mid a, b \in \Q\}.
  \]
  \begin{itemize}
  \item Prove that $\Q(\sqrt{d})$ is a subring of $\C$.
  \item Define a function $N : \Q(\sqrt{d}) \to \Q$ by
    $N(a + b\sqrt{d}) := a^2 - b^2 d$. Prove that $N(zw) = N(z)N(w)$
    and that $N(z) \neq 0$ if $z \in \Q(\sqrt{d})$, $z \neq 0$. The
    function $N$ is a `norm'; it is very useful in the study of
    $\Q(\sqrt{d})$ and of its subrings. (Cf.\ also \Cref{exer:rings:quaternion-norm}.)
  \item Prove that $\Q(\sqrt{d})$ is a field and in fact the smallest
    subfield of $\C$ containing both $\Q$ and $\sqrt{d}$. (Use $N$.)
  \item Prove that $\Q(\sqrt{d}) \cong \Q[t]/(t^2 - d)$. (Cf.\
    \Cref{exam:rings:higher-degree-quotient}.)
  \end{itemize} [\ref{exer:irred:z-sqrt-neg5-not-ufd}, \ref{exer:irred:z-1-plus-sqrt-neg19-over-2-not-euclidean}, \ref{exer:irred:z-sqrt2-euclidean-domain}, \ref{exer:fields:norm-is-multiplicative}]
\end{exer}
\begin{exer}{}{exer:rings:principal-ideal-with-a}
  Let $R$ be a commutative ring, $a \in R$, and
  $f_1(x), \ldots, f_r(x) \in R[x]$.
  \begin{itemize}
  \item Prove the equality of ideals
    $(f_1(x), \ldots, f_r(x), x - a) = (f_1(a), \ldots, f_r(a), x -
    a)$.
  \item Prove the useful substitution trick
    $\dfrac{R[x]}{(f_1(x), \ldots, f_r(x), x - a)} \cong
    \dfrac{R}{(f_1(a), \ldots, f_r(a))}$.
  \end{itemize}
  (Hint: \Cref{exer:rings:image-ideal-not-necessarily}.)
\end{exer}
\begin{exer}{\exerused}{exer:rings:ideal-generated-by-elements}
  Let $R$ be a commutative ring and $a_1, \ldots, a_n$ elements of
  $R$. Prove that
  $\dfrac{R[x_1, \ldots, x_n]}{(x_1 - a_1, \ldots, x_n - a_n)} \cong
  R$. [\S{}\ref{subsec:fields:the-nullstellensatz}]
\end{exer}
\begin{exer}{\exerused}{exer:rings:xk-prime-in-polyring}
  Let $R$ be an integral domain. For all $k = 1, \ldots, n$ prove that
  $(x_1, \ldots, x_k)$ is prime in $R[x_1, \ldots, x_n]$. [\S{}\ref{subsec:rings:prime-and-maximal-ideals}]
\end{exer}
\begin{exer}{\exerused}{exer:rings:maximal-implies-prime-by-hand}
  Prove `by hand' that maximal ideals are prime, without using
  quotient rings. [\S{}\ref{subsec:rings:prime-and-maximal-ideals}]
\end{exer}
\begin{exer}{}{exer:rings:preimage-prime-ideal}
  Let $\varphi : R \to S$ be a homomorphism of commutative rings, and
  let $I \subseteq S$ be an ideal. Prove that if $I$ is a prime ideal
  in $S$, then $\varphi^{-1}(I)$ is a prime ideal in $R$. Show that
  $\varphi^{-1}(I)$ is not necessarily maximal if $I$ is maximal.
\end{exer}
\begin{exer}{}{exer:rings:unique-zero-divisor-prime}
  Let $R$ be a commutative ring, and let $P$ be a prime ideal of
  $R$. Suppose $0$ is the only zero-divisor of $R$ contained in
  $P$. Prove that $R$ is an integral domain.
\end{exer}
\begin{exer}{\lnot}{exer:rings:continuous-functions-compact}
  (If you know a little topology\ldots) Let $K$ be a compact
  topological space, and let $R$ be the ring of continuous real-valued
  functions on $K$, with addition and multiplication defined
  pointwise.
  \begin{enumerate}[label=(\roman*)]
  \item For $p \in K$, let $M_p = \{f \in R \mid f(p) = 0\}$. Prove
    that $M_p$ is a maximal ideal in $R$.
  \item Prove that if $f_1, \ldots, f_r \in R$ have no common zeros,
    then $(f_1, \ldots, f_r) = (1)$. (Hint: Consider
    $f_1^2 + \cdots + f_r^2$.)
  \item Prove that every maximal ideal $M$ in $R$ is of the form $M_p$
    for some $p \in K$. (Hint: You will use the compactness of $K$ and
    (ii).)
  \end{enumerate}
  Conclude that $p \mapsto M_p$ defines a bijection from $K$ to the
  set of maximal ideals of $R$. (The set of maximal ideals of a
  commutative ring $R$ is called the \emph{maximal spectrum} of $R$;
  it is contained in the (prime) spectrum $\operatorname{Spec} R$
  defined in \S{}\ref{subsec:rings:prime-and-maximal-ideals}. Relating commutative rings and `geometric'
  entities such as topological spaces is the business of algebraic
  geometry.)

  The compactness hypothesis is necessary: cf.\
  \Cref{exer:irred:maximal-ideals-not-corresponding-to-points}. [\ref{exer:irred:maximal-ideals-not-corresponding-to-points}]
\end{exer}
\begin{exer}{}{exer:rings:nilradical-in-every-prime}
  Let $R$ be a commutative ring, and let $N$ be its nilradical
  (\Cref{exer:rings:nilpotents-form-ideal}). Prove that $N$ is contained in every prime ideal of
  $R$. (Later on the reader will check that the nilradical is in fact
  the intersection of all prime ideals of $R$: \Cref{exer:irred:nilradical-subset-intersection-of-primes}.)
\end{exer}
\begin{exer}{}{exer:rings:prime-ideal-family}
  Let $R$ be a commutative ring, let $P$ be a prime ideal in $R$, and
  let $I_j$ be ideals of $R$.
  \begin{enumerate}[label=(\roman*)]
  \item Assume that $I_1 \cdots I_r \subseteq P$; prove that
    $I_j \subseteq P$ for some $j$.
  \item By (i), if $P \supseteq \prod_{j=1}^r I_j$, then $P$ contains
    one of the ideals $I_j$. Prove or disprove: if
    $P \supseteq \prod_{j=1}^\infty I_j$, then $P$ contains one of the
    ideals $I_j$.
  \end{enumerate}
\end{exer}
\begin{exer}{}{exer:rings:maximal-iff-simple-quotient}
  Let $M$ be a two-sided ideal in a (not necessarily commutative) ring
  $R$. Prove that $M$ is maximal if and only if $R/M$ is a simple ring
  (cf.\ \Cref{exer:rings:division-ring-two-sided-counterexample}).
\end{exer}
\begin{exer}{\exerused}{exer:rings:maximal-ideal-linear-poly}
  Let $k$ be an algebraically closed field, and let $I \subseteq k[x]$
  be an ideal. Prove that $I$ is maximal if and only if $I = (x - c)$
  for some $c \in k$. [\S{}\ref{subsec:rings:prime-and-maximal-ideals}, \S{}\ref{subsec:irred:adding-roots-algebraically-closed-fields}, \S{}\ref{subsec:fields:algebraic-closure}, \S{}\ref{subsec:fields:the-nullstellensatz}]
\end{exer}
\begin{exer}{}{exer:rings:x2plus1-maximal-in-rx}
  Prove that $(x^2 + 1)$ is maximal in $\R[x]$.
\end{exer}
\begin{exer}{}{exer:rings:krull-dimension-zero}
  A ring $R$ has \emph{Krull dimension $0$} if every prime ideal in
  $R$ is maximal. Prove that fields and Boolean rings (\Cref{exer:rings:boolean-power-set})
  have Krull dimension $0$.
\end{exer}
\begin{exer}{}{exer:rings:z-x-krull-dimension-2}
  Prove that the ring $\Z[x]$ has Krull dimension $\geq 2$. (It is in
  fact exactly $2$; thus it corresponds to a surface from the point of
  view of algebraic geometry.)
\end{exer}

\section{Modules over a ring}
\label{sec:rings:modules-over-a-ring}

I have emphasized the parallels between the basic theory of groups and
the basic theory of rings; there are also important differences. In
$\categ{Grp}$, one takes the quotient of a group, by a (normal
sub)group, and the result is a group: throughout the process one never
leaves $\categ{Grp}$. The situation is in a sense even better in
$\categ{Ab}$, where the normality condition is automatic; one simply
takes the quotient of an abelian group by an abelian (sub)group,
obtaining an abelian group.  The situation in $\categ{Ring}$ is not
nearly as neat. One of the three characters in the story is an ideal:
which is not a ring according to the axioms listed in
\cref{defn:rings:ring}, in all but the most pathological cases. In
other words, the kernel of a homomorphism of rings is (usually) not a
ring. Also, cokernels do not behave as one would hope (cf.\
\Cref{exer:rings:z-in-q-cokernel}). Even if one relaxes the ring definition, giving up the
identity and making ideals subrings, there is no reasonably large
class of examples in which the `ideal' condition is automatically
satisfied by substructures. In short, $\categ{Ring}$ is not a
particularly pleasant category.  \emph{Modules} will fix all these
problems. If $R$ is a ring and $I \subseteq R$ is a two-sided ideal,
then all three structures $R$, $I$, and $R/I$ are modules over
$R$. The category of $R$-modules is the prime example of a
well-behaved category: in this category kernels and cokernels exist
and do precisely what they ought to. The category $\categ{Ab}$ is a
particular case of this construction, since it is the category of
modules over $\Z$ in disguise (\Cref{exam:rings:zmod-is-ab}). The category of modules
over a ring $R$ will share many of the excellent properties of
$\categ{Ab}$; and we will get a brand new, well-behaved category for
each ring $R$. These are all examples\footnote{In fact, it can be
  shown (`Freyd--Mitchell's embedding theorem') that every small
  abelian category is equivalent to a subcategory of the category of
  left-modules over a ring. So to some extent we can understand
  abelian categories by understanding categories of modules well
  enough. We will come back to this in \Cref{chap:homol}.} of the important
notion of \emph{abelian category}.

\subsection{Definition of (left-)R-module}
\label{subsec:rings:definition-of-left-r-module}
In short, $R$-modules are abelian groups endowed with an action of
$R$. To flesh out this idea, recall that `actions' in general denote
homomorphisms into some kind of endomorphism structure: for example,
we defined group actions in \S{}\ref{subsec:groups:actions} as group homomorphisms from a
fixed group to the groups of automorphisms of objects of a category.
We can give an analogous definition of the action of a ring on an
abelian group. Indeed, recall that if $M$ is an abelian group, then
$\End_{\categ{Ab}}(M) := \Hom_{\categ{Ab}}(M, M)$ is a ring in a
natural way (cf.\ \S{}\ref{subsec:rings:endab-g}). A \emph{left-action} of a ring $R$ on $M$
is then simply a homomorphism of rings
$\sigma : R \to \End_{\categ{Ab}}(M)$; we will say that $\sigma$ makes
$M$ into a \emph{left-$R$-module}.  Similarly to the situation with
group actions, it is convenient to spell out this definition.

\begin{prop}{}{prop:rings:module-axioms}
  The datum of a homomorphism $\sigma$ as above is precisely the same
  as the datum of a function $\rho : R \times M \to M$ satisfying the
  following requirements: $(\forall r, s \in R)$
  $(\forall m, n \in M)$
  \begin{itemize}
  \item $\rho(r, m + n) = \rho(r, m) + \rho(r, n)$;
  \item $\rho(r + s, m) = \rho(r, m) + \rho(s, m)$;
  \item $\rho(rs, m) = \rho(r, \rho(s, m))$;
  \item $\rho(1, m) = m$.
  \end{itemize}
\end{prop}

\begin{proof}
  The proof of this claim follows precisely the pattern of the
  discussion in the beginning of \S{}\ref{subsec:groups:actions-on-sets}, so it is left to the
  reader (\Cref{exer:rings:prove-prop5.1}). Of course the relation between $\rho$ and
  $\sigma$ is given by $\rho(r, m) = \sigma(r)(m)$.
\end{proof}

As always, carrying $\rho$ around is inconvenient, so it is common to
omit mention of it: $\rho(r, m)$ is often just denoted $rm$; this
makes the requirements listed above a little more
readable. Summarizing and adopting this shorthand,
\begin{defn}{}{defn:rings:left-R-module}
  A \emph{left-$R$-module structure} on an abelian group $M$ consists
  of a map $R \times M \to M$, $(r, m) \mapsto rm$, such that
  \begin{itemize}
  \item $r(m + n) = rm + rn$;
  \item $(r + s)m = rm + sm$;
  \item $(rs)m = r(sm)$;
  \item $1m = m$.
  \end{itemize}
\end{defn}
Right-$R$-modules are defined analogously. The reader can glance back
at \S{}\ref{sec:groups:group-actions}, especially \Cref{exer:groups:opposite-group}, for a reminder on right- vs.\
left-actions; the issues here are analogous. Thus, for example, a
right-$R$-module structure may be identified with a
left-$R^\circ$-module structure, where $R^\circ$ is the `opposite'
ring obtained by reversing the order of multiplication (Exercise
5.1). However, note that $R$ and $R^\circ$ have no good reasons to be
isomorphic in general (while every group is isomorphic to its
opposite).  These issues become immaterial if $R$ is commutative: then
the identity $R \to R^\circ$ is an isomorphism, and
left-modules/right-modules are identical concepts. The reader will not
miss much by adopting the blanket assumption that all rings mentioned
in this section are commutative. It is occasionally important to make
this hypothesis explicit (for example in dealing with algebras, cf.\
\Cref{exam:rings:ring-hom-defines-module}), but most of the material we are going to review works
verbatim for, say, left-modules over an arbitrary ring as for modules
over a commutative ring. I will write `module' for `left-module', for
convenience; it will be the reader's responsibility to take care of
appropriate changes, if necessary, to adapt the various concepts to
right-modules.

\subsection{The category R-Mod}
\label{subsec:rings:the-category-r-mod}
The reader should spend some time getting familiar with the notion of
module, by proving simple properties such as
\[
  (\forall m \in M) : 0 \cdot m = 0, \qquad (\forall m \in M) : (-1)
  \cdot m = -m,
\]
where $M$ is a module over some ring (\Cref{exer:rings:module-zero}). The following
fancier-sounding (but equally trivial) property is also useful in
developing a feel for modules:
\begin{prop}{}{prop:rings:Z-module}
  Every abelian group is a $\Z$-module, in exactly one way.
\end{prop}

\begin{proof}
  Let $G$ be an abelian group. A $\Z$-module structure on $G$ is a
  ring homomorphism $\Z \to \End_{\categ{Ab}}(G)$. Since $\Z$ is
  initial in $\categ{Ring}$ (\S{}\ref{subsec:rings:ring-homomorphisms}), there exists exactly one such
  homomorphism, proving the statement.
\end{proof}
Thus, `abelian group' and `$\Z$-module' are one and the same
notion. Quite concretely, the action of $n \in \Z$ on an element $a$
of an abelian group simply yields the ordinary `multiple' $na$; this
operation is trivially compatible with the operations in $\Z$.  A
homomorphism of $R$-modules is a homomorphism of (abelian) groups
which is compatible with the module structure. That is, if $M$, $N$
are $R$-modules and $\phi : M \to N$ is a function, then $\phi$ is a
homomorphism of $R$-modules if and only if
\begin{itemize}
\item
  $(\forall m_1 \in M)(\forall m_2 \in M) : \phi(m_1 + m_2) =
  \phi(m_1) + \phi(m_2)$;
\item $(\forall r \in R)(\forall m \in M) : \phi(rm) = r\phi(m)$.
\end{itemize}
It is hopefully clear that the composition of two $R$-module
homomorphisms is an $R$-module homomorphism and that the identity is
an $R$-module homomorphism: $R$-modules form a category which I will
denote\footnote{If $R$ is not commutative, we should agree on whether
  $R$-$\categ{Mod}$ denotes the category of left-modules or of
  right-modules. I will mean `left-modules'.} `$R$-$\categ{Mod}$'.
\begin{exam}{}{exam:rings:zmod-is-ab}
  The category $\Z$-$\categ{Mod}$ of $\Z$-modules is `the same as' the
  category $\categ{Ab}$: indeed, every abelian group is a $\Z$-module
  in exactly one way (\cref{prop:rings:Z-module}), and $\Z$-module
  homomorphisms are simply homomorphisms of abelian groups.
\end{exam}
\begin{exam}{}{exam:rings:k-vect}
  If $R = k$ is a field, $R$-modules are called $k$-vector spaces. I
  will call the category of vector spaces over a field $k$
  `$k$-$\categ{Vect}$'; this is just another name for
  $k$-$\categ{Mod}$. Morphisms in $k$-$\categ{Vect}$ are often
  called\footnote{This term is also used for homomorphisms of
    $R$-modules for more general rings $R$, but not as frequently.}
  linear maps. `Linear algebra' is the study of $k$-$\categ{Vect}$
  (extended to $R$-$\categ{Mod}$ when possible); Chapters~VI and~VIII
  will be devoted to this subject.
\end{exam}
\begin{exam}{}{exam:rings:ring-hom-defines-module}
  Any homomorphism of rings $\alpha : R \to S$ may be used to define
  an interesting $R$-module: define $\rho : R \times S \to S$ by
  $\rho(r, s) := \alpha(r)s$ for all $r \in R$ and $s \in S$. The
  operation on the right is simply multiplication in $S$, and the
  axioms of \cref{defn:rings:left-R-module} are immediate consequences
  of the ring axioms and of the fact that $\alpha$ is a
  homomorphism. For instance, taking $S = R$ and $\alpha = \id_R$
  makes $R$ a (left-)module over itself. It is common to write $rs$
  rather than $\alpha(r)s$.  The ring operation in $S$ and the
  $R$-module structure induced by $\alpha$ are linked even more
  tightly if we require $R$ to be commutative and $\alpha$ to map $R$
  to the center of $S$: that is, if we require $\alpha(r)$, $s$ to
  commute for every $r \in R$, $s \in S$. Indeed, in this case the
  left-module structure defined above and the right-module structure
  defined analogously would coincide; further, with these requirements
  the ring operation in $S$, $(s_1, s_2) \mapsto s_1 s_2$, is
  compatible with the $R$-module structure in the sense
  that\footnote{$(r_1 s_1)(r_2 s_2) = \alpha(r_1) s_1 \alpha(r_2) s_2
    = \alpha(r_1)\alpha(r_2) s_1 s_2 = (r_1 r_2)(s_1 s_2)$,
    $\forall r_1, r_2 \in R$, $\forall s_1, s_2 \in S$.} we can `move'
  the action of $R$ at will through products in $S$.
\end{exam}

Due to their importance, these examples deserve their own official
name:
\begin{defn}{}{defn:rings:R-algebra}
  Let $R$ be a commutative ring. An \emph{$R$-algebra} is a ring
  homomorphism $\alpha : R \to S$ such that $\alpha(R)$ is contained
  in the center of $S$.
\end{defn}
The usual abuse of language leads us to refer to an $R$-algebra by the
name of the target $S$ of the homomorphism. Thus, an $R$-algebra `is'
an $R$-module $S$ with a compatible ring structure, or, if you prefer,
a ring $S$ with a compatible $R$-module structure. An $R$-algebra $S$
is a \emph{division algebra} if $S$ is a division ring.

There is an evident notion of `$R$-algebra homomorphism' (preserving
both the ring and module structure), and we thus get a category
$R$-$\categ{Alg}$. The situation simplifies substantially if $S$ is
itself commutative, in which case the condition on the center is
unnecessary. `Commutative $R$-algebras' form a category, which the
attentive reader will recognize as a coslice category (\Cref{exam:prelims:coslice-category})
in the category of commutative rings.

Also, note that $\Z$-$\categ{Alg}$ is just another name for
$\categ{Ring}$ (why?).  The polynomial rings $R[x_1, \ldots, x_n]$, as
well as all their quotients, are commutative $R$-algebras. This is a
particularly important class of examples; for $R = k$ an algebraically
closed field, these are the rings used in `classical' affine algebraic
geometry (cf.\ \S{}\ref{subsec:fields:a-little-affine-algebraic-geometry}).  The trivial group $0$ has a unique
module structure over any ring $R$ and is a zero-object in
$R$-$\categ{Mod}$, that is, it is both initial and final. As in the
other main categories we have encountered, a bijective homomorphism of
$R$-modules is automatically an isomorphism in $R$-$\categ{Mod}$
(\Cref{exer:rings:bijective-iso}). In these and many other respects, the category
$R$-$\categ{Mod}$ (for any commutative ring $R$) and $\categ{Ab}$ are
similar.  If $R$ is commutative, the similarity goes further: just as
in the category $\categ{Ab}$, each set
$\Hom_{R\text{-}\categ{Mod}}(M, N)$ may itself be seen as an object of
the category (cf.\ \S{}\ref{subsec:groups:homomorphisms-of-abelian-groups})\footnote{Notational convention: One
  often writes $\Hom_R(M, N)$ for $\Hom_{R\text{-}\categ{Mod}}(M,
  N)$. I will not adopt this convention here but I will use it freely
  in later chapters.}.  Indeed, let $M$ and $N$ be $R$-modules. Since
homomorphisms of $R$-modules are in particular homomorphisms of
abelian groups,
$\Hom_{R\text{-}\categ{Mod}}(M, N) \subseteq \Hom_{\categ{Ab}}(M, N)$
as sets (up to natural identifications). The operation making
$\Hom_{\categ{Ab}}(M, N)$ into an abelian group, as in \S{}\ref{subsec:groups:homomorphisms-of-abelian-groups},
clearly preserves $\Hom_{R\text{-}\categ{Mod}}(M, N)$; it follows that
the latter is an abelian group. For $r \in R$ and
$\phi \in \Hom_{R\text{-}\categ{Mod}}(M, N)$, the prescription
$(\forall m \in M) : (r\phi)(m) := r\phi(m)$ defines a
function\footnote{Parse the notation carefully: $r\phi$ is the name of
  a function on the left, while $r\phi(m)$ on the right is the action
  of the element $r \in R$ on the element $\phi(m) \in N$, defined by
  the $R$-module structure on $N$.} $r\phi : M \to N$. This function
is an $R$-module homomorphism if $R$ is commutative, because
$(\forall a \in R)$, $(\forall m \in M)$:
$(r\phi)(am) = r\phi(am) = (ra)\phi(m) = (ar)\phi(m) =
a(r\phi(m))$. Thus, we have a natural action of $R$ on the abelian
group $\Hom_{R\text{-}\categ{Mod}}(M, N)$, and it is immediate to
verify that this makes $\Hom_{R\text{-}\categ{Mod}}(M, N)$ into an
$R$-module. Watch out: if $R$ is not commutative, then in general
$\Hom_{R\text{-}\categ{Mod}}(M, N)$ is `just' an abelian group. More
structure is available if $M$, $N$ are bimodules; cf.\ \S{}\ref{subsec:linrep:bimodules-adjunction-again}.

\subsection{Submodules and quotients}
\label{subsec:rings:submodules-and-quotients}
Since $R$-modules are an `enriched' version of abelian groups, we can
progress through the usual constructions very quickly, by just
pointing out that the analogous constructions in $\categ{Ab}$ are
preserved by the $R$-module structure.

A \emph{submodule} $N$ of an $R$-module $M$ is a subgroup preserved by
the action of $R$. That is, for all $r \in R$ and $n \in N$, the
element $rn$ (defined by the $R$-module structure of $M$) is in fact
in $N$. Put otherwise, and perhaps more transparently, $N$ is itself
an $R$-module, and the inclusion $N \subseteq M$ is an $R$-module
homomorphism.
\begin{exam}{}{}
  We can view $R$ itself as a (left-)$R$-module (cf.\ \Cref{exam:rings:ring-hom-defines-module});
  the submodules of $R$ are then precisely the (left-)ideals of $R$.
\end{exam}
\begin{exam}{}{}
  Both the kernel and the image of a homomorphism $\phi : M \to M'$ of
  $R$-modules are submodules (of $M$, $M'$, respectively).
\end{exam}
\begin{exam}{}{}
  If $r$ is in the center of $R$ and $M$ is an $R$-module, then
  $rM = \{rm \mid m \in M\}$ is a submodule of $M$. If $I$ is any
  (left-)ideal of $R$, then
  $IM = \{\sum_i r_i m_i \mid r_i \in I, m_i \in M\}$ is a submodule
  of $M$.
\end{exam}
If $N$ is a submodule of $M$, then it is in particular a (normal)
subgroup of the abelian group $(M, +)$; thus we may define the
quotient $M/N$ as an abelian group. Of course it would be desirable to
see this as a module, and as usual there is only one reasonable way to
do so: we will want the canonical projection $\pi : M \to M/N$ to be
an $R$-module homomorphism, and this forces
$r(m + N) = r\pi(m) = \pi(rm) = rm + N$ for all $m \in M$. That is, we
are led to define the action of $R$ on $M/N$ by $r(m + N) := rm + N$.
\begin{prop}{}{prop:rings:quotient-module}
  For all submodules $N$, this prescription does define a structure of
  $R$-module on $M/N$.
\end{prop}

\begin{proof}
  The proof of this claim is immediate and is left to the reader.
\end{proof}

The $R$-module $M/N$ is (of course) called the \emph{quotient} of $M$
by $N$.
\begin{exam}{}{}
  If $R$ is a ring and $I$ is a two-sided ideal of $R$, then all three
  of $I$, $R$, and the quotient ring $R/I$ are $R$-modules: $I$ is a
  submodule of $R$, and the rings $R$ and $R/I$ are in fact
  $R$-algebras if $R$ is commutative (cf.\ \Cref{exam:rings:ring-hom-defines-module}).
\end{exam}
\begin{exam}{}{}
  If $R$ is not commutative and $I$ is just a (say) left-ideal, then
  the quotient $R/I$ is not defined as a ring, but it is defined as a
  left-module (the quotient of the module $R$ by the submodule
  $I$). The action of $R$ on $R/I$ is given by left-multiplication:
  $r(a + I) = ra + I$.
\end{exam}
The reader should now expect a universal property for quotients, and
here it is:

\begin{thm}{}{thm:rings:module-quotient-universal}
  Let $N$ be a submodule of an $R$-module $M$. Then for every
  homomorphism of $R$-modules $\phi : M \to P$ such that
  $N \subseteq \ker\phi$, there exists a unique homomorphism of
  $R$-modules $\tilde{\phi} : M/N \to P$ so that the diagram
  \[\begin{tikzcd}
      M \arrow[r, "\phi"] \arrow[d, "\pi"'] & P \\
      M/N \arrow[ur, "\exists!\tilde{\phi}"']
    \end{tikzcd}\] commutes.
\end{thm}
\begin{proof}
  As in previous appearances of such statements, this is an immediate
  consequence of the set-theoretic version (\S{}\ref{subsec:prelims:quotients}) and of easy
  notation matching and compatibility checks. For an even faster
  proof, one can just apply \Cref{thm:groups:quotient-universal-property} and verify that
  $\tilde{\phi}$ is an $R$-module homomorphism.
\end{proof}

Since every submodule $N$ is then the kernel of the canonical
projection $M \to M/N$, our recurring slogan becomes, in the context
of $R$-$\categ{Mod}$:
\[
  \text{kernel} \iff \text{submodule}:
\]
unlike as in $\categ{Grp}$ or $\categ{Ring}$, being a kernel poses no
restriction on the relevant substructures. Put otherwise, `every
monomorphism in $R$-$\categ{Mod}$ is a kernel'; this is one of the
distinguishing features of an abelian category.

\subsection{Canonical decomposition and isomorphism theorems}
\label{subsec:rings:canonical-decomposition-and-isomorphism-theorems}
The discussion now proceeds along the same lines as for (abelian)
groups; the statements of the key facts and a few comments should
suffice, as the proofs are nothing but a rehashing of the proofs of
analogous statements we have encountered previously. Of course the
reader should take the following statements as assignments and provide
all the needed details.

In the context of $R$-modules, the canonical decomposition takes the
following form:
\begin{thm}{}{thm:rings:canonical-decomposition}
  Every $R$-module homomorphism $\phi : M \to M'$ may be decomposed as
  follows:
  \[\begin{tikzcd}
      M \arrow[r, two heads] \arrow[rr, bend left, "\phi"] &
      M/\ker\phi \arrow[r, "\tilde{\phi}", "\sim"'] &
      \image\phi \arrow[r, hook] & M'
    \end{tikzcd}\] where the isomorphism $\tilde{\phi}$ in the middle
  is the homomorphism induced by $\phi$ (as in
  \cref{thm:rings:module-quotient-universal}).
\end{thm}

The `first isomorphism theorem' is the following consequence:
\begin{coro}{}{coro:rings:module-first-iso}
  Suppose $\phi : M \to M'$ is a surjective $R$-module
  homomorphism. Then $M' \cong M / \ker\phi$.
\end{coro}
If $M$ is an $R$-module and $N$ is a submodule of $M$, then there is a
bijection (cf.\ \S{}\ref{subsec:groups:subgroups-of-quotients})
$u : \{\text{submodules } P \text{ of } M \text{ containing } N\} \to
\{\text{submodules of } M/N\}$ preserving inclusions, and the `third
isomorphism theorem' holds:

\begin{prop}{}{prop:rings:third-iso-modules}
  Let $N$ be a submodule of an $R$-module $M$, and let $P$ be a
  submodule of $M$ containing $N$. Then $P/N$ is a submodule of $M/N$,
  and $(M/N) / (P/N) \cong M/P$.
\end{prop}
We also have a version of the `second isomorphism theorem' (cf.\
\Cref{prop:groups:second-isomorphism-theorem}), with simplifications due to the fact that
normality is not an issue in the theory of modules:

\begin{prop}{}{prop:rings:second-iso}
  Let $N$, $P$ be submodules of an $R$-module $M$. Then
  \begin{itemize}
  \item $N + P$ is a submodule of $M$;
  \item $N \cap P$ is a submodule of $P$, and
    $(N + P)/N \cong P/(N \cap P)$.
  \end{itemize}
\end{prop}

More generally, it is hopefully clear that the sum
$\sum_\alpha N_\alpha$ and the intersection $\bigcap_\alpha N_\alpha$
of any family $\{N_\alpha\}_\alpha$ of submodules of an $R$-module $M$
(which are defined as subgroups of the abelian group $M$; cf.\ for
example \Cref{lem:groups:intersection-of-subgroups}) are submodules of $M$.

\subsection*{Exercises}

\begin{exer}{\exerused}{exer:rings:opposite-ring}
  Let $R$ be a ring. The opposite ring $R^\circ$ is obtained from $R$
  by reversing the multiplication: that is, the product $a \bullet b$
  in $R^\circ$ is defined to be $ba \in R$. Prove that the identity
  map $R \to R^\circ$ is an isomorphism if and only if $R$ is
  commutative. Prove that $M_n(R)$ is isomorphic to its opposite (not
  via the identity map!). Explain how to turn right-$R$-modules into
  left-$R$-modules and conversely, if $R \cong R^\circ$. [\S{}\ref{subsec:rings:definition-of-left-r-module},
  \ref{exer:linrep:dual-has-right-module-structure-r-iso-rop}]
\end{exer}
\begin{exer}{\exerused}{exer:rings:prove-prop5.1}
  Prove \cref{prop:rings:module-axioms}. [\S{}\ref{subsec:rings:definition-of-left-r-module}]
\end{exer}
\begin{exer}{\exerused}{exer:rings:module-zero}
  Let $M$ be a module over a ring $R$. Prove that $0 \cdot m = 0$ and
  that $(-1) \cdot m = -m$, for all $m \in M$. [\S{}\ref{subsec:rings:the-category-r-mod}]
\end{exer}
\begin{exer}{\lnot}{exer:rings:schur}
  Let $R$ be a ring. A nonzero $R$-module $M$ is \emph{simple} (or
  \emph{irreducible}) if its only submodules are $\{0\}$ and $M$. Let
  $M$, $N$ be simple modules, and let $\phi : M \to N$ be a
  homomorphism of $R$-modules. Prove that either $\phi = 0$ or $\phi$
  is an isomorphism. (This rather innocent statement is known as
  \emph{Schur's lemma}.) [\ref{exer:rings:simple-endomorphism-division}, \ref{exer:rings:cyclic-module}, \ref{exer:linalg:jordan-holder-for-modules}]
\end{exer}
\begin{exer}{}{exer:rings:hom-R-M}
  Let $R$ be a commutative ring, viewed as an $R$-module over itself,
  and let $M$ be an $R$-module. Prove that
  $\Hom_{R\text{-}\categ{Mod}}(R, M) \cong M$ as $R$-modules.
\end{exer}
\begin{exer}{}{exer:rings:Q-vector-unique}
  Let $G$ be an abelian group. Prove that if $G$ has a structure of
  $\Q$-vector space, then it has only one such structure. (Hint: First
  prove that every element of $G$ has necessarily infinite
  order. Alternative hint: The unique ring homomorphism $\Z \to \Q$ is
  an epimorphism.)
\end{exer}
\begin{exer}{}{exer:rings:field-extension}
  Let $K$ be a field, and let $k \subseteq K$ be a subfield of
  $K$. Show that $K$ is a vector space over $k$ (and in fact a
  $k$-algebra) in a natural way. In this situation, we say that $K$ is
  an \emph{extension} of $k$.
\end{exer}
\begin{exer}{}{exer:rings:R-alg-initial}
  What is the initial object of the category $R$-$\categ{Alg}$?
\end{exer}
\begin{exer}{\lnot}{exer:rings:endomorphism-algebra}
  Let $R$ be a commutative ring, and let $M$ be an $R$-module. Prove
  that the operation of composition on the $R$-module
  $\End_{R\text{-}\categ{Mod}}(M)$ makes the latter an $R$-algebra in
  a natural way. Prove that $M_n(R)$ (cf.\ \Cref{exer:rings:matrix-ring}) is an
  $R$-algebra, in a natural way. [\ref{exer:linalg:ibn-fails-for-endomorphism-ring}, \ref{exer:linalg:mn-r-iso-hom-as-r-algebras}]
\end{exer}
\begin{exer}{}{exer:rings:simple-endomorphism-division}
  Let $R$ be a commutative ring, and let $M$ be a simple $R$-module
  (cf.\ \Cref{exer:rings:schur}). Prove that $\End_{R\text{-}\categ{Mod}}(M)$ is
  a division $R$-algebra.
\end{exer}
\begin{exer}{\exerused}{exer:rings:Rx-module-structures}
  Let $R$ be a commutative ring, and let $M$ be an $R$-module. Prove
  that there is a natural bijection between the set of $R[x]$-module
  structures on $M$ and $\End_{R\text{-}\categ{Mod}}(M)$. [\S{}\ref{subsec:linalg:linear-transformations-of-free-modules-actions-of-polynomial-rings}]
\end{exer}
\begin{exer}{\exerused}{exer:rings:bijective-iso}
  Let $R$ be a ring. Let $M$, $N$ be $R$-modules, and let
  $\phi : M \to N$ be a homomorphism of $R$-modules. Assume $\phi$ is
  a bijection, so that it has an inverse $\phi^{-1}$ as a
  set-function. Prove that $\phi^{-1}$ is a homomorphism of
  $R$-modules. Conclude that a bijective $R$-module homomorphism is an
  isomorphism of $R$-modules. [\S{}\ref{subsec:rings:the-category-r-mod}, \S{}\ref{subsec:linalg:matrices}, \S{}\ref{subsec:homol:abelian-categories}]
\end{exer}
\begin{exer}{}{exer:rings:principal-ideal-iso-R}
  Let $R$ be an integral domain, and let $I$ be a nonzero principal
  ideal of $R$. Prove that $I$ is isomorphic to $R$ as an $R$-module.
\end{exer}
\begin{exer}{\exerused}{exer:rings:prove-second-iso}
  Prove \cref{prop:rings:second-iso}. [\S{}\ref{subsec:rings:canonical-decomposition-and-isomorphism-theorems}]
\end{exer}
\begin{exer}{}{exer:rings:ideal-quotient-iso}
  Let $R$ be a commutative ring, and let $I$, $J$ be ideals of
  $R$. Prove that $I \cdot (R/J) \cong (I + J)/J$ as $R$-modules.
\end{exer}
\begin{exer}{\lnot}{exer:rings:nakayama-nilpotent}
  Let $R$ be a commutative ring, $M$ an $R$-module, and let $a \in R$
  be a nilpotent element, determining a submodule $aM$ of $M$. Prove
  that $M = 0 \iff aM = M$. (This is a particular case of Nakayama's
  lemma, \Cref{exer:linalg:nakayamas-lemma}.) [\ref{exer:linalg:nakayamas-lemma}]
\end{exer}
\begin{exer}{\exerused}{exer:rings:rees-algebra}
  Let $R$ be a commutative ring, and let $I$ be an ideal of
  $R$. Noting that $I^j \cdot I^k \subseteq I^{j+k}$, define a ring
  structure on the direct sum
  \[
    \operatorname{Rees}_R(I) := \bigoplus_{j \geq 0} I^j = R \oplus I
    \oplus I^2 \oplus I^3 \oplus \cdots.
  \]
  The homomorphism sending $R$ identically to the first term in this
  direct sum makes $\operatorname{Rees}_R(I)$ into an $R$-algebra,
  called the \emph{Rees algebra} of $I$. Prove that if $a \in R$ is a
  non-zero-divisor, then the Rees algebra of $(a)$ is isomorphic to
  the polynomial ring $R[x]$ (as an $R$-algebra). [\ref{exer:rings:rees-iso}]
\end{exer}
\begin{exer}{}{exer:rings:rees-iso}
  With notation as in \Cref{exer:rings:rees-algebra}, let $a \in R$ be a
  non-zero-divisor, let $I$ be any ideal of $R$, and let $J$ be the
  ideal $aI$. Prove that
  $\operatorname{Rees}_R(J) \cong \operatorname{Rees}_R(I)$.
\end{exer}

\section{Products, coproducts, etc., in R-Mod}
\label{sec:rings:products-coproducts-etc-in-r-mod}

I have stated several times that categories such as $R$-$\categ{Mod}$
are `well-behaved'. We will explore in this section the sense in which
this can be formalized at this stage. The bottom line is that these
categories enjoy the same nice properties that we have noted along the
way for the category $\categ{Ab}$.

I will also include here some general considerations on finitely
generated modules and algebras.

As in the previous section, I will write `module' for `left-module';
the reader should make appropriate adaptations to the case of
right-modules. Little will be lost by assuming that all rings
appearing here are commutative (thereby removing the distinction
between left- and right-modules).

\subsection{Products and coproducts}
\label{subsec:rings:products-and-coproducts}
As in $\categ{Ab}$, products and coproducts exist, and finite products
and coproducts coincide, in $R$-$\categ{Mod}$. Indeed, recall the
construction of the direct sum of two abelian groups (\S{}\ref{subsec:groups:abelian-groups}): if
$M$ and $N$ are abelian groups, then $M \oplus N$ denotes their
product, with componentwise operation. If $M$ and $N$ are $R$-modules,
we can give an $R$-module structure to $M \oplus N$ by prescribing
$\forall r \in R$: $r(m, n) := (rm, rn)$. This defines the
\emph{direct sum} of $M$, $N$, as an $R$-module. Note that
$M \oplus N$ comes together with several homomorphisms of $R$-modules:
\[
  \pi_M : M \oplus N \to M, \qquad \pi_N : M \oplus N \to N
\]
sending $(m, n)$ to $m$, $n$, respectively, and
\[
  i_M : M \to M \oplus N, \qquad i_N : N \to M \oplus N
\]
sending $m$ to $(m, 0)$ and $n$ to $(0, n)$.
\begin{prop}{}{prop:rings:direct-sum-product-coproduct}
  The direct sum $M \oplus N$ satisfies the universal properties of
  both the product and the coproduct of $M$ and $N$.
\end{prop}

\begin{proof}
  \emph{Product:} Let $P$ be an $R$-module, and let
  $\phi_M : P \to M$, $\phi_N : P \to N$ be two $R$-module
  homomorphisms. The definition of an $R$-module homomorphism
  $\phi_M \times \phi_N : P \to M \oplus N$ is forced by the needed
  commutativity of the diagram
  \[\begin{tikzcd}
      & P \arrow[dl, "\phi_M"'] \arrow[d, dashed, "\phi_M \times \phi_N"] \arrow[dr, "\phi_N"] & \\
      M & M \oplus N \arrow[l, "\pi_M"] \arrow[r, "\pi_N"'] & N
    \end{tikzcd}\] That is,
  $(\forall p \in P) : (\phi_M \times \phi_N)(p) := (\phi_M(p),
  \phi_N(p))$. This is an $R$-module homomorphism, and it is the
  unique one making the diagram commute; therefore $M \oplus N$ works
  as a product of $M$ and $N$.

  \emph{Coproduct:} View the preceding argument through a mirror! Let
  $P$ be an $R$-module, and let $\psi_M : M \to P$, $\psi_N : N \to P$
  be two $R$-module homomorphisms. The definition of an $R$-module
  homomorphism $\psi_M \oplus \psi_N : M \oplus N \to P$ is forced by
  the needed commutativity of the diagram
  \[\begin{tikzcd}
      M \arrow[r, "i_M"] \arrow[dr, "\psi_M"'] & M \oplus N \arrow[d, dashed, "\psi_M \oplus \psi_N"] & N \arrow[l, "i_N"'] \arrow[dl, "\psi_N"] \\
      & P &
    \end{tikzcd}\] That is\footnote{This should look familiar; cf.\
    \Cref{exer:groups:ab-coproduct}.}, $(\forall m \in M)(\forall n \in N)$:
  \begin{align*}
    (\psi_M \oplus \psi_N)(m, n) &= (\psi_M \oplus \psi_N)(m, 0) + (\psi_M \oplus \psi_N)(0, n) \\
                                 &= (\psi_M \oplus \psi_N) \circ i_M(m) + (\psi_M \oplus \psi_N) \circ i_N(n) = \psi_M(m) + \psi_N(n).
  \end{align*}
  This is an $R$-module homomorphism: it is a homomorphism of abelian
  groups by virtue of the commutativity of addition in $P$, and it
  clearly preserves the action of $R$.
\end{proof}
It may seem like a good idea to write $M \times N$ rather than
$M \oplus N$ when viewing the latter as the product of $M$ and $N$;
but in due time (\S{}\ref{sec:linrep:tensor-products-and-the-tor-functors}) we will encounter $M \times N$ again, in
the context of `bilinear maps', and in that context $M \times N$ is
not viewed as an $R$-module.

The reader should also work out the fibered versions of these
constructions; cf.\ Exercises~6.10 and~6.11.

The fact that finite products and coproducts agree in
$R$-$\categ{Mod}$ does not extend to the infinite case (\Cref{exer:rings:product-coproduct-infinite}).

\subsection{Kernels and cokernels}
\label{subsec:rings:kernels-and-cokernels}
The facts that $R$-$\categ{Mod}$ has a zero-object (the $0$-module),
its Hom sets are abelian groups (\S{}\ref{subsec:rings:the-category-r-mod}), and it has (finite)
products and coproducts make $R$-$\categ{Mod}$ an \emph{additive
  category}. The fact that $R$-$\categ{Mod}$ has well-behaved kernels
and cokernels, which we review next, upgrades it to the status of
\emph{abelian category}. (We will come back to these general
definitions in \S{}\ref{sec:homol:un-necessary-categorical-preliminaries}.)

In general, monomorphisms and epimorphisms do not automatically
satisfy good properties, even when objects of a given category are
realized by adding structure to sets. For example, we have seen that
the precise relationship between `surjective morphism' and
`epimorphism' may be rather subtle: epimorphisms are surjective in
$\categ{Grp}$, but for complicated reasons (\S{}\ref{subsec:groups:epimorphisms-and-cokernels}); and there are
epimorphisms that are not surjective in $\categ{Ring}$ (\S{}\ref{subsec:rings:monomorphisms-and-epimorphisms}).  The
situation in $R$-$\categ{Mod}$ is as simple as it can be. Recall that
we have identified universal properties for kernels and cokernels
(cf.\ \S{}\ref{subsec:groups:epimorphisms-and-cokernels}); in the category $R$-$\categ{Mod}$ these would go as
follows: if $\phi : M \to N$ is a homomorphism of $R$-modules, then
$\ker\phi$ is final with respect to the property of factoring
$R$-module homomorphisms $\alpha : P \to M$ such that
$\phi \circ \alpha = 0$:
\[\begin{tikzcd}
    P \arrow[r, dashed, "\exists!\alpha"] \arrow[rr, bend left,
    "\alpha"] \arrow[rrr, bend left=40, "0"] & \ker\phi \arrow[r] & M
    \arrow[r, "\phi"] & N
  \end{tikzcd}\] while $\coker\phi$ is initial with
respect to the property of factoring $R$-module homomorphisms
$\beta : N \to P$ such that $\beta \circ \phi = 0$:
\[\begin{tikzcd}
    M \arrow[r, "\phi"] \arrow[rrr, bend right=40, "0"'] & N \arrow[r,
    "\pi"] \arrow[rr, bend right, "\beta"'] & \coker\phi
    \arrow[r, dashed, "\exists!\beta"'] & P
  \end{tikzcd}\]
\begin{prop}{}{prop:rings:kernels-cokernels}
  The following hold in $R$-$\categ{Mod}$:
  \begin{itemize}
  \item kernels and cokernels exist;
  \item $\phi$ is a monomorphism $\iff$ $\ker\phi$ is trivial $\iff$
    $\phi$ is injective as a set-function;
  \item $\phi$ is an epimorphism $\iff$ $\coker\phi$ is
    trivial $\iff$ $\phi$ is surjective as a set-function.
  \end{itemize}
  Further, every monomorphism identifies its source with the kernel of
  some morphism, and every epimorphism identifies its target with the
  cokernel of some morphism.
\end{prop}
\begin{proof}
  This proposition of course simply generalizes to $R$-$\categ{Mod}$
  facts we know already from our study of $\categ{Ab}$, and a quick
  review should suffice for the careful reader. \emph{Kernels exist:}
  indeed, the `standard' definition of kernel satisfies the universal
  properties spelled out above (same argument as in
  \Cref{prop:groups:kernel-universal-property}). \emph{Cokernels exist:} indeed, let
  $\coker\phi = N / \image\phi$; if
  $\beta : N \to P$ is such that $\beta \circ \phi = 0$, then
  $\image\phi \subseteq \ker\beta$; so $\beta$ must factor
  uniquely through $N/\image\phi$ by the universal property
  of quotients, \cref{thm:rings:module-quotient-universal}. That is,
  $N/\image\phi$ does satisfy the universal property for
  cokernels.

  The proofs of all the implications in the second and third points
  follow familiar patterns from (for example) \Cref{prop:groups:monomorphism-characterization} and
  \Cref{prop:groups:epimorphism-characterization}. The last sentence of
  \cref{prop:rings:kernels-cokernels} simply reiterates the slogan
  `submodule $\iff$ kernel' and its mirror statement (which is just as
  true). Further details are left to the reader.
\end{proof}
By the way, whatever happened to the conditions characterizing
monomorphisms and epimorphisms in $\categ{Set}$ (\Cref{prop:inv-inj-sur})?
In $\categ{Set}$, a function with non-empty source is a monomorphism
if and only if it has a left-inverse, and it is an epimorphism if and
only if it has a right-inverse. We have learned not to expect any of
this to happen in more general categories. Modules are not an
exception: the function $\Z \to \Z$ defined by `multiplication by 2'
is a monomorphism without a left-inverse, and the projection
$\Z \to \Z/2\Z$ is an epimorphism without a right-inverse. We will
come back to this point in \S{}\ref{subsec:rings:split-exact-sequences}.

\subsection{Free modules and free algebras}
\label{subsec:rings:free-modules-and-free-algebras}
The universal property of free $R$-modules is modeled after the
properties defining the other free objects we have encountered: the
goal is to define the $R$-module containing a given set $A$ `in the
most efficient way'. Again, the situation will match the case of
abelian groups closely, so the reader may want to refer back to
\S{}\ref{subsec:groups:free-abelian-groups}.

The universal property formalizing the heuristic requirement goes as
follows: given a set $A$, we are seeking an $R$-module $F_R(A)$,
called a \emph{free $R$-module on the set $A$}, together with a
set-function $j : A \to F_R(A)$, such that for all $R$-modules $M$ and
set-functions $f : A \to M$ there exists a unique $R$-module
homomorphism $\phi : F_R(A) \to M$ such that the diagram
\[\begin{tikzcd}
    F_R(A) \arrow[r, "\phi"] & M \\
    A \arrow[u, "j"] \arrow[ur, "f"']
  \end{tikzcd}\] commutes. Abstract nonsense guarantees that such an
$R$-module is unique up to isomorphism if it exists at all
(\Cref{prop:prelims:initial-final-unique}) and that the function $j : A \to F_R(A)$ is
necessarily injective (cf.\ \Cref{exer:groups:j-injective}). The question is, does it
exist?  The answer will not appear to be very exciting, since it
generalizes directly the case of abelian groups (that is,
$\Z$-modules). Given any set $A$, define the (possibly infinite)
direct sum $N^{\oplus A}$ of an $R$-module $N$ as follows:
\[
  N^{\oplus A} := \{\alpha : A \to N \mid \alpha(a) = 0 \text{ for
    only finitely many elements } a \in A\}.
\]
Of course this agrees with the definition given for abelian groups in
\S{}\ref{subsec:groups:free-abelian-groups}; $N^{\oplus A}$ has an evident $R$-module structure,
obtained by defining, for all $r \in R$ and $a \in A$,
$(r\alpha)(a) := r(\alpha(a))$. For $N = R$ we may define a function
$j : A \to R^{\oplus A}$ by sending $a \in A$ to the function
$j_a : A \to R$:
\[
  (\forall x \in A) : j_a(x) := \begin{cases} 1 & \text{if } x = a, \\
    0 & \text{if } x \neq a. \end{cases}
\]
\begin{prop}{}{prop:rings:free-module}
  $F_R(A) \cong R^{\oplus A}$.
\end{prop}

\begin{proof}
  The proof of this claim matches precisely the proof of
  \Cref{prop:groups:free-abelian-group-concrete-construction} and is left to the reader (\Cref{exer:rings:free-module-proof}).
\end{proof}

The key is that every element of $R^{\oplus A}$ may be written
uniquely as a finite sum $\sum_{a \in A} r_a a$ (shorthand for
$\sum_{a \in A} r_a j(a)$); incidentally, this is how elements of `the
free $R$-module on $A$' are often written — this is legal, by virtue
of \cref{prop:rings:free-module}.

In particular, for $A = \{1, \ldots, n\}$ a finite set,
\cref{prop:rings:free-module} states that the $R$-module
$R^{\oplus n}$, with $j : A \to R^{\oplus n}$ defined by
$j(i) := (0, \ldots, 0, 1, 0, \ldots, 0) \in R^{\oplus n}$ ($1$ in the
$i$-th place), satisfies the universal property for
$F_R(\{1, \ldots, n\})$.  This is all entirely analogous to the story
for $\Z$-modules. The situation becomes a little more interesting if
we switch from $R$-modules to commutative $R$-algebras; the category
is different, so we should expect a different answer. The finite case
is essentially the only one we will need in these notes, so assume
$A = \{1, \ldots, n\}$ is a finite set. In this case, we write $R[A]$
for the polynomial ring $R[x_1, \ldots, x_n]$; we have a set-function
$j : A \to R[A]$, defined by $j(i) = x_i$.

\begin{prop}{}{prop:rings:free-R-algebra}
  $R[A]$ is a free commutative $R$-algebra on the set $A$.
\end{prop}
\begin{proof}
  The statement translates into the following: for every commutative
  $R$-algebra $S$ and every set-function $f : A \to S$, there exists a
  unique $R$-algebra homomorphism $\phi : R[A] \to S$ such that the
  diagram
  \[\begin{tikzcd}
      R[A] \arrow[r, "\phi"] & S \\
      A \arrow[u, "j"] \arrow[ur, "f"']
    \end{tikzcd}\] commutes. Since $S$ is an $R$-algebra, we have a
  fixed homomorphism of rings $\alpha : R \to S$ (cf.\
  \Cref{exam:rings:ring-hom-defines-module}). Then we may construct
  $\phi : R[A] = R[x_1, \ldots, x_n] \to S$ by applying $n$ times the
  `extension property' of \Cref{exam:rings:poly-universal-general}: extend $\alpha$ to $R[x_1]$ so
  as to map $x_1$ to $f(1)$, then to $R[x_1, x_2] = R[x_1][x_2]$ so as
  to map $x_2$ to $f(2)$, etc. Note that each extension is uniquely
  determined by its requirements. This gives $\phi$ as a ring
  homomorphism and shows that it is unique. The reader will verify
  that $\phi$ is also (automatically) an $R$-module homomorphism, and
  hence a homomorphism of $R$-algebras, concluding the proof.
\end{proof}
After the fact, the reader may want to revisit \S{}\ref{subsec:rings:universal-property-of-polynomial-rings} and recognize
that the `universal property of polynomial rings
$\Z[x_1, \ldots, x_n]$' given there was really a version of their role
as free objects in the category of commutative rings, a.k.a.\
commutative $\Z$-algebras.

It is not difficult to identify free objects in the larger category
$R$-$\categ{Alg}$: they consist of `noncommutative polynomial rings'
$R\langle A \rangle$, with variables from the set $A$ but without any
condition relating $ab$ and $ba$ for $a \neq b$ in $A$. More
precisely, $R\langle A \rangle$ is isomorphic to the monoid ring (cf.\
\S{}\ref{subsec:rings:monoid-rings}) over the free monoid on $A$, consisting of all finite strings
of elements in $A$, with operation defined by
concatenation\footnote{This construction is similar to the free group
  on $A$, but without the complication of the presence of inverses and
  of possible cancellations.}. We will encounter this ring again, but
only in the distant future (\Cref{linrep:4.17}).

\subsection{Submodule generated by a subset; Noetherian modules}
\label{subsec:rings:submodule-generated-by-a-subset-noetherian-modules}
Let $M$ be an $R$-module, and let $A \subseteq M$ be a subset of
$M$. By the universal property of free modules, there is a unique
homomorphism of $R$-modules $\phi_A : R^{\oplus A} \to M$. The image
of this homomorphism is a submodule of $M$, the \emph{submodule
  generated by $A$ in $M$}, usually denoted $\langle A \rangle$ (or
$\langle a_1, \ldots, a_n \rangle$ if $A = \{a_1, \ldots, a_n\}$ is
finite). Thus,
\[
  \langle A \rangle = \Bigl\{ \sum_{a \in A} r_a a \Bigm| r_a = 0
  \text{ for only finitely many elements } a \in A \Bigr\}.
\]
It is hopefully clear that $\langle A \rangle$ is the smallest
submodule of $M$ containing $A$.

The module $M$ is \emph{finitely generated} if $M = \langle A \rangle$
for a finite set $A$, that is, if and only if there is a surjective
homomorphism of $R$-modules $R^{\oplus n} \twoheadrightarrow M$ for
some $n$. One of the highlights of \Cref{chap:linalg} will be the
classification of finitely generated modules over PIDs
(\Cref{linalg:5.6}). I have already briefly mentioned the case for $\Z$
(recall that $\Z$ is a PID and $\Z$-modules are nothing but abelian
groups!) back in \S{}\ref{subsec:groups:example-subgroup-generated-by-a-subset}.  Finitely generated $R$-modules are
tremendously important, but they are not as well-behaved as one might
hope at first. For example, it may be that a module $M$ is finitely
generated, but some submodule of $M$ is not finitely generated!

\begin{exam}{}{exam:rings:infinite-poly-ring-fg}
  Let $R = \Z[x_1, x_2, \ldots]$, a polynomial ring on infinitely many
  indeterminates. Then $R$ is finitely generated as an $R$-module:
  indeed, $1$ generates it. However, the ideal $(x_1, x_2, \ldots)$ of
  $R$ generated by all indeterminates is not finitely generated as an
  $R$-module (\Cref{exer:rings:infinite-ideal-not-fg}).
\end{exam}
\begin{defn}{}{defn:rings:noetherian-module}
  An $R$-module $M$ is \emph{Noetherian} if every submodule of $M$ is
  finitely generated as an $R$-module.
\end{defn}
Thus, a ring $R$ is Noetherian in the sense of
\cref{defn:rings:noetherian} if and only if it is Noetherian `as a
module over itself'. The ring in \Cref{exam:rings:infinite-poly-ring-fg} is not Noetherian. We
will study the Noetherian condition more carefully later on
(\S{}\ref{subsec:irred:noetherian-rings-revisited}); but we can already see one reason why this is a good,
`solid' notion.

\begin{prop}{}{prop:rings:noetherian-subquotient}
  Let $M$ be an $R$-module, and let $N$ be a submodule of $M$. Then
  $M$ is Noetherian if and only if both $N$ and $M/N$ are Noetherian.
\end{prop}
\begin{proof}
  If $M$ is Noetherian, then so is $M/N$ (same proof as for
  \Cref{exer:rings:noetherian-homomorphic-image}), and so is $N$ (because every submodule of $N$ is a
  submodule of $M$, so it is finitely generated because $M$ is
  Noetherian). This proves the `only if' part of the statement.

  For the converse, assume $N$ and $M/N$ are Noetherian, and let $P$
  be a submodule of $M$; we have to prove that $P$ is finitely
  generated. Since $P \cap N$ is a submodule of $N$ and $N$ is
  Noetherian, $P \cap N$ is finitely generated. By the `second
  isomorphism theorem', \cref{prop:rings:second-iso},
  $P/(P \cap N) \cong (P + N)/N$, and hence $P/(P \cap N)$ is
  isomorphic to a submodule of $M/N$. Since $M/N$ is Noetherian, this
  shows that $P/(P \cap N)$ is finitely generated. It follows that $P$
  itself is finitely generated, by \Cref{exer:rings:fg-submodule-quotient}.
\end{proof}
\begin{coro}{}{coro:rings:fg-module-over-noetherian-is-noetherian}
  Let $R$ be a Noetherian ring, and let $M$ be a finitely generated
  $R$-module. Then $M$ is Noetherian (as an $R$-module).
\end{coro}

\begin{proof}
  Indeed, by hypothesis there is an onto homomorphism
  $R^{\oplus n} \twoheadrightarrow M$ of $R$-modules; hence (by the
  first isomorphism theorem, \Cref{coro:rings:module-first-iso}) $M$ is isomorphic to a
  quotient of $R^{\oplus n}$. By
  \cref{prop:rings:noetherian-subquotient}, it suffices to prove that
  $R^{\oplus n}$ is Noetherian. This may be done by induction. The
  statement is true for $n = 1$ by hypothesis. For $n > 1$, assume we
  know that $R^{\oplus(n-1)}$ is Noetherian; since $R^{\oplus(n-1)}$
  may be viewed as a submodule of $R^{\oplus n}$, in such a way that
  $R^{\oplus n}/R^{\oplus(n-1)} \cong R$ (\Cref{exer:rings:free-module-quotient}), and $R$ is
  Noetherian, it follows that $R^{\oplus n}$ is Noetherian, again by
  applying \cref{prop:rings:noetherian-subquotient}.
\end{proof}

\subsection{Finitely generated vs. finite type}
\label{subsec:rings:finitely-generated-vs-finite-type}
If $S$ is an $R$-algebra, it may be `finitely generated' in two very
different ways: as an $R$-module and as an $R$-algebra. It is
important to keep these two concepts well distinct, although
unfortunately the language used to express them is very similar.

The following definitions differ in three small details\ldots
\begin{quote}
  ``$S$ is finitely generated as a module over $R$ if there is an onto
  homomorphism of $R$-modules from the free $R$-module on a finite set
  to $S$.''
\end{quote}
\begin{quote}
  ``$S$ is finitely generated as an algebra over $R$ if there is an
  onto homomorphism of $R$-algebras from the free $R$-algebra on a
  finite set to $S$.''
\end{quote}
The mathematical difference is more substantial than it may appear. As
we have seen in \S{}\ref{subsec:rings:free-modules-and-free-algebras}, the free $R$-module over a finite set
$A = \{1, \ldots, n\}$ is isomorphic to $R^{\oplus n}$; the free
commutative $R$-algebra over $A$ is isomorphic to
$R[x_1, \ldots, x_n]$. Thus, a commutative\footnote{We are mostly
  interested in the commutative case, so I will make this hypothesis
  here; the only change in the general case is typographical: $\cdots$
  rather than $[\cdots]$. Also, note that a commutative ring is
  finitely generated as an algebra if and only if it is finitely
  generated as a commutative algebra; cf.\ \Cref{exer:rings:commutative-algebra-fg}.} ring $S$ is
finitely generated as an $R$-module if there is an onto homomorphism
of $R$-modules $R^{\oplus n} \twoheadrightarrow S$ for some $n$; it is
finitely generated as an $R$-algebra if there is an onto homomorphism
of $R$-algebras $R[x_1, \ldots, x_n] \twoheadrightarrow S$ for some
$n$. In other words, $S$ is finitely generated as an $R$-module if and
only if $S \cong R^{\oplus n}/M$ for some $n$ and a submodule $M$ of
$R^{\oplus n}$; it is a finite-type $R$-algebra if and only if
$S \cong R[x_1, \ldots, x_n]/I$ for some $n$ and an ideal $I$ of
$R[x_1, \ldots, x_n]$. We say that $S$ is \emph{finite}\footnote{This
  is particularly unfortunate, since $S$ may very well be an infinite
  set.} in the first case and of \emph{finite type} in the second. It
is clear that `finite' $\implies$ `finite type'; it should be just as
clear that the converse does not hold.
\begin{exam}{}{}
  The polynomial ring $R[x]$ is a finite-type $R$-algebra, but it is
  not finite as an $R$-module.
\end{exam}
The distinction, while macroscopic in general, may evaporate in
special, important cases. For example, one can prove that if $k$ and
$K$ are fields and $k \subseteq K$, then $K$ is of finite type over
$k$ if and only if it is in fact finite as a $k$-module (that is, it
is a finite-dimensional $k$-vector space). This is one version of
Hilbert's Nullstellensatz, a deep result we already mentioned in
\Cref{exam:rings:calculus-of-ideals}5 and that we will prove (in an important class of
examples) in \S{}\ref{subsec:fields:the-nullstellensatz}.

\name{David Hilbert}'s name is associated to another important result
concerning finite-type $R$-algebras: if $R$ is Noetherian (as a ring,
that is, as an $R$-module) and $S$ is a finite-type $R$-algebra, then
$S$ is also Noetherian (as a ring, that is, as an $S$-module). This is
an immediate consequence of the so-called \emph{Hilbert's basis
  theorem}. The proof of Hilbert's basis theorem is completely
elementary: it could be given here as an exercise, with a few key
hints; we will see it in \S{}\ref{subsec:irred:noetherian-rings-revisited}.

\subsection*{Exercises}

\begin{exer}{\exerused}{exer:rings:free-module-proof}
  Prove \cref{prop:rings:free-module}. [\S{}\ref{subsec:rings:free-modules-and-free-algebras}]
\end{exer}
\begin{exer}{}{exer:rings:M-not-iso-MplusM}
  Prove or disprove that if $R$ is a ring and $M$ is a nonzero
  $R$-module, then $M$ is not isomorphic to $M \oplus M$.
\end{exer}
\begin{exer}{}{exer:rings:projection-decomposition}
  Let $R$ be a ring, $M$ an $R$-module, and $p : M \to M$ an
  $R$-module homomorphism such that $p^2 = p$. (Such a map is called a
  \emph{projection}.) Prove that
  $M \cong \ker p \oplus \image p$.
\end{exer}
\begin{exer}{\exerused}{exer:rings:free-module-quotient}
  Let $R$ be a ring, and let $n > 1$. View $R^{\oplus(n-1)}$ as a
  submodule of $R^{\oplus n}$, via the injective homomorphism
  $R^{\oplus(n-1)} \hookrightarrow R^{\oplus n}$ defined by
  $(r_1, \ldots, r_{n-1}) \mapsto (r_1, \ldots, r_{n-1}, 0)$. Give a
  one-line proof that $R^{\oplus n}/R^{\oplus(n-1)} \cong
  R$. [\S{}\ref{subsec:rings:submodule-generated-by-a-subset-noetherian-modules}]
\end{exer}
\begin{exer}{\exerused}{exer:rings:free-module-product-sets}
  (Notation as in \S{}\ref{subsec:rings:free-modules-and-free-algebras}.) For any ring $R$ and any two sets $A_1$,
  $A_2$, prove that
  $(R^{\oplus A_1})^{\oplus A_2} \cong R^{\oplus(A_1 \times
    A_2)}$. [\S{}\ref{subsec:linrep:adjunction-with-hom-and-explicit-computations}]
\end{exer}
\begin{exer}{\lnot}{exer:rings:hom-free-dual}
  Let $R$ be a ring, and let $F = R^{\oplus n}$ be a finitely
  generated free $R$-module. Prove that
  $\Hom_{R\text{-}\categ{Mod}}(F, R) \cong F$. On the other hand, find
  an example of a ring $R$ and a nonzero $R$-module $M$ such that
  $\Hom_{R\text{-}\categ{Mod}}(M, R) = 0$. [\ref{exer:rings:hom-free-product}]
\end{exer}
\begin{exer}{\exerused}{exer:rings:product-coproduct-infinite}
  Let $A$ be any set.
  \begin{itemize}
  \item For any family $\{M_a\}_{a \in A}$ of modules over a ring $R$,
    define the product $\prod_{a \in A} M_a$ and coproduct
    $\bigoplus_{a \in A} M_a$. If $M_a \cong R$ for all $a \in A$,
    these are denoted $R^A$, $R^{\oplus A}$, respectively.
  \item Prove that $\Z^{\N} \not\cong \Z^{\oplus \N}$. (Hint:
    Cardinality.)
  \end{itemize} [\S{}\ref{subsec:rings:products-and-coproducts}, \ref{exer:rings:hom-free-product}]
\end{exer}
\begin{exer}{}{exer:rings:hom-free-product}
  Let $R$ be a ring. If $A$ is any set, prove that
  $\Hom_{R\text{-}\categ{Mod}}(R^{\oplus A}, R)$ satisfies the
  universal property for the product of the family
  $\{R_a\}_{a \in A}$, where $R_a \cong R$ for all $a$; thus,
  $\Hom_{R\text{-}\categ{Mod}}(R^{\oplus A}, R) \cong R^A$. Conclude
  that $\Hom_{R\text{-}\categ{Mod}}(R^{\oplus A}, R)$ is not
  isomorphic to $R^{\oplus A}$ in general (cf.\ Exercises~6.6
  and~6.7).
\end{exer}
\begin{exer}{\lnot}{exer:rings:free-projective}
  Let $R$ be a ring, $F$ a nonzero free $R$-module, and let
  $\phi : M \to N$ be a homomorphism of $R$-modules. Prove that $\phi$
  is onto if and only if for all $R$-module homomorphisms
  $\alpha : F \to N$ there exists an $R$-module homomorphism
  $\beta : F \to M$ such that $\alpha = \phi \circ \beta$. (Free
  modules are projective, as we will see in \Cref{chap:linrep}.) [\ref{exer:rings:7.8},
  \ref{exer:linalg:direct-summand-of-free-over-local-is-free}]
\end{exer}
\begin{exer}{\exerused}{exer:rings:fibered-product}
  (Cf.\ \Cref{exer:prelim:fiber-product}.) Let $M$, $N$, and $Z$ be $R$-modules, and
  let $\mu : M \to Z$, $\nu : N \to Z$ be homomorphisms of
  $R$-modules. Prove that $R$-$\categ{Mod}$ has `fibered products':
  there exists an $R$-module $M \times_Z N$ with $R$-module
  homomorphisms $\pi_M : M \times_Z N \to M$,
  $\pi_N : M \times_Z N \to N$, such that
  $\mu \circ \pi_M = \nu \circ \pi_N$, and which is universal with
  respect to this requirement. The module $M \times_Z N$ may be called
  the \emph{pull-back} of $M$ along $\nu$ (or of $N$ along $\mu$,
  since the construction is symmetric). `Fiber diagrams'
  \[\begin{tikzcd}
      M \times_Z N \arrow[r, "\pi_N"] \arrow[d, "\pi_M"'] \arrow[dr, phantom, "\lrcorner", very near start] & N \arrow[d, "\nu"] \\
      M \arrow[r, "\mu"'] & Z
    \end{tikzcd}\] are commutative, but `even better' than
  commutative; they are often decorated by a square, as shown
  here. [\S{}\ref{subsec:rings:products-and-coproducts}, \ref{exer:rings:fibered-coproduct}, \S{}\ref{subsec:homol:products-coproducts-and-direct-sums}]
\end{exer}
\begin{exer}{\exerused}{exer:rings:fibered-coproduct}
  Define a notion of fibered coproduct of two $R$-modules $M$, $N$,
  along an $R$-module $A$, in the style of \Cref{exer:rings:fibered-product} (and cf.\
  \Cref{exer:prelim:fiber-product}):
  \[\begin{tikzcd}
      A \arrow[r, "\nu"] \arrow[d, "\mu"'] \arrow[dr, phantom, "\ulcorner", very near end] & N \arrow[d] \\
      M \arrow[r] & M \oplus_A N
    \end{tikzcd}\] Prove that fibered coproducts exist in
  $R$-$\categ{Mod}$. The fibered coproduct $M \oplus_A N$ is called
  the \emph{push-out} of $M$ along $\nu$ (or of $N$ along
  $\mu$). [\S{}\ref{subsec:rings:products-and-coproducts}]
\end{exer}
\begin{exer}{}{exer:rings:prove-prop-kernels-cokernels}
  Prove \cref{prop:rings:kernels-cokernels}.
\end{exer}
\begin{exer}{}{exer:rings:homomorphic-image-fg}
  Prove that every homomorphic image of a finitely generated module is
  finitely generated.
\end{exer}
\begin{exer}{\exerused}{exer:rings:infinite-ideal-not-fg}
  Prove that the ideal $(x_1, x_2, \ldots)$ of the ring
  $R = \Z[x_1, x_2, \ldots]$ is not finitely generated (as an ideal,
  i.e., as an $R$-module). [\S{}\ref{subsec:rings:submodule-generated-by-a-subset-noetherian-modules}]
\end{exer}
\begin{exer}{\exerused}{exer:rings:commutative-algebra-fg}
  Let $R$ be a commutative ring. Prove that a commutative $R$-algebra
  $S$ is finitely generated as an algebra over $R$ if and only if it
  is finitely generated as a commutative algebra over $R$. (Cf.\
  \S{}\ref{subsec:rings:finitely-generated-vs-finite-type}.) [\S{}\ref{subsec:rings:finitely-generated-vs-finite-type}]
\end{exer}
\begin{exer}{\exerused}{exer:rings:cyclic-module}
  Let $R$ be a ring. A (left-)$R$-module $M$ is \emph{cyclic} if
  $M = \langle m \rangle$ for some $m \in M$. Prove that simple
  modules (cf.\ \Cref{exer:rings:schur}) are cyclic. Prove that an $R$-module $M$
  is cyclic if and only if $M \cong R/I$ for some (left-)ideal
  $I$. Prove that every quotient of a cyclic module is cyclic. [\ref{exer:rings:hom-cyclic},
  \S{}\ref{subsec:linalg:torsion}]
\end{exer}
\begin{exer}{\lnot}{exer:rings:hom-cyclic}
  Let $M$ be a cyclic $R$-module, so that $M \cong R/I$ for a
  (left-)ideal $I$ (\Cref{exer:rings:cyclic-module}), and let $N$ be another $R$-module.
  \begin{itemize}
  \item Prove that
    $\Hom_{R\text{-}\categ{Mod}}(M, N) \cong \{n \in N \mid (\forall a
    \in I), an = 0\}$.
  \item For $a, b \in \Z$, prove that
    $\Hom_{\categ{Ab}}(\Z/a\Z, \Z/b\Z) \cong \Z/\gcd(a,b)\Z$.
  \end{itemize} [\ref{exer:rings:7.7}]
\end{exer}
\begin{exer}{\exerused}{exer:rings:fg-submodule-quotient}
  Let $M$ be an $R$-module, and let $N$ be a submodule of $M$. Prove
  that if $N$ and $M/N$ are both finitely generated, then $M$ is
  finitely generated. [\S{}\ref{subsec:rings:submodule-generated-by-a-subset-noetherian-modules}]
\end{exer}

\section{Complexes and homology}
\label{sec:rings:complexes-and-homology}

In many contexts, modules arise not `one at a time' but in whole
series: for example, a real manifold of dimension $d$ has one
`homology' group for each dimension from $0$ to $d$. It is necessary
to develop a language capable of dealing with whole sequences of
modules at once. This is the language of homological algebra, of which
we will get a tiny taste in this section, and a slightly heartier
course in \Cref{chap:homol}.

\subsection{Complexes and exact sequences}
\label{subsec:rings:complexes-and-exact-sequences}
A \emph{chain complex} of $R$-modules (or, for simplicity, a
\emph{complex}) is a sequence of $R$-modules and $R$-module
homomorphisms
\[
  \cdots \xrightarrow{d_{i+2}} M_{i+1} \xrightarrow{d_{i+1}} M_i
  \xrightarrow{d_i} M_{i-1} \xrightarrow{d_{i-1}} \cdots
\]
such that $(\forall i) : d_i \circ d_{i+1} = 0$. The notation
$(M_\bullet, d_\bullet)$ may be used to denote a complex, or simply
$M_\bullet$ for simplicity (but do not forget that the homomorphisms
$d_i$ are part of the information carried by a complex).

A complex may be infinite in both directions; `tails' of $0$'s are
(usually) omitted. Several possible alternative conventions may be
used: for example, indices may be increasing rather than decreasing,
giving a \emph{cochain complex} (whose homology is called
\emph{cohomology}; this will be our choice in \Cref{chap:homol}). Such
choices are clearly mathematically immaterial, at least for the simple
considerations which follow.

The homomorphisms $d_i$ are called \emph{boundary maps}, or
\emph{differentials}, due to important examples from geometry. Note
that the defining condition $d_i \circ d_{i+1} = 0$ is equivalent to
the requirement $\image d_{i+1} \subseteq \ker d_i$.
Almost the whole point about complexes is to `measure' the difference
between $\image d_{i+1}$ and $\ker d_i$, which is called
the \emph{homology} of the complex (cf.\ \S{}\ref{subsec:rings:homology-and-the-snake-lemma}). We say that a
complex is \emph{exact} `at $M_i$' if it has no homology there; that
is, $\image d_{i+1} = \ker d_i$. For example, if $M_i$ is a
trivial module (usually denoted simply by $0$), then the complex is
necessarily exact at $M_i$, since then
$\image d_{i+1} = \ker d_i = 0$. A complex is \emph{exact}
and is often called an \emph{exact sequence} if it is exact at all its
modules.
\begin{exam}{}{}
  A complex $\cdots \to 0 \to L \xrightarrow{\alpha} M \to \cdots$ is
  exact at $L$ if and only if $\alpha$ is a monomorphism. Indeed,
  exactness at $L$ is equivalent to $\ker\alpha = $ image of the
  trivial homomorphism $0 \to L$, that is, to $\ker\alpha = 0$. This
  is equivalent to the injectivity of $\alpha$
  (\cref{prop:rings:kernels-cokernels}).
\end{exam}
\begin{exam}{}{exam:rings:exact-at-n-epimorphism}
  A complex $\cdots \to M \xrightarrow{\beta} N \to 0 \to \cdots$ is
  exact at $N$ if and only if $\beta$ is an epimorphism. Indeed, the
  complex is exact at $N$ if and only if $\image\beta = $
  kernel of the trivial homomorphism $N \to 0$, that is,
  $\image\beta = N$.
\end{exam}
\begin{defn}{}{defn:rings:short-exact-sequence}
  A \emph{short exact sequence} is an exact complex of the form
  \[
    0 \to L \xrightarrow{\alpha} M \xrightarrow{\beta} N \to 0.
  \]
\end{defn}
As seen in the previous two examples, exactness at $L$ and $N$ is
equivalent to $\alpha$ being injective and $\beta$ being
surjective. The extra piece of data carried by a short exact sequence
is the exactness at $M$, that is,
$\image\alpha = \ker\beta$; by the first isomorphism
theorem (\Cref{coro:rings:module-first-iso}), we then have
$N \cong M/\ker\beta = M/\image\alpha$.

All in all, we have good material to work on some more Pavlovian
conditioning: at the sight of a short exact sequence as above, the
reader should instinctively identify $L$ with a submodule of $M$ (via
the injective map $\alpha$) and $N$ with the quotient $M/L$ (via the
isomorphism induced by the surjective map $\beta$, under the auspices
of the first isomorphism theorem).  Short exact sequences abound in
nature. For example, a single homomorphism $\phi : M \to M'$ gives
rise immediately to a short exact sequence
\[
  0 \to \ker\phi \to M \to \image\phi \to 0.
\]
In fact, one important reason to focus on short exact sequences is
that this observation allows us to break up every exact complex into a
large number of short exact sequences. The diagonal sequences in the
following diagram are short exact sequences, and they interlock nicely
by the exactness of the horizontal complex. This observation
simplifies many arguments; cf.\ for example \Cref{exer:rings:7.5}.

\subsection{Split exact sequences}
\label{subsec:rings:split-exact-sequences}
A particular case of short exact sequence arises by considering the
second projection from a direct sum: $M_1 \oplus M_2 \to M_2$; there
is then an exact sequence
\[
  0 \to M_1 \to M_1 \oplus M_2 \to M_2 \to 0,
\]
obtained by identifying $M_1$ with the kernel of the projection. These
short exact sequences are said to `split'; more generally, a short
exact sequence
\[
  0 \to M_1 \to N \to M_2 \to 0
\]
`splits' if it is isomorphic to one of these sequences in the sense
that there is a commutative diagram
\[\begin{tikzcd}
    0 \arrow[r] & M_1 \arrow[r] \arrow[d, "\sim"'] & N \arrow[r] \arrow[d, "\sim"'] & M_2 \arrow[r] \arrow[d, "\sim"'] & 0 \\
    0 \arrow[r] & M_1' \arrow[r] & M_1' \oplus M_2' \arrow[r] & M_2'
    \arrow[r] & 0
  \end{tikzcd}\] in which the vertical maps are all
isomorphisms\footnote{In fact, this last requirement is somewhat
  redundant; cf.\ \Cref{exer:rings:7.11}.}.

\begin{exam}{}{}
  The exact sequence of $\Z$-modules
  $0 \to \Z \xrightarrow{\cdot 2} \Z \to \Z/2\Z \to 0$ is not split.
\end{exam}
Splitting sequences give us the opportunity to go back to a question
we left dangling at the end of \S{}\ref{subsec:rings:kernels-and-cokernels}: what should we make of the
condition of `having a left- (resp., right-)inverse' for a
homomorphism? We realized that this condition is stronger than the
requirement of being a monomorphism (resp., an epimorphism); can we
give a more explicit description of such morphisms?

\begin{prop}{}{prop:rings:split-sequence}
  Let $\phi : M \to N$ be an $R$-module homomorphism. Then
  \begin{itemize}
  \item $\phi$ has a left-inverse if and only if the sequence
    $0 \to M \xrightarrow{\phi} N \to \coker\phi \to 0$
    splits.
  \item $\phi$ has a right-inverse if and only if the sequence
    $0 \to \ker\phi \to M \xrightarrow{\phi} N \to 0$ splits.
  \end{itemize}
\end{prop}
\begin{proof}
  I will prove the first part and leave the other as an exercise to
  the reader (\Cref{exer:rings:7.6}). If the sequence splits, then $\phi$ may
  be identified with the embedding of $M$ into a direct sum
  $M \oplus M'$, and the projection $M \oplus M' \to M$ gives a
  left-inverse of $\phi$. Conversely, assume that $\phi$ has a
  left-inverse $\psi$:
  \[\begin{tikzcd}
      M \arrow[r, "\phi"] & N \arrow[l, bend left, "\psi"] \\
      & M \arrow[u, equal] \arrow[ul, "\id"]
    \end{tikzcd}\] Then I claim that $N$ is isomorphic to
  $M \oplus \ker\psi$ and that $\phi$ corresponds to the
  identification of $M$ with the first factor:
  $M \to M \oplus \ker\psi \cong N$. The isomorphism
  $M \oplus \ker\psi \to N$ is given by $(m, k) \mapsto \phi(m) + k$;
  its inverse $N \to M \oplus \ker\psi$ is
  $n \mapsto (\psi(n), n - \phi\psi(n))$. The element
  $n - \phi\psi(n)$ is in $\ker\psi$ as it should be, since
  $\psi(n - \phi\psi(n)) = \psi(n) - \psi\phi\psi(n) = \psi(n) -
  \psi(n) = 0$. All necessary verifications are immediate and are left
  to the reader.
\end{proof}
Because of \cref{prop:rings:split-sequence}, $R$-module homomorphisms
with a left-inverse are called \emph{split monomorphisms}, and
homomorphisms with a right-inverse are called \emph{split
  epimorphisms}.

We will come back to split exact sequences (in the more demanding
context of $\categ{Grp}$) in \S{}\ref{subsec:groups2:exact-sequences-of-groups-extension-problem} and then later again when we
return to modules and, more generally, to abelian categories
(Chapters~VIII and~IX).

\subsection{Homology and the snake lemma}
\label{subsec:rings:homology-and-the-snake-lemma}

\begin{defn}{}{defn:rings:homology}
  The \emph{$i$-th homology} of a complex
  $M_\bullet : \cdots \xrightarrow{d_{i+2}} M_{i+1}
  \xrightarrow{d_{i+1}} M_i \xrightarrow{d_i} M_{i-1}
  \xrightarrow{d_{i-1}} \cdots$ of $R$-modules is the $R$-module
  \[
    H_i(M_\bullet) := \frac{\ker d_i}{\image d_{i+1}}.
  \]
\end{defn}

That is, $H_i(M_\bullet)$ is a module capturing the failure of
exactness at $M_i$. Of course
$H_i(M_\bullet) = 0 \iff \image d_{i+1} = \ker d_i \iff$
the complex $M_\bullet$ is exact at $M_i$: that is, the homology
modules are a measure of the `failure of a complex from being exact'.
\begin{exam}{}{exam:rings:homology-generalizes-ker-coker}
  In fact, homology should be thought of as a (vast) generalization of
  the notions of kernel and cokernel. Indeed, consider the (very)
  particular case in which $M_\bullet$ is the complex
  $0 \to M_1 \xrightarrow{\phi} M_0 \to 0$. Then
  $H_1(M_\bullet) \cong \ker\phi$ and
  $H_0(M_\bullet) \cong \coker\phi$.
\end{exam}
I will end this very brief excursion into more abstract territories by
indicating how a commutative diagram involving two short exact
sequences generates a `long exact sequence' in homology. This is
actually a particular case of a more general construction ---
according to which a suitable commutative diagram involving three
complexes yields a really long `long exact homology sequence'. We will
come back to this general construction when we deal more extensively
with homological algebra in \Cref{chap:homol}. The reader is also likely to
learn about it in a course on algebraic topology, where this fact is
put to impressive use in studying invariants of manifolds.  In the
simple form we will analyze, this is affectionately known as the
\emph{snake lemma}. Consider two short exact sequences linked by
homomorphisms, so as to form a commutative diagram\footnote{In fact,
  it is better to view this diagram as three (very short) complexes
  linked by $R$-module homomorphisms $\alpha_i$, $\beta_i$ so that
  `the rows are exact'. One can define a category of complexes, and
  this diagram is nothing but a `short exact sequence of complexes'.}:
\[\begin{tikzcd}
    0 \arrow[r] & L_1 \arrow[r, "\alpha_1"] \arrow[d, "\lambda"'] & M_1 \arrow[r, "\beta_1"] \arrow[d, "\mu"'] & N_1 \arrow[r] \arrow[d, "\nu"'] & 0 \\
    0 \arrow[r] & L_0 \arrow[r, "\alpha_0"'] & M_0 \arrow[r,
    "\beta_0"'] & N_0 \arrow[r] & 0
  \end{tikzcd}\]
\begin{lem}{The snake lemma}{lem:rings:snake-lemma}
  With notation as above, there is an exact sequence
  \[
    0 \to \ker\lambda \to \ker\mu \to \ker\nu \xrightarrow{\delta}
    \coker\lambda \to \coker\mu \to
    \coker\nu \to 0.
  \]
\end{lem}
\begin{rem}{}{}
  Most of the homomorphisms in this sequence are induced in a
  completely straightforward way from the corresponding homomorphisms
  $\lambda$, $\mu$, $\nu$. The one `surprising' homomorphism is the
  one denoted $\delta$; I will discuss its definition below.
\end{rem}
\begin{rem}{}{rem:rings:snake-lemma-homology-restatement}
  In view of \Cref{exam:rings:homology-generalizes-ker-coker}, we could have written the sequence in this
  statement as
  \[
    0 \to H_1(L_\bullet) \to H_1(M_\bullet) \to H_1(N_\bullet)
    \xrightarrow{\delta} H_0(L_\bullet) \to H_0(M_\bullet) \to
    H_0(N_\bullet) \to 0
  \]
  where $L_\bullet$ is the complex
  $0 \to L_1 \xrightarrow{\lambda} L_0 \to 0$, etc. The snake lemma
  generalizes to arbitrary complexes $L_\bullet$, $M_\bullet$,
  $N_\bullet$, producing a `long exact homology sequence' of which
  this is just the tail end. As mentioned above, we will discuss this
  rather straightforward generalization later (\S{}\ref{subsec:homol:the-long-exact-cohomology-sequence}).
\end{rem}
\begin{rem}{}{rem:rings:snake-lemma-relaxed-hypotheses}
  A popular version of the snake lemma does not assume that $\alpha_1$
  is injective and $\beta_0$ is surjective: that is, we could consider
  a commutative diagram of exact sequences
  \[\begin{tikzcd}
      L_1 \arrow[r, "\alpha_1"] \arrow[d, "\lambda"'] & M_1 \arrow[r, "\beta_1"] \arrow[d, "\mu"'] & N_1 \arrow[r] \arrow[d, "\nu"'] & 0 \\
      0 \arrow[r] & L_0 \arrow[r, "\alpha_0"'] & M_0 \arrow[r,
      "\beta_0"'] & N_0
    \end{tikzcd}\] The lemma will then state that there is `only' an
  exact sequence
  $\ker\lambda \to \ker\mu \to \ker\nu \xrightarrow{\delta}
  \coker\lambda \to \coker\mu \to
  \coker\nu$.
\end{rem}
Proving the snake lemma is something that should not be done in
public, and it is notoriously useless to write down the details of the
verification for others to read: the details are all essentially
obvious, but they lead quickly to a notational quagmire. Such proofs
are collectively known as the sport of \emph{diagram chase}, best
executed by pointing several fingers at different parts of a diagram
on a blackboard, while enunciating the elements one is manipulating
and stating their fate\footnote{Real purists chase diagrams in
  arbitrary categories, thus without the benefit of talking about
  `elements', and we will practice this skill later on
  (\Cref{chap:homol}). For example, the snake lemma can be proven by
  appealing to universal property after universal property of kernels
  and cokernels, without ever choosing elements anywhere. But the
  performing technique of pointing fingers at a board while
  monologuing through the argument remains essentially the
  same.}. Nevertheless, I should explain where the `connecting'
homomorphism $\delta$ comes from, since this is the heart of the
statement of the snake lemma and of its proof.

Here is the whole diagram, including kernels and cokernels; thus,
columns are exact (as well as the two original sequences, placed
horizontally):
\[\begin{tikzcd}
    & \ker\lambda \arrow[r] \arrow[d] & \ker\mu \arrow[r] \arrow[d] & \ker\nu \arrow[d] \arrow[dll, "\delta", rounded corners, to path={ -- ([xshift=2ex]\tikztostart.east) |- ([yshift=-0.5ex]\tikztotarget.north) -- (\tikztotarget)}] \\
    0 \arrow[r] & L_1 \arrow[r, "\alpha_1"] \arrow[d, "\lambda"'] & M_1 \arrow[r, "\beta_1"] \arrow[d, "\mu"'] & N_1 \arrow[r] \arrow[d, "\nu"'] & 0 \\
    0 \arrow[r] & L_0 \arrow[r, "\alpha_0"'] \arrow[d] & M_0 \arrow[r, "\beta_0"'] \arrow[d] & N_0 \arrow[r] \arrow[d] & 0 \\
    & \coker\lambda \arrow[r] & \coker\mu
    \arrow[r] & \coker\nu \arrow[r] & 0
  \end{tikzcd}\] By the way, I trust that the reader now sees why this
lemma is called the snake lemma.  \emph{Definition of the connecting
  homomorphism $\delta$.} Let $a \in \ker\nu$. I claim that $a$ can be
mapped through the diagram all the way to
$\coker\lambda$: view $a$ as an element $b$ of $N_1$;
since $\beta_1$ is surjective, $\exists\, c \in M_1$ mapping to $b$;
let $d = \mu(c)$ be the image of $c$ in $M_0$; by commutativity $d$
maps to $\nu(b) = 0$ in $N_0$, so
$d \in \ker\beta_0 = \image\alpha_0$; therefore
$\exists\, e \in L_0$ mapping to $d$; let
$f \in \coker\lambda$ be the image of $e$. Set
$\delta(a) := f$.

Is this well-defined? The choice of $c$ is not unique. Suppose $c'$
also maps to $b$: then $\beta_1(c' - c) = 0$, so by exactness
$\exists\, g \in L_1$ with $c' - c = \alpha_1(g)$. Since columns are
complexes, $\lambda(g)$ maps to
$\mu(c' - c) = \mu(\alpha_1(g)) = \alpha_0(\lambda(g))$, so changing
$c$ to $c'$ modifies $e$ to $e + \lambda(g)$ and hence $f$ to
$f + 0 = f$ (since $g$ dies in $\coker\lambda$). Thus
$\delta$ is well-defined.

This is a tiny part of the proof of the snake lemma, but it probably
suffices to demonstrate why reading a written-out version of a diagram
chase may be supremely uninformative. The rest of the proof (left to
the reader — I am not listing this as an official exercise for fear
that someone might actually turn a solution in for grading) amounts to
many, many similar arguments. The definition of the maps induced on
kernels and cokernels is substantially less challenging than the
definition of the connecting morphism $\delta$ described
above. Exactness at most spots in the sequence is also reasonably
straightforward; most of the work will go into proving exactness at
$\ker\nu$ and $\coker\lambda$.

Dear reader: don't shy away from trying this, for it is excellent,
indispensable practice. Miss this opportunity and you will forever
feel unsure about such manipulations.  The snake lemma streamlines
several facts, which would not be hard to prove individually, but
become really straightforward once the lemma is settled. For example,

\begin{coro}{}{}
  In the same situation presented in the snake lemma (notation as in
  \S{}\ref{subsec:rings:homology-and-the-snake-lemma}), assume that $\mu$ is surjective and $\nu$ is
  injective. Then $\lambda$ is surjective and $\nu$ is an isomorphism.
\end{coro}

\begin{proof}
  Indeed, $\mu$ surjective $\implies$ $\coker\mu = 0$;
  $\nu$ injective $\implies$ $\ker\nu = 0$
  (\cref{prop:rings:kernels-cokernels}). Feeding this information into
  the sequence of the snake lemma gives an exact sequence
  \[0 \to \ker\lambda \to \ker\mu \to 0 \to
    \coker\lambda \to 0 \to \coker\nu \to
    0.\] Exactness implies
  $\coker\lambda = \coker\nu = 0$
  (\Cref{exer:rings:7.1}); hence $\lambda$ and $\nu$ are surjective, with the
  stated consequences.
\end{proof}

Several more such statements may be proved just as easily; the reader
should experiment to his or her heart's content.

\subsection*{Exercises}
\begin{exer}{\exerused}{exer:rings:7.1}
  Assume that the complex $\cdots \to 0 \to M \to 0 \to \cdots$ is
  exact. Prove that $M \cong 0$. [\S{}\ref{subsec:rings:homology-and-the-snake-lemma}]
\end{exer}
\begin{exer}{}{exer:rings:7.2}
  Assume that the complex $\cdots \to 0 \to M \to M' \to 0 \to \cdots$
  is exact. Prove that $M \cong M'$.
\end{exer}
\begin{exer}{}{exer:rings:7.3}
  Assume that the complex
  $\cdots \to 0 \to L \to M \xrightarrow{\phi} M' \to N \to 0 \to
  \cdots$ is exact. Show that, up to natural identifications,
  $L = \ker\phi$ and $N = \coker\phi$.
\end{exer}
\begin{exer}{}{exer:rings:7.4}
  Construct short exact sequences of $\Z$-modules
  \[0 \to \Z^{\oplus\N} \to \Z^{\oplus\N} \to \Z \to 0\] and
  \[0 \to \Z^{\oplus\N} \to \Z^{\oplus\N} \to \Z^{\oplus\N} \to 0.\]
  (Hint: David Hilbert's Grand Hotel.)
\end{exer}
\begin{exer}{\exerused}{exer:rings:7.5}
  Assume that the complex $\cdots \to L \to M \to N \to \cdots$ is
  exact and that $L$ and $N$ are Noetherian. Prove that $M$ is
  Noetherian. [\S{}\ref{subsec:rings:complexes-and-exact-sequences}]
\end{exer}
\begin{exer}{\exerused}{exer:rings:7.6}
  Prove the ``split epimorphism'' part of
  \cref{prop:rings:split-sequence}. [\S{}\ref{subsec:rings:split-exact-sequences}]
\end{exer}
\begin{exer}{\exerused}{exer:rings:7.7}
  Let $0 \to M \to N \to P \to 0$ be a short exact sequence of
  $R$-modules, and let $L$ be an $R$-module.
  \begin{enumerate}[label=(\roman*)]
  \item Prove that there is an exact sequence\footnote{In general,
      this will be a sequence of abelian groups; if $R$ is
      commutative, so that each $\Hom_{R\text{-Mod}}$ is an $R$-module
      (\S{}\ref{subsec:rings:the-category-r-mod}), then it will be an exact sequence of $R$-modules.}
    \[\Hom_{R\text{-Mod}}(P,L) \to \Hom_{R\text{-Mod}}(N,L) \to
      \Hom_{R\text{-Mod}}(M,L).\]
  \item Redo \Cref{exer:rings:hom-cyclic}. (Use the exact sequence
    $0 \to I \to R \to R/I \to 0$.)
  \item Construct an example showing that the rightmost homomorphism
    in (i) need not be onto.
  \item Show that if the original sequence splits, then the rightmost
    homomorphism in (i) is onto.
  \end{enumerate} [\ref{exer:rings:7.9}, \ref{exer:linrep:onto-hom-explicit-restriction-extension-scalars}, \S{}\ref{subsec:linrep:adjunction-again}]
\end{exer}
\begin{exer}{\exerused}{exer:rings:7.8}
  Prove that every exact sequence $0 \to M \to N \to F \to 0$ of
  $R$-modules, with $F$ free, splits. (Hint: \Cref{exer:rings:free-projective}.)
  [\S{}\ref{subsec:linrep:duality-and-exactness}]
\end{exer}
\begin{exer}{}{exer:rings:7.9}
  Let $0 \to M \to N \to F \to 0$ be a short exact sequence of
  $R$-modules, with $F$ free, and let $L$ be an $R$-module. Prove that
  there is an exact sequence
  \[\Hom_{R\text{-Mod}}(F,L) \to \Hom_{R\text{-Mod}}(N,L) \to
    \Hom_{R\text{-Mod}}(M,L) \to 0.\] (Cf.\ \Cref{exer:rings:7.7}.)
\end{exer}
\begin{exer}{\exerused}{exer:rings:7.10}
  In the situation of the snake lemma, assume that $\lambda$ and $\nu$
  are isomorphisms. Use the snake lemma and prove that $\mu$ is an
  isomorphism. This is called the ``short five-lemma,'' as it follows
  immediately from the five-lemma (cf.\ \Cref{exer:rings:7.14}), as well as
  from the snake lemma. [\ref{exer:linrep:why-ext-is-called-ext-part1}, \ref{exer:homol:short-five-lemma}]
\end{exer}
\begin{exer}{\exerused}{exer:rings:7.11}
  Let
  \[0 \to M_1 \to N \to M_2 \to 0 \tag{$*$}\] be an exact sequence of
  $R$-modules. (This may be called an ``extension'' of $M_2$ by
  $M_1$.) Suppose there is any $R$-module homomorphism
  $N \to M_1 \oplus M_2$ making the diagram
  \[\begin{tikzcd}
      0 \ar[r] & M_1 \ar[r] \ar[d,equal] & N \ar[r] \ar[d] & M_2 \ar[r] \ar[d,equal] & 0 \\
      0 \ar[r] & M_1 \ar[r] & M_1 \oplus M_2 \ar[r] & M_2 \ar[r] & 0
    \end{tikzcd}\] commute, where the bottom sequence is the standard
  sequence of a direct sum. Prove that $(*)$ splits. [\S{}\ref{subsec:rings:split-exact-sequences}]
\end{exer}
\begin{exer}{\lnot}{exer:rings:7.12}
  Practice your diagram chasing skills by proving the ``four-lemma'':
  if
  \[\begin{tikzcd}
      A_1 \ar[r] \ar[d,"\alpha"'] & B_1 \ar[r] \ar[d,"\beta"'] & C_1 \ar[r] \ar[d,"\gamma"'] & D_1 \ar[d,"\delta"'] \\
      A_0 \ar[r] & B_0 \ar[r] & C_0 \ar[r] & D_0
    \end{tikzcd}\] is a commutative diagram of $R$-modules with exact
  rows, $\alpha$ is an epimorphism, and $\beta$, $\delta$ are
  monomorphisms, then $\gamma$ is a monomorphism. [\ref{exer:rings:7.13}, \ref{exer:homol:four-lemma-in-abelian-categories}]
\end{exer}
\begin{exer}{}{exer:rings:7.13}
  Prove another\footnote{It is in fact unnecessary to prove both
    versions, but to realize this one has to view the matter from the
    more general context of abelian categories; cf.\ \Cref{exer:homol:four-lemma-in-abelian-categories}.}
  version of the ``four-lemma'' of \cref{exer:rings:7.12}: if
  \[\begin{tikzcd}
      B_1 \ar[r] \ar[d,"\beta"'] & C_1 \ar[r] \ar[d,"\gamma"'] & D_1 \ar[r] \ar[d,"\delta"'] & E_1 \ar[d,"\epsilon"'] \\
      B_0 \ar[r] & C_0 \ar[r] & D_0 \ar[r] & E_0
    \end{tikzcd}\] is a commutative diagram of $R$-modules with exact
  rows, $\beta$ and $\delta$ are epimorphisms, and $\epsilon$ is a
  monomorphism, then $\gamma$ is an epimorphism.
\end{exer}
\begin{exer}{\lnot}{exer:rings:7.14}
  Prove the ``five-lemma'': if
  \[\begin{tikzcd}
      A_1 \ar[r] \ar[d,"\alpha"'] & B_1 \ar[r] \ar[d,"\beta"'] & C_1 \ar[r] \ar[d,"\gamma"'] & D_1 \ar[r] \ar[d,"\delta"'] & E_1 \ar[d,"\epsilon"'] \\
      A_0 \ar[r] & B_0 \ar[r] & C_0 \ar[r] & D_0 \ar[r] & E_0
    \end{tikzcd}\] is a commutative diagram of $R$-modules with exact
  rows, $\beta$ and $\delta$ are isomorphisms, $\alpha$ is an
  epimorphism, and $\epsilon$ is a monomorphism, then $\gamma$ is an
  isomorphism. (You can avoid the needed diagram chase by pasting
  together results from previous exercises.) [\ref{exer:rings:7.10}]
\end{exer}
\begin{exer}{\lnot}{exer:rings:7.15}
  Consider the following commutative diagram of $R$-modules:
  \[\begin{tikzcd}
      0 \ar[r] & L_2 \ar[r] \ar[d] & M_2 \ar[r] \ar[d,"\alpha"'] & N_2 \ar[r] \ar[d] & 0 \\
      0 \ar[r] & L_1 \ar[r] \ar[d] & M_1 \ar[r] \ar[d,"\beta"'] & N_1 \ar[r] \ar[d] & 0 \\
      0 \ar[r] & L_0 \ar[r] & M_0 \ar[r] & N_0 \ar[r] & 0
    \end{tikzcd}\] Assume that the three rows are exact and the two
  rightmost columns are exact. Prove that the left column is
  exact. Second version: assume that the three rows are exact and the
  two leftmost columns are exact; prove that the right column is
  exact. This is the ``nine-lemma.'' (You can avoid a diagram chase by
  applying the snake lemma; for this, you will have to turn the
  diagram by $90^\circ$.) [\ref{exer:rings:7.16}]
\end{exer}
\begin{exer}{}{exer:rings:7.16}
  In the same situation as in \cref{exer:rings:7.15}, assume that the
  three rows are exact and that the leftmost and rightmost columns are
  exact.
  \begin{itemize}
  \item Prove that $\alpha$ is a monomorphism and $\beta$ is an
    epimorphism.
  \item Is the central column necessarily exact? (Hint: No. Place
    $\Z \oplus \Z$ in the middle, and surround it artfully with six
    copies of $\Z$ and two $0$'s.)
  \item Assume further that the central column is a complex (that is,
    $\beta \circ \alpha = 0$); prove that it is then necessarily
    exact.
  \end{itemize}
\end{exer}
\begin{exer}{\lnot}{exer:rings:7.17}
  Generalize the previous two exercises as follows. Consider a
  (possibly infinite) commutative diagram of $R$-modules:
  \[\begin{tikzcd}
      & \vdots \ar[d] & \vdots \ar[d] & \vdots \ar[d] & \\
      0 \ar[r] & L_{i+1} \ar[r] \ar[d] & M_{i+1} \ar[r] \ar[d] & N_{i+1} \ar[r] \ar[d] & 0 \\
      0 \ar[r] & L_i \ar[r] \ar[d] & M_i \ar[r] \ar[d] & N_i \ar[r] \ar[d] & 0 \\
      0 \ar[r] & L_{i-1} \ar[r] \ar[d] & M_{i-1} \ar[r] \ar[d] & N_{i-1} \ar[r] \ar[d] & 0 \\
      & \vdots & \vdots & \vdots &
    \end{tikzcd}\] in which the central column is a complex and every
  row is exact. Prove that the left and right columns are also
  complexes. Prove that if any two of the columns are exact, so is the
  third. (The first part is straightforward. The second part will take
  you a couple of minutes now due to the needed diagram chases, and a
  couple of seconds later, once you learn about the long exact
  (co)homology sequence in \S{}\ref{subsec:homol:the-long-exact-cohomology-sequence}.) [\ref{exer:homol:two-exact-columns-implies-third-abelian-category}]
\end{exer}

%%% Local Variables:
%%% mode: LaTeX
%%% TeX-engine: luatex
%%% TeX-master: "../master"
%%% End:
